% Definicja składni języka Kotlin dla pakietu listings (brak wsparcia wbudowanego)
\lstdefinelanguage{Kotlin}{
 keywords={package, import, class, object, interface, data, fun, val, var,
 private, public, internal, const, companion, return, if, else, when, for,
 while, do, break, continue, in, is, as, null, true, false, override,
 suspend, this, super, try, catch, finally, throw, by, out, sealed, enum},
 keywordstyle=\bfseries\color{EEUltramarine},
 ndkeywords={Int, Float, Long, Boolean, String, List, Map, Bitmap, Flow,
 Mutex, ArrayDeque, Dispatchers, Unit},
 ndkeywordstyle=\color{EEBlueDark},
 comment=[l]{//},
 morecomment=[s]{/*}{*/},
 commentstyle=\color{EEBlueDark},
 stringstyle=\color{EETangerine},
 morestring=[b]",
 sensitive=true
}
% Poprawna obsługa polskich znaków diakrytycznych wewnątrz listingów
\lstset{
 extendedchars=true,
 literate=%
 {ą}{{\k{a}}}1 {ć}{{\'c}}1 {ę}{{\k{e}}}1 {ł}{{\l{}}}1 {ń}{{\'n}}1
 {ó}{{\'o}}1 {ś}{{\'s}}1 {ź}{{\'z}}1 {ż}{{\.z}}1
 {Ą}{{\k{A}}}1 {Ć}{{\'C}}1 {Ę}{{\k{E}}}1 {Ł}{{\L{}}}1 {Ń}{{\'N}}1
 {Ó}{{\'O}}1 {Ś}{{\'S}}1 {Ź}{{\'Z}}1 {Ż}{{\.Z}}1
}

\chapter{Wstęp}\label{chap:wstep}
Przedmiotem pracy jest automatyczna identyfikacja gatunków grzybów na urządzeniach mobilnych pracujących autonomicznie, bez połączenia z serwerem. Grzybobranie – w Polsce i szerzej w Europie Środkowej – pozostaje powszechną praktyką rekreacyjną i kulinarną. Jednocześnie niesie realne ryzyko zdrowotne, gdy owocniki zostaną źle rozpoznane. Większość ciężkich zatruć nie wynika ze świadomego spożycia gatunków trujących, lecz z powierzchownej oceny cech morfologicznych przez osoby bez przygotowania mykologicznego. Ta rozbieżność między łatwą dostępnością grzybów w lasach a ograniczoną wiedzą przeciętnego zbieracza uzasadnia podjęcie tego tematu.

Za przeważającą część śmiertelnych zatruć grzybami na świecie, szacowaną na ponad 90\%, odpowiadają gatunki zawierające amatoksyny, przede wszystkim muchomor sromotnikowy (\textit{Amanita phalloides}) \cite{garcia2015amanita, mengs2012legalon}. Toksyny te cechują się wysoką hepatotoksycznością, a dawka śmiertelna $\alpha$-amanityny dla człowieka wynosi około 0,1~mg/kg masy ciała, co oznacza, że pojedynczy owocnik może stanowić zagrożenie dla życia osoby dorosłej \cite{garcia2015amanita}. Dodatkowym czynnikiem pogarszającym rokowanie jest wielogodzinny okres utajenia, w którym objawy kliniczne nie występują, przez co pomoc medyczna bywa udzielana ze znacznym opóźnieniem. Według danych europejskich i północnoamerykańskich ośrodków toksykologicznych liczbę zgonów spowodowanych zatruciami grzybami szacuje się na kilkadziesiąt do około dwustu przypadków rocznie \cite{mengs2012legalon}, przy czym rzeczywista skala zjawiska pozostaje trudna do precyzyjnego ustalenia z powodu niepełnej sprawozdawczości. W warunkach polskich zatrucia grzybami wykazują wyraźną sezonowość i nasilają się w okresie jesiennym, a ich bezpośrednią przyczyną jest najczęściej pomylenie gatunku jadalnego z jego trującym odpowiednikiem.

Z tych powodów sens ma narzędzie wspierające wstępną identyfikację grzybów bezpośrednio w terenie, na typowym smartfonie. System ma wspierać użytkownika przy ocenie owocników ułożonych na płaskiej powierzchni. Zwłaszcza w kwestii bezpieczeństwa mykologicznego i wyraźnego sygnalizowania gatunków potencjalnie trujących.

\section{Motywacja podjęcia tematu}
Motywacja do opracowania systemu wynika z trzech przesłanek: skali problemu zatruć, trudności rozróżniania gatunków bliźniaczych oraz ograniczeń dotychczasowych rozwiązań mobilnych. Kolejne podrozdziały omawiają je po kolei i uzasadniają wybór zastosowanych technologii.

\subsection{Problem gatunków bliźniaczych}
Kluczową trudnością w automatyzacji rozpoznawania grzybów jest zjawisko gatunków bliźniaczych (ang. \textit{look-alikes}), czyli par lub grup gatunków o zbliżonych cechach makroskopowych, spośród których część jest jadalna, a część silnie toksyczna. Wizualne podobieństwo owocników sprawia, że rozróżnienie oparte wyłącznie na kształcie kapelusza, barwie czy pokroju bywa zawodne nawet dla doświadczonych zbieraczy, a subtelne cechy diagnostyczne, takie jak budowa blaszek, obecność pierścienia czy pochwy u podstawy trzonu, są często pomijane.

Do najbardziej reprezentatywnych przykładów należy pieprznik jadalny (\textit{Cantharellus cibarius}), popularnie zwany kurką, mylony z lisówką pomarańczową (\textit{Hygrophoropsis aurantiaca}) oraz z toksycznym gatunkiem \textit{Omphalotus olearius}, którego spożycie wywołuje ostre dolegliwości żołądkowo-jelitowe, a nie zatrucie śmiertelne. Młode owocniki pieczarki polnej (\textit{Agaricus campestris}) do złudzenia przypominają wczesne stadia muchomora sromotnikowego (\textit{Amanita phalloides}) oraz jadowitego (\textit{Amanita virosa}); ta pomyłka bywa śmiertelna. We wczesnym stadium rozwoju oba muchomory przypominają jadalne gatunki o białych kapeluszach. Możliwie wiarygodne rozróżnianie takich par jest zatem jednym z kluczowych kryteriów projektowych systemu, w którym błąd klasyfikacji może mieć konsekwencje zagrażające życiu.

Skala problemu wzrasta dodatkowo wraz z liczbą rozpoznawanych taksonów. Zadanie klasyfikacji drobnoziarnistej (ang. \textit{fine-grained classification}) na zbiorze obejmującym $147$ klas grzybów, odpowiadających $145$ odrębnym gatunkom biologicznym (dwie pary klas to synonimy taksonomiczne odziedziczone po etykietach iNaturalist; tab.~\ref{tab:zal2_gatunki}), wymaga od modelu różnicowania klas na podstawie subtelnych, lokalnych cech morfologicznych – układu blaszek, faktury powierzchni kapelusza czy odcienia miazgi po przełamaniu – przy jednoczesnym zachowaniu odporności na wewnątrzklasową zmienność wynikającą z wieku owocnika, wilgotności podłoża oraz oświetlenia scenerii leśnej. W przeciwieństwie do zadań klasyfikacji gruboziarnistej, w których klasy różnią się istotnie w przestrzeni cech wysokopoziomowych, granice decyzyjne pomiędzy gatunkami bliźniaczymi bywają wąskie i nieliniowe, a błąd predykcji koncentruje się właśnie w ich sąsiedztwie. Zwiększenie liczby klas z rzędu kilkudziesięciu do $147$ klas pogłębia to zjawisko, ponieważ przestrzeń cech musi jednocześnie separować znacznie więcej par potencjalnie mylonych taksonów, co wymaga architektury o odpowiednio dużej pojemności reprezentacyjnej. Zagadnienie to rozwinięto w rozdziale \ref{chap:przeglad_rozwiazan}.

\subsection{Uzasadnienie wyboru detekcji wieloobiektowej}
Naturalnym podejściem do rozpoznawania gatunku grzyba na podstawie obrazu wydaje się klasyfikacja całego obrazu (ang. \textit{image classification}), przypisująca pojedynczą etykietę klasy do całej sceny. Podejście to zakłada jednak obecność dokładnie jednego dominującego obiektu w kadrze i traci użyteczność w typowym scenariuszu domowym, w którym zebrane owocniki są rozkładane obok siebie na stole lub w koszu. W takich warunkach kilka egzemplarzy, potencjalnie należących do różnych gatunków, występuje jednocześnie w jednym kadrze, a ich fragmenty wzajemnie się przysłaniają i nakładają na siebie, co uniemożliwia poprawne przypisanie jednej etykiety do całego obrazu.

Rozwiązaniem tego problemu jest detekcja wieloobiektowa (ang. \textit{multi-object detection}), która jednocześnie lokalizuje każdy owocnik za pomocą ramki otaczającej oraz przypisuje mu odrębną klasę gatunkową. Rozwiązanie takie pozwala na indywidualną ocenę każdego egzemplarza w scenie i – co ważne w kontekście bezpieczeństwa – na selektywne oznaczenie obiektów sklasyfikowanych jako trujące. Ponieważ narzędzie ma działać na strumieniu obrazu z kamery w czasie rzeczywistym, wybrano jednoprzebiegowy detektor z rodziny YOLO (ang. \textit{You Only Look Once}), której pierwotna koncepcja realizuje detekcję jako pojedyncze przejście sieci neuronowej i zapewnia wysoką przepustowość przetwarzania \cite{redmon2016yolo}. Jako model bazowy przyjęto architekturę YOLO11 Medium (\texttt{YOLO11m}), stanowiącą kompromis pomiędzy dokładnością a kosztem obliczeniowym, adekwatny do ograniczeń urządzeń mobilnych przy zadaniu rozróżniania wielu gatunków \cite{ultralytics2024yolo11}.

\subsection{Konieczność działania w trybie offline}
Całkowita niezależność od sieci to wymaganie twarde, nie preferencja. Grzybobranie odbywa się z reguły w głębi lasu, gdzie zasięg GSM i LTE bywa ograniczony albo zerowy. W takich warunkach rozwiązania oparte na architekturze klient–serwer, przesyłające obraz do zdalnego interfejsu chmurowego (ang. \textit{cloud-based API}), stają się bezużyteczne dokładnie w tym miejscu i momencie, gdy są najbardziej potrzebne.

Dlatego przyjęto, że cały proces inferencji musi być realizowany lokalnie, na urządzeniu użytkownika (ang. \textit{on-device inference}). Wymóg ten narzuca dodatkowe ograniczenia dotyczące rozmiaru modelu oraz zapotrzebowania na pamięć i moc obliczeniową. Podporządkowano im dobór architektury sieci oraz proces jej optymalizacji. Aby umożliwić efektywne uruchomienie modelu na sprzęcie mobilnym, przeprowadzono jego konwersję do formatu TensorFlow Lite z zastosowaniem kwantyzacji wag \cite{tflite}; szczegóły podano w rozdz.~\ref{chap:implementacja}.

\section{Cel i zakres pracy}
Celem pracy jest zaprojektowanie, implementacja i ewaluacja autonomicznego systemu detekcji oraz klasyfikacji gatunków grzybów na platformie Android, działającego offline i w czasie rzeczywistym. Czas rzeczywisty rozumiany jest w pracy jako interaktywny odczyt z podglądu kamery: od wykrycia owocnika w kadrze do wyświetlenia jego ramki lub ostrzeżenia mija wyraźnie mniej niż $1$\,s (podrozdz.~\ref{sec:wymagania_niefunkcjonalne}). Kadencja analizy $4$\,FPS ($250$\,ms) jest docelową częstotliwością przetwarzania klatek, a nie definicją czasu rzeczywistego. Realizacja obejmuje model uczenia maszynowego oraz aplikację mobilną, która ma wspomagać identyfikację owocników w terenie, a nie zastępować ocenę eksperta.

Osiągnięcie celu głównego wymagało zrealizowania następujących zadań szczegółowych:
\begin{enumerate}[label=\textbf{Z\arabic*.}]
 \item \textbf{Opracowanie i wytrenowanie modelu detekcji} wieloobiektowej YOLO11m obejmującego $147$ klas grzybów ($145$ odrębnych gatunków biologicznych), zwłaszcza par gatunków bliźniaczych, przy wykorzystaniu uczenia transferowego na bazie wag wstępnie wytrenowanych na zbiorze COCO \cite{lin2014coco}.
 \item \textbf{Zaprojektowanie półautomatycznego potoku anotacji} zbioru danych z wykorzystaniem modelu wizyjno-językowego GroundingDINO \cite{liu2023grounding}, ograniczającego pracochłonność ręcznego etykietowania ramek otaczających.
 \item \textbf{Przeprowadzenie optymalizacji i konwersji modelu} do postaci umożliwiającej lokalną inferencję na urządzeniu mobilnym, obejmującej dwuetapową ścieżkę PyTorch $\to$ TensorFlow SavedModel $\to$ TensorFlow Lite oraz dobór precyzji numerycznej FP16.
 \item \textbf{Zintegrowanie modelu z natywną aplikacją} systemu Android, napisaną w języku Kotlin z wykorzystaniem biblioteki CameraX oraz deklaratywnego interfejsu Jetpack Compose, wraz z implementacją asymetrycznej, programowej logiki bezpieczeństwa.
 \item \textbf{Przeprowadzenie ilościowej oceny skuteczności} rozwiązania na izolowanym zbiorze testowym, z naciskiem na wykrywalność gatunków trujących oraz macierz ryzyka bezpieczeństwa mykologicznego.
\end{enumerate}

Zakres pracy obejmuje stworzenie kompletnego potoku przetwarzania danych (ang. \textit{machine learning pipeline}), automatyzację procesu anotacji zbioru danych z wykorzystaniem modelu GroundingDINO \cite{liu2023grounding}, trening głębokiej sieci neuronowej YOLO11 Medium (\texttt{YOLO11m}), jej konwersję do formatu TensorFlow Lite (ścieżką PyTorch $\to$ TensorFlow SavedModel $\to$ TFLite), a także integrację zoptymalizowanego modelu z aplikacją mobilną wykorzystującą interfejs CameraX. Poza zakresem pracy pozostają zagadnienia mykologiczne wykraczające poza identyfikację wizualną, takie jak analiza chemiczna czy mikroskopowa cech diagnostycznych owocników.

\section{Hipotezy badawcze}
Realizację celu głównego oparto na dwóch hipotezach badawczych, weryfikowanych ilościowo w rozdziale \ref{chap:badania}:
\begin{description}[topsep=2pt,itemsep=2pt]
 \item[$H_1$] Jednoprzebiegowa detekcja wieloobiektowa YOLO11m w precyzji FP16 umożliwia autonomiczną identyfikację $147$ klas grzybów ($145$ odrębnych gatunków biologicznych) na urządzeniu mobilnym, bez dostępu do sieci, w czasie rzeczywistym rozumianym jako interaktywny odczyt z podglądu kamery w czasie krótszym niż $1$\,s.
 \item[$H_2$] Programowa, asymetryczna logika progów decyzyjnych ($0{,}18$ dla klas trujących, $0{,}60$ dla klas jadalnych) ogranicza wskaźnik krytycznego ryzyka (CFPR) na izolowanym zbiorze testowym względem pojedynczego progu globalnego.
\end{description}
Hipoteza $H_1$ dotyczy warstwy algorytmiczno-sprzętowej: weryfikuje się ją pomiarem $mAP$ wag PyTorch na izolowanym teście, porównaniem $mAP$ na próbie $2\,500$ kadrów między \texttt{.pt} a TFLite FP16 (Interpreter i dekoder zgodny z \texttt{MushroomDetector.kt}, nie LiteRT \texttt{val}; podrozdz.~\ref{sec:tflite_parity}) oraz pomiarem na S23 Ultra czasu \texttt{detect()}, opóźnienia end-to-end klatki i czasu do pierwszego wyświetlenia ramki (podrozdz.~\ref{sec:latency_e2e}). Hipoteza $H_2$ dotyczy pary progów $0{,}18$/$0{,}60$ względem jednego progu globalnego $0{,}30$ na \emph{tym samym} zbiorze testowym. Oba CFPR policzono skryptem \texttt{test\_h2\_baseline.py} na wagach PyTorch: $2{,}06\%$ wobec $3{,}66\%$. Dla artefaktu TFLite FP16 z dekoderem aplikacji (\texttt{test\_tflite\_cfpr.py}) wynoszą one $2{,}11\%$ wobec $3{,}64\%$ (podrozdz.~\ref{sec:safety_matrix}). Bufor konsensusu klatek ($4$ z $5$) nie wchodzi w treść $H_2$; opisano go w rozdz.~\ref{chap:implementacja} jako mechanizm stabilizacji czasowej nakładki.

\section{Układ pracy}
Praca składa się z sześciu rozdziałów, zgodnych z kolejnością prac nad projektem. Rozdział bieżący przedstawia motywację, cel i zakres. W rozdziale drugim omówiono stan wiedzy: aplikacje mobilne do identyfikacji grzybów, metody detekcji i klasyfikacji obiektów oraz optymalizację modeli pod urządzenia o ograniczonych zasobach. Trzeci rozdział dotyczy projektu systemu. Wymagań funkcjonalnych i niefunkcjonalnych oraz doboru architektury.

Czwarty opisuje implementację: przygotowanie danych, anotację, uczenie, konwersję modelu i aplikację mobilną. Piąty zawiera badania skuteczności i wydajności, w tym analizę gatunków bliźniaczych oraz jakościową ocenę detekcji w warunkach terenowych na danych zewnętrznych. Rozdział szósty podsumowuje wyniki i wskazuje kierunki dalszych prac.


\chapter{Przegląd istniejących rozwiązań}\label{chap:przeglad_rozwiazan}
Autonomiczny system identyfikacji owocników na urządzeniu wbudowanym wymaga uprzedniego przeglądu trzech obszarów: dostępnych aplikacji, metod detekcji i klasyfikacji w obrazie oraz technik dostosowania modeli do ograniczonego sprzętu. Takie tło pozwala ocenić istniejące podejścia, wskazać ich ograniczenia i uzasadnić decyzje projektowe omówione dalej.

\section{Analiza istniejących aplikacji mobilnych}
Dwie aplikacje wyznaczają górną granicę możliwości komercyjnych rozwiązań. iNaturalist \cite{inaturalist} gromadzi miliony obserwacji przyrodniczych i klasyfikuje je modelami trenowanymi na tej samej skali; Picture Mushroom \cite{picturemushroom} zawęża zakres do grzybów i osiąga imponującą trafność dla zdjęć studyjnych. Obie buduje się jednak na tym samym założeniu: smartfon to tylko aparat. Obraz wędruje na serwer, tam model inferuje i wraca wynik.

Ten wzorzec klient--serwer sprawdza się w parku z LTE. W lesie LTE zwykle nie ma. Polskie kompleksy leśne (Puszcza Białowieska, Bory Tucholskie, lasy Sudetów) to typowe tereny grzybobrania i typowe martwe strefy GSM. W głębi drzewostanu sygnał jest słaby albo go nie ma wcale. Gdy grzybiarz wyciąga telefon nad znaleziskiem, chmura mu nie odpowie.

Trzy konsekwencje tego wzorca są bezpośrednio szkodliwe dla scenariusza terenowego:
\begin{enumerate}
 \item Aplikacja traci funkcjonalność dokładnie tam, gdzie jest najbardziej potrzebna: w lesie bez zasięgu GSM/LTE.
 \item Dwukierunkowa transmisja zdjęcia wysokiej rozdzielczości wyklucza analizę strumienia wideo w czasie rzeczywistym.
 \item Żadna z aplikacji nie obsługuje jednoczesnej detekcji wielu owocników w kadrze -- jedno zdjęcie, jeden wynik.
\end{enumerate}

Porównanie właściwości aplikacji komercyjnych i systemu opracowanego w tej pracy zestawiono w tab.~\ref{tab:aplikacje_porownanie}. Porównanie prowadzono według pięciu kryteriów istotnych dla scenariusza terenowego: modelu wykonania inferencji, dostępności bez łączności sieciowej, zdolności do analizy strumienia wideo, obsługi wielu obiektów w jednym kadrze oraz obecności dedykowanej warstwy sygnalizacji ryzyka toksykologicznego.

\begin{table}[htbp]
 \centering
 \caption{Porównanie funkcjonalne popularnych aplikacji identyfikacyjnych z systemem opracowanym w pracy. Źródło: opracowanie własne na podstawie \cite{inaturalist, picturemushroom}.}
 \label{tab:aplikacje_porownanie}
 \begin{tabular}{|p{5.4cm}|c|c|c|}
 \hline
 \textbf{Kryterium} & \textbf{iNaturalist} & \textbf{Picture Mushroom} & \textbf{System własny} \\
 \hline \hline
 Model wykonania inferencji & klient--serwer & klient--serwer & na urządzeniu \\
 \hline
 Praca bez dostępu do sieci & nie & nie & tak \\
 \hline
 Analiza strumienia wideo & nie & nie & tak ($4$\,FPS) \\
 \hline
 Wiele obiektów w kadrze & nie & nie & tak \\
 \hline
 Warstwa ostrzegania toksykologicznego & brak & informacyjna & asymetryczna \\
 \hline
 \end{tabular}
\end{table}

Chmura zawodzi dokładnie w środowisku, dla którego narzędzie jest projektowane. Dlatego w tej pracy przyjęto odwrotny paradygmat: przetwarzanie krawędziowe (ang. \textit{edge AI}), gdzie cały model działa lokalnie na urządzeniu. To oznacza konieczność zmieszczenia wystarczająco dobrego detektora w budżecie pamięci i mocy obliczeniowej smartfona. Techniki tego dostosowania opisano w kolejnych podrozdziałach.

\section{Przegląd metod detekcji i klasyfikacji obiektów}
Głębokie sieci neuronowe zastąpiły ręcznie projektowane deskryptory obrazu – ich zdolność do hierarchicznej ekstrakcji cech pozwala odróżnić blaszki muchomora od blaszek pieczarki bez jawnego kodowania zasad mykologii. Poniżej opisano różnicę między klasyfikacją a detekcją, porównano rodziny detektorów i uzasadniono dobór architektury zastosowanej w pracy.

\subsection{Splotowe sieci neuronowe i transformatory}
Splotowe sieci neuronowe (CNN) działają na operacji splotu obrazu z wyuczalnymi filtrami \cite{lecun1998gradient}. Pierwsze warstwy wychwytują krawędzie i tekstury, głębsze składają je w coraz bardziej abstrakcyjne cechy, aż do poziomu rozróżnialnego gatunkowo układu blaszek czy faktury kapelusza. Współdzielenie wag nadaje niezmienniczość translacyjną, a pooling redukuje wymiarowość. Klasyczną ilustrację tej hierarchii pokazuje rys.~\ref{fig:lenet5_web} z architekturą LeNet-5 \cite{lecun1998gradient}. Połączenia rezydualne (ResNet) usunęły barierę zanikającego gradientu i umożliwiły efektywne uczenie sieci o dziesiątkach warstw \cite{he2016resnet}.

\begin{figure}[htbp]
 \centering
 \includegraphics[width=0.78\textwidth]{lenet5_official}
 \caption{Architektura klasycznej splotowej sieci neuronowej LeNet-5. Źródło: \cite{lecun1998gradient}.}
 \label{fig:lenet5_web}
\end{figure}

Transformatory wizyjne (ViT) \cite{dosovitskiy2021vit} zaadaptowały mechanizm samouwagi (ang. \textit{self-attention}) z NLP \cite{vaswani2017attention} do analizy obrazu. Samouwaga modeluje globalne zależności między odległymi fragmentami obrazu, czego sieć splotowa z ograniczonym polem receptywnym nie zrobi tanio. W detekcji obiektów mechanizmy uwagi działają w rodzinie DETR, na której opiera się model GroundingDINO \cite{liu2023grounding} – kluczowy dla auto-anotacji w tej pracy.

\subsection{Detekcja otwartego słownika: GroundingDINO}
YOLO zna tylko klasy, które dostał w treningu. Do budowy zbioru to fatalna cecha: trzeba najpierw narysować setki tysięcy ramek, zanim w ogóle ruszy uczenie gatunków. Ręczne BBOX-y dla pełnej bazy byłyby nierealne w skali pracy inżynierskiej.

GroundingDINO omija ten impas \cite{liu2023grounding}. To detektor otwartego słownika: zamiast indeksu klasy dostaje tekst. Enkoder Swin patrzy na obraz, BERT czyta frazę, a \textit{cross-modality attention} liczy zgodność regionu z tekstem (rys.~\ref{fig:grounding_dino_web}). W potoku anotacji zapytanie brzmiało po prostu \texttt{"mushroom"}. Model wskazywał owocnik bez wiedzy, czy to borowik, czy muchomor – gatunek brał się z folderu iNaturalist, nie z głowicy. Zero-shotowa lokalizacja oddzieliła „gdzie jest grzyb” od „jaki to gatunek”. Bez tego zbioru $147$ klas nie byłoby.

\begin{figure}[htbp]
 \centering
 \includegraphics[width=0.88\textwidth]{grounding_dino_arch_official}
 \caption{Architektura detektora otwartego słownika Grounding DINO. Źródło: \cite{liu2023grounding}.}
 \label{fig:grounding_dino_web}
\end{figure}

\subsection{Klasyfikacja obrazu a detekcja wieloobiektowa}
Klasyfikator widzi cały kadr i wydaje jeden werdykt. Formalnie odwzorowanie takie przypisuje obrazowi jedną etykietę spośród $K$ klas:
\begin{equation}
f_{\mathrm{cls}}:\; \mathbb{R}^{H\times W\times C} \rightarrow \{1, 2, \dots, K\},
\end{equation}
gdzie
\begin{description}[topsep=2pt,itemsep=-3pt]
 \item[$H, W, C$] wysokość, szerokość oraz liczba kanałów obrazu wejściowego,
 \item[$K$] liczba rozpoznawanych klas.
\end{description}
Softmax zamienia wynikowy wektor logitów w rozkład prawdopodobieństwa, a odpowiedzią jest klasa o najwyższym wyniku. Jeden kadr, jedna odpowiedź. Nie ma mowy o położeniu. Nie ma mowy o kilku gatunkach jednocześnie – a dokładnie taki jest las: kosz, w którym leżą kurka, podgrzybek i niejadalna gołąbka naraz.

Detekcja wieloobiektowa rozrywa to ograniczenie. Wynikiem nie jest etykieta, lecz zbiór krotek:
\begin{equation}
f_{\mathrm{det}}:\; \mathbb{R}^{H\times W\times C} \rightarrow \big\{(\mathbf{b}_i, c_i, s_i)\big\}_{i=1}^{N},
\end{equation}
gdzie
\begin{description}[topsep=2pt,itemsep=-3pt]
 \item[$\mathbf{b}_i$] ramka otaczająca (BBOX) $i$-tego obiektu, opisana wektorem współrzędnych $(x_i, y_i, w_i, h_i)$,
 \item[$c_i$] etykieta klasy przypisana obiektowi,
 \item[$s_i$] wartość pewności detekcji (ang. \textit{confidence score}),
 \item[$N$] liczba wykrytych obiektów.
\end{description}
Miarą jakości dopasowania ramki predykowanej $B_p$ do ramki referencyjnej $B_{gt}$ jest współczynnik przecięcia nad sumą (IoU), zdefiniowany jako:
\begin{equation}
\mathrm{IoU}(B_p, B_{gt}) = \frac{\lvert B_p \cap B_{gt} \rvert}{\lvert B_p \cup B_{gt} \rvert}.
\end{equation}
Jeden owocnik, jedna ramka i etykieta, co ilustruje rys.~\ref{fig:class_vs_det_web}. Klasyfikator wydaje jeden werdykt na cały kadr; detektor opisuje każdy obiekt z osobna. W kępie opieńki (\textit{Armillaria mellea}), gdzie kilkanaście trzonów styka się w ściółce, tylko detekcja pozwala odróżnić jadalny owocnik od sąsiedniego, który nim nie jest.

\begin{figure}[htbp]
 \centering
 \includegraphics[width=0.9\textwidth]{classification_vs_detection_official}
 \caption{Porównanie klasyfikacji całego obrazu a detekcji wieloobiektowej. Źródło: opracowanie własne na podstawie \cite{redmon2016yolo, lin2014coco}.}
 \label{fig:class_vs_det_web}
\end{figure}

Trening od zera na $147$ klasach przekroczyłby budżet czasu i pamięci VRAM. Dlatego wagi startują z COCO, a następnie model jest dostrajany w~całości, bez zamrażania warstw (\texttt{freeze=0}; rozdz.~\ref{chap:implementacja}) \cite{alhadhrami2018transfer, lin2014coco}. Inicjalizacja z~COCO dostarcza gotowych niskopoziomowych filtrów krawędzi i~tekstur, które w~praktyce zmieniają się najmniej, podczas gdy wyższe warstwy i~głowica adaptują się do cech mykologicznych. Rys.~\ref{fig:transfer_learning_web} ilustruje ogólną ideę uczenia transferowego (w wariancie z~zamrożonym szkieletem); w~niniejszej pracy szkieletu nie zamrażano. Zastosowano wagi \texttt{yolo11m.pt} (rozdz.~\ref{chap:implementacja}).

\begin{figure}[htbp]
 \centering
 \includegraphics[width=0.85\textwidth]{transfer_learning_official}
 \caption{Ogólna idea uczenia transferowego w wariancie z zamrożonym szkieletem. W niniejszej pracy nie zamrażano warstw (\texttt{freeze=0}); schemat ilustruje jedynie koncepcję. Źródło: \cite{alhadhrami2018transfer}.}
 \label{fig:transfer_learning_web}
\end{figure}

\subsection{Detektory dwuetapowe i jednoetapowe}
Dwa przejścia sieci lub jedno – tyle się sprowadza podział rodzin detekcyjnych. Faster R-CNN \cite{ren2015fasterrcnn} generuje najpierw propozycje regionów, potem dopiero klasyfikuje i prostuje ramki. Precyzja wysoka; koszt wysoki. Na serwerze z GPU to akceptowalny kompromis. Na Snapdragonie 8 Gen 2 w lesie rozwiązanie dwuetapowe odpada: setki megabajtów, dwa przebiegi na klatkę i thermal throttling po kilku minutach.

Detektory jednoetapowe – SSD \cite{liu2016ssd} i cała rodzina YOLO \cite{redmon2016yolo} – regresują ramki i klasy bezpośrednio na siatce obrazu w jednym przebiegu. Jeden tensor wyjściowy. Przewidywalne opóźnienie. Model mieszczący się w kilkudziesięciu megabajtach flash.

W tab.~\ref{tab:detektory} zestawiono obie rodziny. Na urządzeniu z procesorem Snapdragon jednoetapowy detektor lepiej spełnia budżet opóźnienia: jeden przebieg sieci, przewidywalny czas inferencji i mniejszy artefakt binarny.

\begin{table}[h!]
 \centering
 \caption{Porównanie detektorów dwuetapowych i jednoetapowych. Źródło: opracowanie własne na podstawie \cite{ren2015fasterrcnn, redmon2016yolo, liu2016ssd}.}
 \label{tab:detektory}
 \begin{tabular}{|l|c|c|}
 \hline
 \textbf{Kryterium} & \textbf{Dwuetapowe} & \textbf{Jednoetapowe} \\
 \hline \hline
 Reprezentatywny model & Faster R-CNN & YOLO, SSD \\
 \hline
 Liczba przebiegów sieci & dwa & jeden \\
 \hline
 Dokładność detekcji & bardzo wysoka & wysoka \\
 \hline
 Szybkość inferencji & niska & bardzo wysoka \\
 \hline
 Przydatność na urządzeniach mobilnych & ograniczona & wysoka \\
 \hline
 \end{tabular}
\end{table}

\subsection{Ewolucja architektury YOLO}
Redmon i wsp. w 2016 roku postawili jedną tezę: detekcja to regresja na siatce: jeden przebieg wystarczy \cite{redmon2016yolo}. Teza słuszna. Przez osiem lat kolejne generacje zmniejszały lukę wobec detektorów dwuetapowych bez rezygnacji z tej zasady: YOLOv3 dołożył predykcję wieloskalową \cite{redmon2018yolov3}, YOLOv8 usunął kotwice i ujednolicił API Ultralytics \cite{ultralytics2023yolov8}, YOLO11 przeprojektował szkielet i szyję \cite{ultralytics2024yolo11}.

Historia wersji jest tu mniej istotna niż jeden konkretny wniosek inżynierski: YOLO11m z $\approx20{,}1$~mln parametrów mieści się w budżecie TFLite FP16 ($38{,}58$~MiB) i jest wystarczająco głęboki, by odróżniać $147$ klas, w tym gatunki bliźniacze – łącznie z parami takimi jak \textit{Amanita phalloides}/\textit{Amanita citrina} czy \textit{Galerina marginata}/\textit{Kuehneromyces mutabilis}. Nano jest za płytki: ablacja to dokumentuje (rozdz.~\ref{chap:badania}). Wariantu Large nie trenowano i nie mierzono na telefonie; odrzucono go na podstawie większej liczby parametrów względem budżetu artefaktu, bez osobnego eksperymentu.

Współczesne architektury YOLO dzielą się funkcjonalnie na trzy zasadnicze części:
\begin{enumerate}
 \item Szkielet (ang. \textit{backbone}) jest siecią splotową odpowiedzialną za ekstrakcję hierarchicznych map cech z obrazu wejściowego.
 \item Szyja (ang. \textit{neck}) agreguje cechy z różnych poziomów rozdzielczości, najczęściej z wykorzystaniem struktury piramidy cech typu PANet.
 \item Głowica (ang. \textit{head}) generuje końcowe predykcje ramek oraz klas, w nowszych wersjach w wariancie bezkotwicowym (ang. \textit{anchor-free}).
\end{enumerate}

Oficjalny schemat Ultralytics (rys.~\ref{fig:yolo11_web}) pokazuje ten podział wprost: wejście $640\times640\times3$ przechodzi przez szkielet z blokami \texttt{Conv}/\texttt{C3k2}, następnie przez szyję z \texttt{SPPF}, \texttt{C2PSA}, \texttt{Upsample} i \texttt{Concat}, a trzy gałęzie \texttt{Detect} na skalach $80\times80$, $40\times40$ i $20\times20$ zasilają zadania Object Detection \cite{ultralytics2024yolo11}.

\begin{figure}[htbp]
 \centering
 \includegraphics[width=0.88\textwidth]{yolo11_arch_official}
 \caption{Struktura detektora YOLO11 z podziałem na Backbone, Neck i Decoupled Head. Źródło: \cite{ultralytics2024yolo11}.}
 \label{fig:yolo11_web}
\end{figure}

Trzyskalowa głowica nie jest ozdobą diagramu. To odpowiedź na geometrię owocników. Mapa $80\times80$ łapie drobne kapelusze i młode owocniki na ściółce; $40\times40$ obejmuje typowy rozmiar w kadrze telefonu; $20\times20$ obejmuje duże skupiska, np.\ kępy opieńki miodowej (\textit{Armillaria mellea}), gdzie wiele trzonów styka się i nakłada. Bez fuzji skal model albo sklejałby sąsiednie owocniki w jedną ramkę, albo gubił drobne okazy między liśćmi.

Szkielet oparto na blokach \texttt{C3k2} (ang. \textit{Cross Stage Partial with 2 kernels}) – wariancie CSP z dwoma ścieżkami cech i lekkimi splotami \cite{ultralytics2024yolo11}. Jedna ścieżka liczy, druga omija część transformacji; wynik jest konkatenowany. Mniej redundantnego gradientu, mniej parametrów, a tekstury blaszek i łusek nadal przechodzą wyżej. Moduł \texttt{SPPF} na dnie szkieletu zbiera kontekst przestrzenny sekwencyjnym max-poolingiem zamiast kilku równoległych jąder SPP. Taniej obliczeniowo przy podobnym zasięgu recepcji. Blok \texttt{C2PSA} w szyi dokłada uwagę przestrzenną przed rozesłaniem cech do trzech detektorów.

Głowica YOLO11 jest bezkotwicowa i rozłączna: regresja ramki i klasyfikacja gatunku idą osobnymi podsieciami na tej samej mapie cech. Wspólna głowica w starszych YOLO pchała gradient lokalizacji przeciwko gradientowi klasy – tu konflikt jest mniejszy \cite{ultralytics2024yolo11}. Brak kotwic oznacza też brak ręcznego doboru \textit{anchor boxes} pod wąskie trzony i płaskie kapelusze; sieć regresuje ramkę względem środka komórki siatki.

Trzy człony budują łączną stratę: geometria ramki, precyzja brzegów i klasyfikacja gatunkowa. Geometrię obsługuje Complete IoU (CIoU) – klasyczne IoU to za mało, bo nie karze za przesunięty środek ani złe proporcje \cite{ultralytics2024yolo11}:
\begin{equation}
\mathcal{L}_{\mathrm{CIoU}} = 1 - \mathrm{IoU} + \frac{\rho^2(b, b^{gt})}{c^2} + \alpha v,
\end{equation}
gdzie
\begin{description}[topsep=2pt,itemsep=-3pt]
 \item[$\rho^2(b, b^{gt})$] kwadrat odległości euklidesowej pomiędzy środkami ramki predykowanej $b$ i referencyjnej $b^{gt}$,
 \item[$c$] przekątna najmniejszego prostokąta otaczającego obie ramki,
 \item[$v$] człon penalizujący niezgodność współczynnika proporcji (ang. \textit{aspect ratio}),
\end{description}
\begin{equation}
v = \frac{4}{\pi^2} \left(\arctan\frac{w^{gt}}{h^{gt}} - \arctan\frac{w}{h} \right)^2,
\end{equation}
gdzie $\alpha$ jest współczynnikiem wagowym równoważącym wpływ członu $v$ względem członu przecięcia, wyznaczanym adaptacyjnie jako:
\begin{equation}
\alpha = \frac{v}{(1-\mathrm{IoU}) + v}.
\end{equation}
Trzy sygnały w jednej liczbie. Człon $v$ z arcustangensami rozróżnia smukły trzon maślaka od płaskiego kapelusza koźlarza. Dokładnie ta geometria decyduje o parach bliźniaczych. Gdy ramki się nie nakładają, gradient klasycznego IoU zanika; CIoU podtrzymuje sygnał przez $\rho^2/c^2$ i skraca pierwsze, powolne epoki treningu \cite{ultralytics2024yolo11}. Geometrię $\rho$ i $c$ ilustruje rys.~\ref{fig:ciou_geometria}.

\begin{figure}[htbp]
 \centering
 \begin{tikzpicture}[scale=1.0]
 % najmniejszy prostokat otaczajacy obie ramki
 \draw[dashed, gray] (0,0) rectangle (6.2,3.8);
 \draw[dashed, gray, thick] (0,0) -- (6.2,3.8);
 \node[gray, font=\small] at (4.55,1.35) {$c$};

 % ramka referencyjna
 \draw[EEUltramarine, very thick] (0,0) rectangle (3.6,2.6);
 \fill[EEUltramarine] (1.8,1.3) circle (2pt);
 \node[EEUltramarine, font=\small, anchor=south west] at (0.1,0.1) {$b^{gt}$};

 % ramka predykowana
 \draw[EETangerine, very thick] (2.6,1.2) rectangle (6.2,3.8);
 \fill[EETangerine] (4.4,2.5) circle (2pt);
 \node[EETangerine, font=\small, anchor=north east] at (6.1,3.7) {$b$};

 % odleglosc srodkow
 \draw[-latex, semithick] (1.8,1.3) -- (4.4,2.5);
 \node[font=\small, anchor=south east] at (3.35,1.78) {$\rho$};

 % wymiary
 \draw[|-|, gray] (0,-0.35) -- node[below, font=\scriptsize, gray] {$w^{gt}$} (3.6,-0.35);
 \draw[|-|, gray] (-0.35,0) -- node[left, font=\scriptsize, gray] {$h^{gt}$} (-0.35,2.6);
 \draw[|-|, gray] (2.6,4.15) -- node[above, font=\scriptsize, gray] {$w$} (6.2,4.15);
 \draw[|-|, gray] (6.55,1.2) -- node[right, font=\scriptsize, gray] {$h$} (6.55,3.8);
 \end{tikzpicture}
 \caption{Interpretacja geometryczna składników miary CIoU: odległość środków ramek $\rho$, przekątna $c$ najmniejszego prostokąta otaczającego obie ramki oraz wymiary wykorzystywane w członie proporcji $v$. Źródło: opracowanie własne.}
 \label{fig:ciou_geometria}
\end{figure}

Precyzję krawędzi ramki podnosi DFL (\textit{Distribution Focal Loss}): zamiast regresować jedną liczbę, sieć modeluje każdą współrzędną jako rozkład dyskretny na sąsiednich węzłach \cite{ultralytics2024yolo11}:
\begin{equation}
\mathcal{L}_{\mathrm{DFL}}(\mathcal{S}_i, \mathcal{S}_{i+1}) = -\big(y_{i+1}-y\big)\log(\mathcal{S}_i) - \big(y - y_i\big)\log(\mathcal{S}_{i+1}),
\end{equation}
$y$ to rzeczywista współrzędna krawędzi, $y_i$/$y_{i+1}$ – sąsiednie węzły dyskretyzacji, $\mathcal{S}_i$/$\mathcal{S}_{i+1}$ – prognozowane prawdopodobieństwa przynależności. Interpolacja liniowa odtwarza wartość ciągłą z zachowaniem różniczkowalności. CIoU i DFL tworzą człon lokalizacyjny; Binary Cross-Entropy obsługuje klasyfikację gatunkową. Trzy składniki, jedna funkcja celu.

Inferencja: YOLO11m produkuje tysiące nakładających się kandydatów. NMS iteracyjnie wybiera ramkę o najwyższej pewności, odrzuca wszystkie pokrywające się z nią powyżej progu IoU i powtarza. Jeden owocnik. Jedna ramka wynikowa.

\section{Metody optymalizacji modeli dla urządzeń mobilnych}
Model wytrenowany na RTX~4070 musi działać na Snapdragonie bez serwera, bez chłodzenia i bez degradacji po godzinie w terenie. To nie jest proste skalowanie. Poniżej opisano ograniczenia sprzętowe telefonu, logikę kwantyzacji oraz drogę formatu modelu od PyTorch do pliku \texttt{.tflite} w pamięci flash.

\subsection{Ograniczenia sprzętowe urządzeń mobilnych}
Snapdragon nie jest RTX-em. Urządzenie mobilne dysponuje kilkoma gigabajtami RAM, ograniczoną przepustowością magistrali pamięci i nie posiada aktywnego chłodzenia – SoC po kilku minutach intensywnych obliczeń wchodzi w thermal throttling. Model musi zatem być mały i rzadko aktywowany. Stąd decyzja o FP16 i ograniczeniu inferencji do 4~FPS.

\subsection{Kwantyzacja potreningowa}
Podstawową techniką dostosowania wytrenowanego modelu do platformy mobilnej jest kwantyzacja potreningowa (ang. \textit{post-training quantization}, PTQ), polegająca na zastąpieniu 32-bitowej reprezentacji zmiennoprzecinkowej wag i aktywacji reprezentacją o mniejszej liczbie bitów \cite{jacob2018quantization}. Odwzorowanie wartości rzeczywistej na wartość skwantowaną opisuje przekształcenie afiniczne:
\begin{equation}
r = S\,(q - Z),
\end{equation}
gdzie
\begin{description}[topsep=2pt,itemsep=-3pt]
 \item[$r$] pierwotna wartość zmiennoprzecinkowa,
 \item[$q$] odpowiadająca jej wartość całkowitoliczbowa,
 \item[$S$] dodatni współczynnik skali (ang. \textit{scale}),
 \item[$Z$] przesunięcie punktu zerowego (ang. \textit{zero-point}).
\end{description}
Współczynnik skali dla reprezentacji $n$-bitowej wyznacza się na podstawie zakresu wartości pierwotnych zgodnie z zależnością:
\begin{equation}
S = \frac{r_{\max} - r_{\min}}{2^{n} - 1}.
\end{equation}

Kwantyzacja całkowitoliczbowa, opisana powyższym przekształceniem afinicznym, różni się fundamentalnie od redukcji precyzji w obrębie samej reprezentacji zmiennoprzecinkowej. Liczba zmiennoprzecinkowa w formacie IEEE~754 jest reprezentowana jako:
\begin{equation}
x = (-1)^{s} \cdot 2^{\,e-\mathrm{bias}} \cdot \left(1 + \frac{m}{2^{p}}\right),
\end{equation}
gdzie $s$ oznacza bit znaku, $e$ pole wykładnika, $m$ pole mantysy o precyzji $p$~bitów, a $\mathrm{bias}$ jest przesunięciem wykładnika. FP32: $p=23$ ($\approx1{,}19\times10^{-7}$ względnie). FP16: $p=10$ ($\approx9{,}77\times10^{-4}$). W obu przypadkach błąd rośnie z wielkością liczby. INT8 ma stały krok $S$ na cały zakres warstwy. Małe aktywacje giną w kwancie.

Nie jest to wyłącznie kwestia akademicka. Gradienty jasności na blaszkach (\textit{lamellae}), pierścieniu i łuskach kapelusza to właśnie te „małe” sygnały. One oddzielają jadalną czubajkę od trującej czubajeczki i młodego muchomora od pieczarki. INT8 spłaszcza je do stałego kwantu $S$ warstwy. FP16 zachowuje skalowanie wykładnicze, więc te detale nie giną w siatce całkowitoliczbowej.

Dlatego zrezygnowano z pełnej kwantyzacji INT8. Decyzja jest projektowa, nie empiryczna: w pracy nie raportowano $mAP$ ani opóźnienia wariantu INT8. Szacunkowy rozmiar artefaktu INT8 ($\approx20$\,MiB) i zysk na NPU wynikają z szerokości reprezentacji, nie z pomiaru na telefonie. FP16 dla YOLO11m daje $38{,}58$\,MiB – połowa FP32 – przy zachowaniu drobnych cech teksturalnych potrzebnych do par bliźniaczych. Stąd decyzja produkcyjna: tylko FP16.

Różnicę w podziale przestrzeni bitowej pomiędzy trzema rozważanymi reprezentacjami zilustrowano na rys.~\ref{fig:bity_formatow}. Redukcja FP32~$\rightarrow$~FP16 zachowuje strukturę zmiennoprzecinkową i skraca zarówno pole wykładnika, jak i mantysy, natomiast przejście do INT8 całkowicie eliminuje pole wykładnika, zastępując skalowanie wykładnicze pojedynczym, stałym dla całej warstwy współczynnikiem $S$.

\begin{figure}[htbp]
 \centering
 \begin{tikzpicture}[
 pole/.style={draw, minimum height=7mm, align=center, font=\scriptsize, inner sep=1pt},
 znak/.style={pole, fill=black!10},
 wykl/.style={pole, fill=EEBlue!25},
 mant/.style={pole, fill=EETangerine!30},
 calk/.style={pole, fill=EEUltramarine!20}
 ]
 \def\jb{0.34}
 % FP32
 \node[font=\small, anchor=east] at (-0.25,0) {FP32};
 \node[znak, minimum width=\jb cm] (a1) at (\jb/2, 0) {$s$};
 \node[wykl, minimum width=8*\jb cm] (a2) at (\jb+8*\jb/2, 0) {wykładnik: $8$ b};
 \node[mant, minimum width=23*\jb cm] (a3) at (9*\jb+23*\jb/2, 0) {mantysa: $23$ b};
 \node[font=\scriptsize, anchor=west] at (32*\jb+0.2, 0) {$32$ b};

 % FP16
 \node[font=\small, anchor=east] at (-0.25,-1.1) {FP16};
 \node[znak, minimum width=\jb cm] at (\jb/2, -1.1) {$s$};
 \node[wykl, minimum width=5*\jb cm] at (\jb+5*\jb/2, -1.1) {wykł.: $5$ b};
 \node[mant, minimum width=10*\jb cm] at (6*\jb+10*\jb/2, -1.1) {mantysa: $10$ b};
 \node[font=\scriptsize, anchor=west] at (16*\jb+0.2, -1.1) {$16$ b};

 % INT8
 \node[font=\small, anchor=east] at (-0.25,-2.2) {INT8};
 \node[calk, minimum width=8*\jb cm] at (8*\jb/2, -2.2) {liczba całkowita: $8$ b};
 \node[font=\scriptsize, anchor=west] at (8*\jb+0.2, -2.2) {$8$ b $+$ wspólne $S$, $Z$};
 \end{tikzpicture}
 \caption{Podział przestrzeni bitowej w formatach FP32, FP16 oraz INT8. Źródło: opracowanie własne na podstawie \cite{jacob2018quantization}.}
 \label{fig:bity_formatow}
\end{figure}

Ilościowe zestawienie właściwości numerycznych trzech rozważanych reprezentacji wag przedstawiono w tab.~\ref{tab:kwantyzacja}. Zestawienie obejmuje szerokość pola mantysy, wynikającą z niej względną precyzję reprezentacji, charakter propagacji błędu kwantyzacji oraz szacunkowy rozmiar artefaktu modelu YOLO11m w danym formacie. Wpływ kwantyzacji na jakość decyzji klasyfikacyjnych ocenia się w kategoriach macierzy pomyłek, której elementy ($TP$, $FP$, $FN$, $TN$) zdefiniowano na rys.~\ref{fig:conf_matrix_web}; miary te wykorzystano później przy ewaluacji bezpieczeństwa mykologicznego w rozdziale \ref{chap:badania}.

\begin{figure}[htbp]
 \centering
 \includegraphics[width=0.75\textwidth]{confusion_matrix_official}
 \caption{Definicje macierzy pomyłek (TP, FP, FN, TN). Źródło: \cite{fawcett2006roc}.}
 \label{fig:conf_matrix_web}
\end{figure}

\begin{table}[htbp]
 \centering
 \caption{Porównanie formatów reprezentacji numerycznej wag rozważanych na etapie konwersji modelu. Źródło: opracowanie własne na podstawie \cite{jacob2018quantization, tflite}.}
 \label{tab:kwantyzacja}
 \begin{tabular}{|p{4.6cm}|c|c|c|}
 \hline
 \textbf{Właściwość} & \textbf{FP32} & \textbf{FP16} & \textbf{INT8} \\
 \hline \hline
 Szerokość słowa & $32$\,b & $16$\,b & $8$\,b \\
 \hline
 Pole mantysy $p$ & $23$\,b & $10$\,b & brak \\
 \hline
 Precyzja względna & $1{,}19\cdot10^{-7}$ & $9{,}77\cdot10^{-4}$ & krok stały $S$ \\
 \hline
 Skalowanie błędu & proporcjonalne & proporcjonalne & absolutne \\
 \hline
 Rozmiar TFLite (YOLO11m) & ok.\ $77$\,MiB & $38{,}58$\,MiB & ok.\ $20$\,MiB \\
 \hline
 Wymaga kalibracji & nie & nie & tak \\
 \hline
 Zastosowanie w pracy & referencja & \textbf{produkcja} & odrzucony \\
 \hline
 \end{tabular}
\end{table}

\subsection{Formaty pośrednie i docelowe modelu}
PyTorch nie trafia na telefon wprost. Graf obliczeniowy wymaga konwersji do formatu krawędziowego. Powszechnym formatem pośrednim wymiany modeli jest ONNX, który oddziela środowisko treningowe od docelowego i pozwala przełączać się między frameworkami \cite{onnx}; w~niniejszej pracy nie wykorzystano go jednak w~finalnym potoku. Zastosowano prostszą, dwuetapową ścieżkę: Ultralytics eksportuje model do TensorFlow SavedModel, a~następnie konwerter \texttt{tf.lite.TFLiteConverter} kompiluje binarny artefakt TFLite, który Interpreter ładuje w czasie wykonania na urządzeniu \cite{tflite}. Tu stosuje się kwantyzację FP16. Schemat ścieżki SavedModel do TFLite FP16 pokazuje rys.~\ref{fig:tflite_flow_web}; wynikowy artefakt \texttt{.tflite} trafia bezpośrednio do aktywów aplikacji Android (rozdz.~\ref{chap:implementacja}).

\begin{figure}[htbp]
 \centering
 \includegraphics[width=0.9\textwidth]{tflite_conversion_flow_official}
 \caption{Schemat konwersji SavedModel do TFLite FP16. Źródło: \cite{tflite}.}
 \label{fig:tflite_flow_web}
\end{figure}


\chapter{Projekt systemu rozpoznawania}\label{chap:projekt_systemu}
\section{Wymagania funkcjonalne i niefunkcjonalne}
Punkt wyjścia projektu to jeden scenariusz: grzybiarz w Puszczy Białowieskiej, brak LTE, pełny kosz i telefon w dłoni. Wymagania wynikają z tego scenariusza wprost. Funkcjonalne opisują, co aplikacja robi dla użytkownika; niefunkcjonalne określają, ile zasobów można na to zużyć.

\subsection{Wymagania funkcjonalne}
Interakcja użytkownika z systemem zamyka się w kilku przypadkach użycia. Aplikacja musi spełniać pięć wymagań funkcjonalnych:
\begin{enumerate}
 \item Uruchamiać podgląd obrazu z kamery w czasie rzeczywistym jako główny widok roboczy.
 \item Automatycznie wykrywać wiele owocników jednocześnie w kadrze bez ręcznego wskazywania obiektów przez użytkownika.
 \item Nakładać ramki otaczające (BBOX) z etykietami gatunkowymi w wersji polskiej i łacińskiej.
 \item Generować wyraźne, krytyczne ostrzeżenia wizualne przy wykryciu gatunku trującego, wyróżnione kontrastową barwą względem detekcji jadalnych.
 \item Aktualizować wyniki detekcji na bieżąco wraz ze zmianą obserwowanej sceny.
\end{enumerate}
Reguła nadrzędna: priorytet czułości dla klas toksycznych. Ostrzeżenie o gatunku trującym wyprzedza komfort użytkowania. Lepiej oznaczyć jadalną czubajkę kania jako podejrzaną niż przeoczyć muchomora zielonego. Zakres interakcji przedstawia rys.~\ref{fig:przypadki_uzycia}.

\begin{figure}[htbp]
 \centering
 \begin{tikzpicture}[
 pu/.style={draw, ellipse, align=center, font=\scriptsize, minimum width=3.2cm,
 minimum height=1.0cm, inner sep=1pt, fill=EEBlue!10},
 rel/.style={-latex, semithick},
 ext/.style={dashed, -latex, semithick, gray}
 ]
 % aktor
 \draw[semithick] (0,0) circle (0.22);
 \draw[semithick] (0,-0.22) -- (0,-0.9);
 \draw[semithick] (-0.4,-0.45) -- (0.4,-0.45);
 \draw[semithick] (0,-0.9) -- (-0.3,-1.5);
 \draw[semithick] (0,-0.9) -- (0.3,-1.5);
 \node[font=\scriptsize, align=center] at (0,-1.9) {Grzybiarz\\(użytkownik)};

 % granica systemu
 \draw[gray, rounded corners=3pt] (2.0,-4.3) rectangle (10.6,2.1);
 \node[gray, font=\scriptsize, anchor=north west] at (2.15,2.0) {Aplikacja mobilna};

 \node[pu] (skan) at (5.0,0.9) {Skanowanie owocników\\w podglądzie kamery};
 \node[pu] (ostrz) at (5.0,-0.4) {Odbiór ostrzeżenia\\o gatunku trującym};
 \node[pu] (zapis) at (5.0,-1.7) {Zapis znaleziska\\do historii};
 \node[pu] (hist) at (5.0,-3.0) {Przegląd historii\\znalezisk};

 \draw[rel] (0.45,-0.3) -- (skan);
 \draw[rel] (0.45,-0.6) -- (ostrz);
 \draw[rel] (0.45,-1.1) -- (zapis);
 \draw[rel] (0.45,-1.5) -- (hist);

 \node[pu, fill=EETangerine!18] (offl) at (8.6,0.9) {Praca offline};
 \node[pu, fill=EETangerine!18] (kadr) at (8.6,-1.7) {Kadrowanie\\i zapis miniatury};
 \draw[ext] (offl) -- node[right, font=\scriptsize, gray, xshift=1pt] {} (skan);
 \draw[ext] (kadr) -- (zapis);
 \node[gray, font=\scriptsize, anchor=west] at (6.7,1.55) {};
 \end{tikzpicture}
 \caption{Diagram przypadków użycia aplikacji mobilnej; relacje przerywane oznaczają zależności rozszerzające. Źródło: opracowanie własne.}
 \label{fig:przypadki_uzycia}
\end{figure}

\subsection{Wymagania niefunkcjonalne}\label{sec:wymagania_niefunkcjonalne}
Telefon w lesie, bez zasięgu, bateria w połowie. Cztery ograniczenia twarde:
\begin{enumerate}
 \item Praca offline jest bezwzględna: zero odwołań sieciowych; aplikacja działa w martwej strefie GSM/LTE bez degradacji funkcji.
 \item Rozmiar artefaktu pozostaje poniżej $40$\,MiB, tak aby model mieścił się w pamięci flash i nie groził wyrzutem procesu przez Android Low-Memory Killer.
 \item Minimalny poziom API to 26 (Android 8.0), co obejmuje zdecydowaną większość urządzeń w obiegu.
 \item Projektowy niski pobór energii: ograniczenie tempa inferencji do $4$\,FPS ma na celu redukcję nagrzewania i ochronę baterii podczas pracy w terenie. Bezpośredni pomiar poboru mocy nie wchodził w zakres pracy.
\end{enumerate}

Wymagania te uzupełnia definicja czasu rzeczywistego, przyjęta w pracy w znaczeniu użytkowym, a nie w sensie twardych systemów czasu rzeczywistego. Aplikacja działa w czasie rzeczywistym, jeżeli zbieracz kierujący telefon na owocnik otrzymuje odczyt na podglądzie kamery bez zatrzymywania obrazu i w czasie wyraźnie krótszym niż $1$\,s. Kryterium to rozbito na dwie mierzalne wielkości: (i) opóźnienie end-to-end pojedynczej klatki -- od przekazania klatki przez analizator CameraX do opublikowania nowego stanu nakładki -- oraz (ii) czas reakcji -- od nadejścia pierwszej klatki, na której owocnik został wykryty, do wyświetlenia jego ramki lub ostrzeżenia toksycznego. Obie wielkości muszą pozostać poniżej $1$\,s. Zielona asercja „Jadalny” z założenia wymaga zgodności w $4$ z $5$ analizowanych klatek (rozdz.~\ref{chap:implementacja}); jej czas pojawienia się raportowany jest osobno i nie wchodzi do kryterium (ii), ponieważ opóźnienie asercji jadalności jest zamierzonym kosztem bezpieczeństwa.

Cztery ograniczenia i definicja czasu reakcji dają trzy budżety inżynierskie.

\textbf{Budżet opóźnienia}: czas reakcji poniżej $1$\,s przy kadencji analizy $250$\,ms ($4$\,FPS). Kadencja jest docelową częstotliwością analizy klatek, wynikającą z kompromisu między płynnością nakładki a nagrzewaniem SoC i zużyciem baterii; nie jest definicją czasu rzeczywistego. Jeżeli pełny cykl klatki -- konwersja klatki CameraX do bitmapy, letterboxing $640\times640$, \texttt{detect()}, filtr progów i \texttt{DetectionTracker} -- trwa dłużej niż $250$\,ms, \texttt{Mutex.tryLock()} odrzuca kolejne klatki i chwilowa kadencja spada poniżej $4$\,FPS. Nie narusza to kryterium interaktywności, dopóki opóźnienie end-to-end i czas reakcji pozostają poniżej $1$\,s. Na S23 Ultra zmierzono zarówno sam \texttt{detect()} -- średnio $156{,}5$\,ms na zimnym SoC i $245{,}9$\,ms ($p_{95}=278{,}3$\,ms) po obciążeniu (tab.~\ref{tab:latency_s23}) -- jak i pełny cykl end-to-end oraz czas reakcji (podrozdz.~\ref{sec:latency_e2e}).

\textbf{Budżet binarny modelu}: poniżej $40$\,MiB. Artefakt TFLite FP16 YOLO11m zajmuje $38{,}58$\,MiB. Margines $1{,}42$\,MiB chroni przed limitami Android App Bundle i zmniejsza obciążenie transferu przy instalacji w terenie.

\textbf{Budżet RAM}: poniżej $250$\,MB procesu. Składowe: wagi rezydentne ${\approx}39$\,MiB + bufor wejściowy $640{\times}640{\times}3{\times}4$\,B~${\approx}4{,}69$\,MiB + bufor wyjściowy $151{\times}8400{\times}4$\,B~${\approx}4{,}84$\,MiB + bitmapa podglądu CameraX/Compose + narzut ART. Łącznie margines względem domyślnego limitu sterty API~26+ ogranicza ryzyko wyrzutu procesu przez Low-Memory Killer podczas kilkugodzinnej sesji terenowej.

Zbiorcze zestawienie wyprowadzonych budżetów inżynierskich wraz z ich składowymi oraz sposobem weryfikacji przedstawiono w tab.~\ref{tab:budzety}.

\begin{table}[htbp]
 \centering
 \caption{Budżety inżynierskie systemu wyprowadzone z wymagań niefunkcjonalnych. Źródło: opracowanie własne.}
 \label{tab:budzety}
 \begin{tabular}{|p{2.9cm}|c|p{4.4cm}|p{3.1cm}|}
 \hline
 \textbf{Budżet} & \textbf{Limit} & \textbf{Składowe} & \textbf{Weryfikacja} \\
 \hline \hline
 Opóźnienie (czas reakcji) & $<1$\,s & kadencja analizy $250$\,ms ($4$\,FPS); pełny cykl klatki: CameraX, letterbox, \texttt{detect()}, tracker & tab.~\ref{tab:latency_s23} i~\ref{tab:latency_e2e} \\
 \hline
 Rozmiar binarny modelu & $<40$\,MiB & artefakt TFLite FP16 $38$--$39$\,MiB $+$ margines pakietu & pomiar rozmiaru pliku \\
 \hline
 Pamięć operacyjna RAM & $<250$\,MB & wagi $\approx39$\,MiB, bufor wejściowy $4{,}69$\,MiB, bufor wyjściowy $4{,}84$\,MiB, podgląd i ART & analityczne oszacowanie budżetu pamięci \\
 \hline
 \end{tabular}
\end{table}

Trzy budżety, jedno źródło: YOLO11m w formacie TFLite FP16. Dlaczego akurat ten model – uzasadnia kolejny podrozdział.

Sprzęt użyty do treningu i testów mobilnych zestawia tab.~\ref{tab:sprzet}. Stos programowy przedstawiono w tab.~\ref{tab:stos_programowy}. Obliczenia treningowe wykonywano wyłącznie na stacji z CUDA; weryfikację inferencji krawędziowej przeprowadzono na fizycznym urządzeniu mobilnym.

\begin{table}[htbp]
 \centering
 \caption{Specyfikacja środowiska sprzętowego wykorzystanego w pracy. Źródło: opracowanie własne.}
 \label{tab:sprzet}
 \begin{tabular}{|p{3.4cm}|p{4.6cm}|p{4.6cm}|}
 \hline
 \textbf{Komponent} & \textbf{Stacja robocza (trening)} & \textbf{Platforma mobilna (test)} \\
 \hline \hline
 Jednostka centralna & AMD Ryzen~5~3600, $6$/$12$, $3{,}60$\,GHz & Qualcomm Snapdragon~8 Gen~2 \\
 \hline
 Pamięć operacyjna & $32$\,GB DDR4 $3200$\,MHz & $12$\,GB LPDDR5X \\
 \hline
 Akcelerator & NVIDIA GeForce RTX~4070, $12$\,GB VRAM & Adreno~740 (delegat GPU TFLite, z rezerwą CPU) \\
 \hline
 Pamięć masowa & $1$\,TB NVMe PCIe Gen4 & pamięć wewnętrzna UFS \\
 \hline
 System operacyjny & Windows~11 & Android~16 \\
 \hline
 \end{tabular}
\end{table}

\begin{table}[htbp]
 \centering
 \caption{Stos programowy zastosowany na poszczególnych etapach realizacji projektu. Źródło: opracowanie własne.}
 \label{tab:stos_programowy}
 \begin{tabular}{|p{3.3cm}|p{5.0cm}|p{4.3cm}|}
 \hline
 \textbf{Etap potoku} & \textbf{Technologia} & \textbf{Rola w systemie} \\
 \hline \hline
 Akwizycja danych & Python $3.11$, \texttt{requests.Session} & pobieranie obserwacji iNaturalist \\
 \hline
 Auto-anotacja & GroundingDINO, \texttt{autodistill} & lokalizacja zero-shot owocników \\
 \hline
 Trening & PyTorch $2.x$, CUDA~$12.x$, Ultralytics & uczenie transferowe YOLO11m \\
 \hline
 Konwersja & TensorFlow, \texttt{tf.lite.TFLiteConverter} & eksport SavedModel do TFLite FP16 \\
 \hline
 Inferencja mobilna & TensorFlow Lite Interpreter, GpuDelegate & lokalna detekcja na urządzeniu \\
 \hline
 Interfejs użytkownika & Kotlin, Jetpack Compose, CameraX & podgląd, nakładka AR, gesty \\
 \hline
 Persystencja & Room (SQLite), Kotlin Coroutines & historia znalezisk \\
 \hline
 \end{tabular}
\end{table}

Dobór bibliotek wynika z podziału na stację roboczą (akwizycja, anotacja, trening, konwersja) oraz telefon (inferencja TFLite, interfejs Compose, persystencja Room). Dzięki temu graf uczenia nie trafia na urządzenie mobilne, a aplikacja nie zależy od Pythona w czasie wykonania.

\section{Architektura systemu i dobór sieci neuronowej}
Budżety z poprzedniego podrozdziału wymuszają konkretne decyzje architektoniczne: jeden przebieg sieci, artefakt poniżej $40$\,MiB, praca offline. Poniżej opisano wynikającą z tych decyzji strukturę systemu i model, który je spełnia.

\subsection{Trójwarstwowa architektura systemu}
Trzy moduły, trzy odpowiedzialności. CameraX dostarcza klatki i skaluje je do $640\times640$. Interpreter TFLite uruchamia model FP16 i zwraca surowy tensor. Kotlin odrzuca słabe kandydatury, rysuje ramki AR i wyzwala ostrzeżenie toksykologiczne.

Granica między przygotowaniem modelu a czasem wykonania przebiega na artefakcie \texttt{.tflite}. Karta graficzna pracuje tylko raz: trening i konwersja na stacji roboczej. Telefon wykonuje wyłącznie inferencję i logikę decyzyjną -- żadnego połączenia sieciowego, żadnego serwera. Na rys.~\ref{fig:warstwy} pokazano ten podział.

\begin{figure}[htbp]
 \centering
 \begin{tikzpicture}[
 w/.style={draw, rounded corners=3pt, align=center, font=\small,
 text width=12.4cm, inner sep=5pt},
 modul/.style={draw, rounded corners=2pt, align=center, font=\scriptsize,
 text width=3.6cm, minimum height=8mm, inner sep=2pt},
 strz/.style={-latex, thick}
 ]
 % warstwa 1
 \node[w, fill=black!5] (w1) at (0,0) {\textbf{Warstwa przygotowania danych i uczenia} \\[2pt] {\footnotesize stacja robocza: AMD Ryzen 5 3600, NVIDIA RTX 4070, Python}};
 \node[modul, fill=EEBlue!12] (m11) at (-4.2,-1.15) {Akwizycja\\i walidacja obrazów};
 \node[modul, fill=EEBlue!12] (m12) at (0,-1.15) {Auto-anotacja\\GroundingDINO};
 \node[modul, fill=EEBlue!12] (m13) at (4.2,-1.15) {Trening YOLO11m\\(PyTorch, AMP)};

 % warstwa 2
 \node[w, fill=black!5] (w2) at (0,-2.9) {\textbf{Warstwa konwersji modelu} \\[2pt] {\footnotesize PyTorch $\rightarrow$ TensorFlow SavedModel $\rightarrow$ TensorFlow Lite FP16}};
 \node[modul, fill=EETangerine!22] (m21) at (0,-4.05) {Artefakt \texttt{.tflite} FP16, $38{,}58$\,MiB};

 % warstwa 3
 \node[w, fill=black!5] (w3) at (0,-5.8) {\textbf{Warstwa wykonawcza (urządzenie mobilne)} \\[2pt] {\footnotesize Android, Kotlin, Jetpack Compose, CameraX, TFLite Interpreter}};
 \node[modul, fill=EEUltramarine!14] (m31) at (-4.2,-6.95) {Akwizycja klatek\\CameraX};
 \node[modul, fill=EEUltramarine!14] (m32) at (0,-6.95) {Inferencja\\i dekodowanie tensora};
 \node[modul, fill=EEUltramarine!14] (m33) at (4.2,-6.95) {Prezentacja, ostrzeżenia\\i baza Room};

 \draw[strz] (m12) -- (m13);
 \draw[strz] (m11) -- (m12);
 \draw[strz] (m13) -- ++(0,-0.75) -- (4.2,-2.15) -- (0,-2.15) -- (w2);
 \draw[strz] (m21) -- ++(0,-0.75) -- (0,-5.05) -- (w3);
 \draw[strz] (m31) -- (m32);
 \draw[strz] (m32) -- (m33);
 \end{tikzpicture}
 \caption{Podział systemu na warstwę przygotowania danych i uczenia, warstwę konwersji modelu oraz warstwę wykonawczą urządzenia mobilnego. Źródło: opracowanie własne.}
 \label{fig:warstwy}
\end{figure}
Przepływ danych przez system to jednokierunkowy potok o dziewięciu krokach (rys.~\ref{fig:potok_architektura}). Klasa \texttt{ImageAnalysis} CameraX dostarcza klatki ze strategią \texttt{STRATEGY\_KEEP\_ONLY\_LATEST}. Starsze klatki są odrzucane, zanim zdążą się kolejkować. Bramka czasowa w \texttt{onFrame} przepuszcza klatkę co $250$\,ms ($4$\,FPS) albo co $1{,}5$\,s (\textit{heartbeat}) po ustabilizowaniu sceny, a \texttt{Mutex.tryLock()} odrzuca klatkę, gdy poprzednia inferencja jeszcze trwa. Interpreter TFLite uruchamia model FP16 i zwraca tensor $[1,\,151,\,8400]$: 4 współrzędne ramki i $147$ wyników klas dla każdego z $8400$ punktów predykcyjnych na mapach cech. Fast-fail obcina listę do 50 najwyżej ocenionych kandydatów; NMS ($\mathrm{IoU}>0{,}45$) usuwa nadmiary. Należy zaznaczyć, że zaimplementowany w aplikacji mechanizm NMS (\texttt{MushroomDetector.kt}) jest z zasady niewrażliwy na klasę (ang. \textit{class-agnostic}). Oznacza to, że pętla porównuje wyłącznie IoU geometryczny bez sprawdzania indeksu klasy. Ramka gatunku B o niższym wyniku pewności zostanie bezwzględnie usunięta, jeśli jej IoU z nakładającą się ramką gatunku A przekroczy zadany próg. Jest to celowe uproszczenie ukierunkowane na wydajność mobilną.

Kroki tego potoku -- od akwizycji klatki po nakładkę i zapis -- ilustruje rys.~\ref{fig:potok_architektura}.

\begin{figure}[htbp]
 \centering
 \begin{tikzpicture}[
 blok/.style={draw, rounded corners=2pt, align=center, text width=9.2cm,
 inner sep=3pt, font=\small, minimum height=7.5mm},
 polblok/.style={draw, rounded corners=2pt, align=center, text width=4.3cm,
 inner sep=3pt, font=\footnotesize, minimum height=8.5mm},
 akw/.style={blok, fill=EEBlue!12},
 inf/.style={blok, fill=EEUltramarine!12},
 bez/.style={blok, fill=EETangerine!18},
 poltox/.style={polblok, fill=EETangerine!25},
 poled/.style={polblok, fill=EEBlue!20},
 strz/.style={-latex, thick}
 ]
 \node[akw] (n1) at (0,0) {1. CameraX \texttt{ImageAnalysis} – \texttt{STRATEGY\_KEEP\_ONLY\_LATEST}};
 \node[akw] (n2) at (0,-1.15) {2. Bramka $250$\,ms ($4$\,FPS) lub $1{,}5$\,s; \texttt{tryLock} przy zajętej inferencji};
 \node[inf] (n3) at (0,-2.30) {3. Interpreter TensorFlow Lite – model YOLO11m FP16};
 \node[inf] (n4) at (0,-3.45) {4. Surowy tensor wyjściowy $[1,\,151,\,8400]$};
 \node[inf] (n5) at (0,-4.60) {5. Szybka selekcja top-$50$ + bramka wstępna ($0{,}18$ toksyczne, $0{,}20$ inne)};
 \node[inf] (n6) at (0,-5.75) {6. Tłumienie niemaksymalne (class-agnostic NMS), $\mathrm{IoU}>0{,}45$};
 
 % Rozwidlenie sciezki toksycznej i jadalnej
 \node[poltox] (n7a) at (-2.45,-7.15) {7a. Ścieżka toksyczna:\\ostrzeżenie $\ge0{,}18$\\(bez konsensusu $4$/$5$)};
 \node[poled] (n7b) at (2.45,-7.15) {7b. Ścieżka jadalna:\\wymaga $\ge0{,}60$ oraz\\bufora konsensusu $4$/$5$};
 
 \node[bez] (n8) at (0,-8.55) {8. Filtr prezentacji \texttt{displayStatus} (ostrzeżenie / „Jadalny” / szary NEUTRAL)};
 \node[blok, fill=black!6] (n9) at (0,-9.70) {9. Nakładka AR (Compose) oraz baza Room (\texttt{MushroomFind}, flaga \texttt{voteConfirmed})};

 \draw[strz] (n1) -- (n2);
 \draw[strz] (n2) -- (n3);
 \draw[strz] (n3) -- (n4);
 \draw[strz] (n4) -- (n5);
 \draw[strz] (n5) -- (n6);
 \draw[strz] (n6.south) -- ++(0,-0.2) -| (n7a.north);
 \draw[strz] (n6.south) -- ++(0,-0.2) -| (n7b.north);
 \draw[strz] (n7a.south) |- ++(0,-0.2) -- (n8.north);
 \draw[strz] (n7b.south) |- ++(0,-0.2) -- (n8.north);
 \draw[strz] (n8) -- (n9);
 \end{tikzpicture}
 \caption{Potok przetwarzania klatki obrazu z uwzględnieniem bramkowania wstępnego i asymetrii ścieżek decyzyjnych. Źródło: opracowanie własne.}
 \label{fig:potok_architektura}
\end{figure}

Wstępnie przefiltrowane kandydatury podlegają bramkowaniu pewności na poziomie dekodera ($0{,}18$ dla klas toksycznych, $0{,}20$ dla pozostałych), co eliminuje szum tła przed wejściem do pętli NMS. Następnie ścieżki decyzyjne ulegają asymetrycznemu rozdzieleniu: detekcje gatunków trujących o pewności $\ge 0{,}18$ generują ostrzeżenie bez oczekiwania na konsensus $4$ z $5$ (ramka pojawia się po dwóch zgodnych klatkach). Z kolei asercja jadalności wymaga podwójnego warunku: wysokiego progu pewności ($\ge 0{,}60$) oraz domknięcia bufora konsensusu klatek (\textit{multi-frame voting buffer}, $4$ z $5$ ostatnich klatek w module \texttt{DetectionTracker}). Wszystkie pozostałe predykcje otrzymują szary status neutralny. Wynik trafia równolegle do nakładki AR Jetpack Compose, rysującej ramki i etykiety, oraz – przy ręcznym zapisie – do bazy Room, gdzie encja \texttt{MushroomFind} przechowuje gatunek, kategorię toksyczności, wartość pewności, znacznik czasu i ścieżkę wycinka obrazu. Polityka progów dotyczy warstwy wyświetlania: ostrzeżenie toksyczne od $0{,}18$, asercja jadalności od $0{,}60$; nie implikuje wyższego $R$ detektora na klasach trujących.

\subsection{Uzasadnienie wyboru modelu YOLO11 Medium}
Wybór modelu to czterowymiarowy kompromis: dokładność, rozmiar binarny, opóźnienie inferencji i zużycie energii, wszystko przy spełnieniu budżetów z poprzedniego podrozdziału. Przeanalizowano cztery alternatywy:
\begin{enumerate}
 \item YOLOv8n, reprezentujący wcześniejszą generację rodziny YOLO, o gorszym stosunku dokładności do liczby parametrów niż jego następca \cite{ultralytics2023yolov8}.
 \item Lżejszy wariant YOLO11 Nano (\texttt{YOLO11n}), atrakcyjny ze względu na minimalny rozmiar artefaktu TFLite, lecz o niedostatecznej pojemności reprezentacyjnej w zadaniu rozróżniania wielu gatunków bliźniaczych \cite{ultralytics2024yolo11}.
 \item Detektor MobileNet SSD, o niższej precyzji detekcji, niewystarczającej do rozróżniania gatunków bliźniaczych \cite{liu2016ssd}.
 \item Rodzina EfficientDet, której czas inferencji jest zbyt długi względem założeń pracy w reżimie czasu rzeczywistego.
\end{enumerate}
Wybrano YOLO11 Medium (\texttt{YOLO11m}) \cite{ultralytics2024yolo11}. $\approx20{,}1$~mln parametrów daje pojemność wystarczającą do modelowania subtelnych cech morfologicznych $147$ klas, a konwersja do TFLite FP16 ($38{,}58$~MiB) mieści artefakt w pamięci flash telefonu. Na S23 Ultra zmierzono czas \texttt{detect()} (tab.~\ref{tab:latency_s23}) oraz pełny cykl klatki i czas reakcji (podrozdz.~\ref{sec:latency_e2e}). Ablacja pokazuje, że Nano ma wyraźnie niższe $mAP$ na tym samym teście; osobnej analizy par bliźniaczych dla Nano nie przeprowadzono. Wejście modelu to $640\times640$ pikseli; dekoder zwraca listę detekcji (gatunek, pewność, ramka). Ilościową ocenę skuteczności przedstawia rozdz.~\ref{chap:badania}.


\chapter{Implementacja rozwiązania}\label{chap:implementacja}
\section{Przygotowanie zbioru danych i półautomatyczna anotacja}
Baza obejmuje $147$ klas i nie ma gotowych etykiet lokalizacyjnych. Ręczne rysowanie ramek dla całego zbioru jest nierealne: to kilka miesięcy pracy dla jednej osoby. GroundingDINO rozwiązał ten problem bez ręcznego etykietowania geometrii na pełnej skali: zero-shotowa lokalizacja frazą \texttt{"mushroom"} wygenerowała etykiety automatycznie, a nazwa gatunku brała się z folderu iNaturalist, nie z głowicy modelu. Jakość lokalizacji potwierdzono ręcznym audytem geometrii $1\,065$ ramek oraz spisem pominięć na $250$ kadrach izolowanego testu (podrozdz.~\ref{sec:grounding_dino}). Potok: masowa akwizycja z filtrem \texttt{quality\_grade=research}, auto-anotacja, oczyszczenie duplikatów, podział korpusu uczącego na \texttt{train}/\texttt{val} oraz niezależny zbiór testowy izolowany po identyfikatorze obserwacji \texttt{obs\_id}. Liczebności obrazów i ramek zestawia tab.~\ref{tab:podzial_zbioru}; te same wartości obowiązują w rozdz.~\ref{chap:badania}. Schemat przedstawia rys.~\ref{fig:potok_danych}.

\begin{figure}[htbp]
 \centering
 \begin{tikzpicture}[
 etap/.style={draw, rounded corners=2pt, align=center, font=\scriptsize,
 text width=2.95cm, minimum height=9mm, inner sep=2.5pt},
 src/.style={etap, fill=black!7},
 acq/.style={etap, fill=EEBlue!14},
 ann/.style={etap, fill=EEUltramarine!14},
 cln/.style={etap, fill=EETangerine!20},
 strz/.style={-latex, semithick}
 ]
 \node[src] (api) at (0,0) {iNaturalist REST API\\\texttt{quality\_grade}\\\texttt{=research}};
 \node[acq] (dl) at (3.9,0) {Skrypt \texttt{gather\_data}\\$105$ taksonów};
 \node[acq] (val) at (7.8,0) {Cel $\geq 1000$ zdjęć\\skip / retry HTTP};
 \node[acq] (purge) at (11.7,0) {Rozszerzenie\\do $147$ klas};

 \node[ann] (dino) at (11.7,-2.1) {GroundingDINO\\zero-shot};
 \node[ann] (reg) at (7.8,-2.1) {Rejestr klas\\\texttt{class\_registry}\\ID $0$--$146$};
 \node[cln] (nms) at (3.9,-2.1) {Czyszczenie\\duplikatów ramek};
 \node[cln] (split) at (0,-2.1) {Izolacja testu\\po \texttt{obs\_id}};

 \node[src] (yaml) at (0,-3.9) {\texttt{data.yaml}\\\texttt{nc: 147}};

 \draw[strz] (api) -- (dl);
 \draw[strz] (dl) -- (val);
 \draw[strz] (val) -- (purge);
 \draw[strz] (purge) -- (dino);
 \draw[strz] (dino) -- (reg);
 \draw[strz] (reg) -- (nms);
 \draw[strz] (nms) -- (split);
 \draw[strz] (split) -- (yaml);
 \end{tikzpicture}
 \caption{Potok przygotowania zbioru: akwizycja z iNaturalist, dociągnięcie klas do progu $1000$ zdjęć, rozszerzenie do $147$ klas, auto-anotacja oraz izolacja zbioru testowego po \texttt{obs\_id}. Źródło: opracowanie własne.}
 \label{fig:potok_danych}
\end{figure}

Liczebność podziałów podano w tab.~\ref{tab:podzial_zbioru}. Obrazy i ramki zliczono plik po pliku na fizycznym zbiorze (\texttt{images/} oraz \texttt{labels/} każdego podziału). Korpus uczący (\texttt{train}/\texttt{val}) pochodzi z masowej akwizycji i auto-anotacji. Skrypt weryfikacji przecięcia zliczył w korpusie uczącym $129\,607$ unikalnych \texttt{obs\_id}. Zbiór testowy zbudowano osobno: z iNaturalist pobrano nowe obserwacje, których \texttt{obs\_id} nie występuje w \texttt{train} ani \texttt{val} (skrypt \texttt{harvest\_clean\_test.py}), a ramki wygenerowano tym samym protokołem GroundingDINO. Weryfikacja przecięcia dała $0$ wspólnych \texttt{obs\_id} na $32\,727$ obrazach testowych. Cel operacyjny to do $250$ zdjęć na klasę; dla $23$ taksonów iNaturalist nie dostarczył tylu niewidzianych obserwacji research-grade, więc te klasy weszły do testu w mniejszej liczbie (od $8$ zdjęć dla \textit{Hygrophorus marzuolus} do $217$ dla \textit{Scleroderma verrucosum}). W dalszej części pracy obowiązują wartości z tej tabeli.

\begin{table}[htbp]
 \centering
 \caption{Liczba obrazów i ramek otaczających w podziałach \texttt{train}/\texttt{val}/\texttt{test} ($147$ klas). Test jest niezależnym zbiorem, izolowanym po \texttt{obs\_id}, a nie wycinkiem $10\%$ korpusu uczącego. Źródło: opracowanie własne.}
 \label{tab:podzial_zbioru}
 \begin{tabular}{|l|c|r|r|}
 \hline
 \textbf{Podział} & \textbf{Rola} & \textbf{Obrazy} & \textbf{Ramki} \\
 \hline \hline
 \texttt{train} & uczenie & $287\,423$ & $602\,251$ \\
 \hline
 \texttt{val} & dobór hiperparametrów & $31\,223$ & $64\,648$ \\
 \hline
 \texttt{test} & ewaluacja (izolacja \texttt{obs\_id}) & $32\,727$ & $73\,272$ \\
 \hline
 łącznie & --- & $351\,373$ & $740\,171$ \\
 \hline
 \end{tabular}
\end{table}

\paragraph{Uwaga o podziale train/val.}
Skrypt \texttt{split\_dataset.py} wydzielał losowo, osobno dla każdej klasy dodanej w etapie 2 (identyfikatory $\ge 105$), po $10\%$ plików do \texttt{val} i do pierwotnego \texttt{test}, bez grupowania po \texttt{obs\_id}. Podział pierwszych $105$ klas nie jest odtworzony tym skryptem, a pierwotny \texttt{test} zastąpiono później zbiorem izolowanym. Różne zdjęcia tej samej obserwacji mogą zatem trafić do obu podziałów, co może nieznacznie zawyżać metryki walidacyjne doboru hiperparametrów. Izolacja po \texttt{obs\_id} obowiązuje wyłącznie dla zbioru testowego (weryfikacja: $0$ wspólnych \texttt{obs\_id} z korpusem uczącym). Nie jest to wyciek do ewaluacji końcowej, ale pozostaje ograniczeniem metodologicznym podziału \texttt{train}/\texttt{val}.

\subsection{Akwizycja danych z serwisu iNaturalist}
Zbiór uczący pochodzi z iNaturalist \cite{inaturalist} – platformy citizen science z publicznym interfejsem REST API (\url{https://api.inaturalist.org/v1/observations}). Kluczowy filtr: \texttt{quality\_grade=research}. Obserwacja na tym poziomie przeszła weryfikację przez co najmniej dwóch niezależnych identyfikatorów ze społeczności. Skrypt \texttt{gather\_data.py} odpytuje API stronicowanymi żądaniami dla każdego taksonu z listy $105$ gatunków, podmienia w adresie miniaturę \texttt{square} na wariant \texttt{large} i zapisuje plik lokalnie pod nazwą \texttt{\{obs\_id\}\_\{photo\_id\}.jpg}. Format anotacji lokalizacyjnych powstaje dopiero w kolejnym etapie (GroundingDINO); na tym kroku zbierane są same obrazy. Rys.~\ref{fig:coco_sample_web} pokazuje wyłącznie konwencję ramki otaczającej na przykładzie zbioru COCO \cite{lin2014coco}; nie są to fotografie grzybów z niniejszego zbioru.

\begin{figure}[htbp]
 \centering
 \includegraphics[width=0.9\textwidth]{coco_dataset_sample_official}
 \caption{Przykładowe ramki BBOX ze zbioru COCO. Źródło: \cite{lin2014coco}.}
 \label{fig:coco_sample_web}
\end{figure}

Pobieranie jest sekwencyjne w ramach gatunku: jedna sesja HTTP (\texttt{requests.Session}) obsługuje wszystkie taksony, a między zapisami plików obowiązuje pauza $0{,}3$\,s, żeby nie przeciążać API. Cel operacyjny to $\geq 1000$ niepustych plików na klasę (do $20$ stron $\times$ $50$ obserwacji). Jeśli katalog gatunku już spełnia próg, skrypt pomija takson (\texttt{Skip}) i przechodzi dalej. Dzięki temu przerwany proces można wznawiać bez dublowania pracy.

Błędy sieciowe i kody $429$/$500$/$503$ nie kończą całego przebiegu. Dla każdej strony API przewidziano do trzech ponowień z pauzą $5$\,s. Uszkodzony pojedynczy link do zdjęcia jest pomijany w bloku \texttt{except} – pętla idzie dalej. Po przejściu całej listy skrypt uruchamia kolejny cykl audytu, aż wszystkie $105$ klas osiągną próg; dopiero później zbiór rozszerzono o dodatkowe taksony do finalnych $147$ klas.

Rdzeń mechanizmu akwizycji (sesja HTTP, stronicowanie, wznowienie oraz cel $1000$ zdjęć) przedstawiono na listingu \ref{lst:scraper}.

\begin{lstlisting}[language=Python,caption={Akwizycja obrazów z iNaturalist (\texttt{gather\_data.py}): sesja HTTP, wznowienie pobierania i cel $1000$ zdjęć na gatunek},label={lst:scraper}]
BASE_DATA_DIR = "E:/Kiril_Horobets/mushroom_data"
TARGET_PER_CLASS = 1000
PAGES, PER_PAGE, MAX_RETRIES = 20, 50, 3

def download_species(session, species):
    folder = f"{BASE_DATA_DIR}/{species.replace(' ', '_')}"
    os.makedirs(folder, exist_ok=True)
    existing = count_valid_files(folder)  # pliki > 0 B
    if existing >= TARGET_PER_CLASS:
        print(f"Skip: {species} ({existing} photos)")
        return True

    for page in range(1, PAGES + 1):
        url = (
            "https://api.inaturalist.org/v1/observations"
            f"?taxon_name={species.replace(' ', '+')}"
            f"&has[]=photos&quality_grade=research"
            f"&per_page={PER_PAGE}&page={page}"
        )
        results = []
        for retry in range(MAX_RETRIES):
            response = session.get(
                url, headers={"User-Agent": "Mozilla/5.0"}, timeout=15
            )
            if response.status_code == 200:
                results = response.json().get("results", [])
                break
            if response.status_code in (429, 500, 503):
                time.sleep(5)
        if not results:
            break

        for obs in results:
            for photo in obs.get("photos", []):
                img_url = photo["url"].replace("square", "large")
                path = os.path.join(folder, f"{obs['id']}_{photo['id']}.jpg")
                if os.path.exists(path):
                    continue
                try:
                    data = session.get(img_url, timeout=10).content
                    with open(path, "wb") as f:
                        f.write(data)
                    time.sleep(0.3)
                except Exception:
                    continue
    return count_valid_files(folder) >= TARGET_PER_CLASS
\end{lstlisting}

Funkcja \texttt{download\_species} najpierw liczy niepuste pliki w katalogu gatunku. Przy $\geq 1000$ zwraca \texttt{True} bez kontaktu z API – to warunek wznowienia po przerwie. W przeciwnym razie stronicuje obserwacje (\texttt{quality\_grade=research}, \texttt{has[]=photos}). Adres miniatury \texttt{square} zamieniany jest na \texttt{large}. Istniejący plik o tej samej nazwie nie jest pobierany ponownie.

Pętla \texttt{MAX\_RETRIES} chroni przed chwilowymi błędami serwera ($429$, $500$, $503$). Zapis obrazu i pauza $0{,}3$\,s ograniczają ryzyko limitu zapytań. Główna pętla \texttt{main()} owija to w cykle audytu: dopóki którakolwiek z $105$ klas nie osiągnie progu, skrypt odczekuje $60$\,s i skanuje listę od nowa. Sesja \texttt{requests.Session} utrzymuje połączenia TCP między taksonami.

Zbiorcze zestawienie parametrów skryptu akwizycji przedstawiono w tab.~\ref{tab:scraper}.

\begin{table}[htbp]
 \centering
 \caption{Parametry skryptu \texttt{gather\_data.py}. Źródło: opracowanie własne.}
 \label{tab:scraper}
 \begin{tabular}{|p{5.3cm}|p{3.1cm}|p{4.2cm}|}
 \hline
 \textbf{Parametr} & \textbf{Wartość} & \textbf{Uzasadnienie} \\
 \hline \hline
 Interfejs HTTP & \texttt{requests.Session} & ponowne użycie połączeń TCP \\
 \hline
 Liczba taksonów (lista bazowa) & $105$ & gatunki + look-alikes \\
 \hline
 Stopień weryfikacji obserwacji & \texttt{research} & wyłącznie rekordy zweryfikowane \\
 \hline
 Rekordów na stronę / liczba stron & $50$ / $20$ & do $1000$ obserwacji na takson \\
 \hline
 Cel liczebności klasy & $\geq 1000$ plików & próg \texttt{Skip} / sukces cyklu \\
 \hline
 Ponowienia HTTP & $3\times$ ($5$\,s) & kody $429$, $500$, $503$ \\
 \hline
 Pauza między zdjęciami & $0{,}3$\,s & ograniczenie obciążenia API \\
 \hline
 Katalog wyjściowy & \texttt{E:/Kiril\_Horobets/}\newline\texttt{mushroom\_data} & lokalny dysk stacji uczelnianej \\
 \hline
 Liczba klas po rozszerzeniu & $147$ & finalny rejestr produkcyjny \\
 \hline
 \end{tabular}
\end{table}

\subsection{Otwarto-słownikowa auto-anotacja GroundingDINO i rejestr klas}\label{sec:grounding_dino}
GroundingDINO dostaje tekst, nie indeks klasy. Ontologia \texttt{CaptionOntology(\{"mushroom": "mushroom"\})} mówi modelowi: znajdź obiekt semantycznie zgodny z pojęciem \texttt{mushroom} \cite{liu2023grounding}. Gatunek (\textit{Boletus edulis}, \textit{Amanita muscaria}, cokolwiek) pochodzi z nazwy katalogu iNaturalist, nie z głowicy klasyfikacyjnej. Zero-shot rozdziela dwa zadania: lokalizację owocnika od etykiety taksonomicznej. Predykcje trafiają do formatu YOLO: pięć wartości znormalizowanych – \texttt{class\_id cx cy w h}.

Mapowanie nazw gatunków na identyfikatory numeryczne obsługuje rejestr \texttt{class\_registry.txt}. Funkcja \texttt{get\_or\_create\_id()} czyta istniejące mapowania i przy nowym taksonie przydziela kolejny inkrementowany identyfikator. Reguła \textit{append-only}. Żadnej alfabetycznej renumeracji, która unieważniłaby wcześniejsze etykiety. Rejestr końcowy: $147$ klas ($145$ odrębnych gatunków biologicznych; tab.~\ref{tab:zal2_gatunki}), identyfikatory od 0 do 146, jedno źródło prawdy dla warstwy anotacji, pliku YAML i aplikacji mobilnej.

Ograniczenia VRAM wymuszają przetwarzanie partiami po 300 obrazów (\texttt{BATCH\_SIZE = 300}). Każda partia izoluje się w katalogach przejściowych \texttt{staging\_in} i \texttt{staging\_out} – niepełny przebieg inferencji nie zanieczyszcza zbioru wynikowego, a po przerwie demon wznawia pracę bezpiecznie. Po każdej partii \texttt{gc.collect()} zwalnia obiekty Pythona; dopiero potem \texttt{torch.cuda.empty\_cache()} oddaje bloki CUDA. Kolejność istotna. Potok ma dwustopniowy mechanizm awaryjny: gdy inferencja całej partii zawiedzie, demon przełącza się na tryb pojedynczego elementu (\textit{single-item fallback}). Uszkodzony obraz, który wywołuje wyjątek w trybie izolowanym, trwale usuwa \texttt{os.remove} – nigdy więcej nie blokuje procesu.

Po zakończeniu etykietowania skrypt regeneruje plik \texttt{data.yaml}: synchronizuje ścieżkę zbioru, liczbę klas \texttt{nc: 147} i tablicę nazw z aktualnym stanem rejestru. Konfiguracja przekazywana do Ultralytics zawsze odpowiada fizycznej strukturze katalogów \texttt{images/} i \texttt{labels/}.

Zasadniczą pętlę demona auto-anotacji, obejmującą inicjalizację ontologii, przetwarzanie partiami, mechanizm awaryjny oraz zwalnianie pamięci karty graficznej, przedstawiono na listingu \ref{lst:autolabel}.

\begin{lstlisting}[language=Python,caption={Auto-anotacja zero-shot modelem GroundingDINO z przetwarzaniem partiami i zwalnianiem pamięci VRAM (\texttt{scripts/auto\_label.py})},label={lst:autolabel}]
BATCH_SIZE = 300

base_model = GroundingDINO(ontology=CaptionOntology({"mushroom": "mushroom"}))

for folder_name in os.listdir(BASE_DIR):
 if folder_name in IGNORE_DIRS:
 continue
 src_path = os.path.join(BASE_DIR, folder_name)
 class_id = get_or_create_id(folder_name)
 files_to_process = collect_pending_images(src_path)

 for i in range(0, len(files_to_process), BATCH_SIZE):
 batch_files = files_to_process[i: i + BATCH_SIZE]
 clear_staging()

 for f in batch_files:
 shutil.copy2(os.path.join(src_path, f), os.path.join(STAGING_IN, f))

 try:
 with torch.inference_mode():
 base_model.label(input_folder=STAGING_IN, output_folder=STAGING_OUT)

 generated_labels = index_generated_labels(STAGING_OUT)
 for f in batch_files:
 save_single_label(f, src_path, class_id, generated_labels)

 except Exception as batch_error:
 # Batch failed: retry one file at a time
 print(f"[WARN] Batch failed ({batch_error}), falling back to single-item")
 for single_f in batch_files:
 clear_staging()
 try:
 shutil.copy2(os.path.join(src_path, single_f),
 os.path.join(STAGING_IN, single_f))
 with torch.inference_mode():
 base_model.label(input_folder=STAGING_IN,
 output_folder=STAGING_OUT)
 save_single_label(single_f, src_path, class_id,
 index_generated_labels(STAGING_OUT))
 except Exception as single_error:
 print(f"[ERR] Dropping broken file {folder_name}/{single_f}")
 os.remove(os.path.join(src_path, single_f))
 finally:
 clear_staging()
 gc.collect()
 if torch.cuda.is_available():
 torch.cuda.empty_cache()
\end{lstlisting}

Linia 3 inicjalizuje model nauczyciela raz. Poza pętlą po katalogach. Wielokrotne ładowanie wag do VRAM byłoby kosztowne. Ontologia \texttt{CaptionOntology} mapuje zapytanie \texttt{"mushroom"} na etykietę wynikową; model otwartego słownika lokalizuje owocniki wszystkich $147$ klas jednym pojęciem językowym.

Linia 9 przypisuje identyfikator z rejestru \textit{append-only} na podstawie nazwy katalogu – nie z predykcji GroundingDINO. Model odpowiada wyłącznie za geometrię ramki. \texttt{collect\_pending\_images} (linia 10) pomija obrazy z istniejącymi etykietami: proces jest idempotentny i odporny na przerwania.

Pętla w linii 12 dzieli pliki na partie po 300. \texttt{clear\_staging()} (linia 14) czyści katalogi przejściowe przed każdą partią. Indeksowanie wyników w linii 23 obejmuje wyłącznie etykiety bieżącego przebiegu. \texttt{torch.inference\_mode()} (linia 20) wyłącza budowę grafu gradientów i redukuje zapotrzebowanie na pamięć.

Blok \texttt{except} (linie 27--41) realizuje fallback: wyjątek partii nie przerywa procesu. Pliki wracają w trybie izolowanym – uszkodzony obraz odpada na linii 41, $299$ pozostałych jest zachowanych. Obraz, który zawiedzie w trybie izolowanym, nigdy nie wraca do przetwarzania.

Blok \texttt{finally} (linie 42--46) działa bezwarunkowo. Kolejność wywołań ma znaczenie: \texttt{gc.collect()} zwalnia referencje Pythona do tensorów, dopiero potem \texttt{torch.cuda.empty\_cache()} oddaje bloki CUDA alokatorowi. Odwrotna kolejność nie zwalniałaby bloków wciąż referencjonowanych przez żywe obiekty, co prowadzi do fragmentacji VRAM i awarii demona przy długiej pracy.

Rozkład $147$ klas objętych rejestrem względem kategorii toksykologicznej, zgodny z mapowaniem zaimplementowanym w rejestrze aplikacji mobilnej oraz z wykazem zamieszczonym w załączniku, przedstawiono w tab.~\ref{tab:toksycznosc} oraz, w ujęciu graficznym, na rys.~\ref{fig:rozklad_kategorii}. Gatunki objęte agresywnym progiem ostrzegawczym $0{,}18$ stanowią łącznie $23{,}1\%$ rejestru ($34$ klasy), natomiast pozostałe $76{,}9\%$ podlega progom zachowawczym.

\begin{figure}[htbp]
 \centering
 \begin{tikzpicture}[scale=1.0]
 \def\sk{0.075} % cm na jeden gatunek
 \foreach \y/\n/\etyk/\kol in {
 0/69/{Jadalny}/{EEBlue!45},
 -0.9/44/{Niejadalny}/{black!25},
 -1.8/27/{Trujący}/{EETangerine!70},
 -2.7/7/{Śmiertelnie trujący}/{red!55}}{
 \node[font=\scriptsize, anchor=east] at (-0.15,\y) {\etyk};
 \fill[\kol, draw=black!50] (0,\y-0.28) rectangle (\n*\sk,\y+0.28);
 \node[font=\scriptsize, anchor=west] at (\n*\sk+0.12,\y) {\n};
 }
 % os
 \draw[semithick] (0,-3.15) -- (80*\sk,-3.15);
 \foreach \x in {0,20,40,60,80}{
 \draw[semithick] (\x*\sk,-3.15) -- (\x*\sk,-3.28)
 node[below, font=\scriptsize] {\x};
 }
 \node[font=\scriptsize, anchor=north] at (40*\sk,-3.72) {liczba klas};
 \end{tikzpicture}
 \caption{Liczebność kategorii toksykologicznych w rejestrze $147$ klas. Źródło: opracowanie własne.}
 \label{fig:rozklad_kategorii}
\end{figure}

\begin{table}[htbp]
 \centering
 \caption{Rozkład $147$ klas objętych systemem względem kategorii toksykologicznej. Źródło: opracowanie własne.}
 \label{tab:toksycznosc}
 \begin{tabular}{|l|c|c|p{5.6cm}|}
 \hline
 \textbf{Kategoria} & \textbf{Liczba} & \textbf{Udział} & \textbf{Traktowanie w polityce} \\
 \hline \hline
 Jadalny & $69$ & $46{,}9\%$ & asercja od $0{,}60$ (wymaga konsensusu $4$/$5$, utrwalanego w bazie) \\
 \hline
 Niejadalny & $44$ & $29{,}9\%$ & status neutralny \\
 \hline
 Trujący & $27$ & $18{,}4\%$ & ostrzeżenie od $0{,}18$ \\
 \hline
 Śmiertelnie trujący & $7$ & $4{,}8\%$ & ostrzeżenie od $0{,}18$ \\
 \hline \hline
 \textbf{Razem} & $\mathbf{147}$ & $\mathbf{100{,}0\%}$ & --- \\
 \hline
 \end{tabular}
\end{table}

Do kategorii najwyższego ryzyka zaliczono siedem taksonów: \textit{Amanita phalloides}, \textit{Amanita virosa}, \textit{Amanita pantherina}, \textit{Galerina marginata}, \textit{Cortinarius rubellus}, \textit{Entoloma sinuatum} oraz \textit{Gyromitra esculenta}. Gatunki te stanowią $4{,}8\%$ rejestru. Cytowane źródło dotyczy przede wszystkim zatruć amatoksynowych, z dominującą rolą muchomora sromotnikowego \cite{garcia2015amanita}, a nie łącznego udziału całej siódemki w statystyce Europy Środkowej. Najniższy próg przypisano im ze względu na wysoki koszt błędu na tych klasach, a nie dlatego, że jedno źródło dokumentuje ich łączną większość zatruć.

Automatyczna anotacja zero-shot nie zwalnia z obowiązku kontroli jakości etykiet. Geometria ramki jest jedynym sygnałem lokalizacyjnym przekazywanym dalej do YOLO11m, więc systematyczny błąd nauczyciela przeniósłby się na cały trening. Po zakończeniu przebiegu GroundingDINO przeprowadzono więc ręczny audyt lokalizacji na próbie auto-anotowanych obrazów, etykietowanych tym samym protokołem co korpus. Skontrolowano $N=1\,065$ obrazów. Audyt ocenia jakość nauczyciela, a nie zastępuje ewaluacji detektora na bieżącym harvestcie $32\,727$ zdjęć. Próba nie była równomierna względem rejestru $147$ klas: celowo zwiększono udział taksonów o najwyższym ryzyku toksykologicznym, tak aby błąd nauczyciela na klasach krytycznych dla bezpieczeństwa nie umknął w próbie zdominowanej przez gatunki jadalne. W materiale znalazło się $175$ obrazów klas śmiertelnie trujących oraz $285$ obrazów klas trujących; łącznie $460$ kadrów ($43{,}2\%$ próby) wobec $17{,}7\%$ udziału tych kategorii w ówczesnym rejestrze. Kategorie przypisano według rejestru z dnia audytu (18 sierpnia 2026~r.); późniejsza korekta klasyfikacji $11$ gatunków (31 sierpnia 2026~r.) zmieniła udział klas toksycznych do $23{,}1\%$, ale nie wpływa na ocenę geometrii ramek. Pozostałe $605$ obrazów pochodziło z klas jadalnych i niejadalnych.

Każdą ramkę oceniano binarnie jako poprawną albo niepoprawną względem widocznego pokroju owocnika (habitus: kapelusz oraz -- gdy widoczny -- trzon). Ramkę uznawano za niepoprawną, gdy obrys wciągał mech, ściółkę liściową, sąsiednie drewno albo pomijał istotną część owocnika.

Audyt $1\,065$ kadrów dotyczył \emph{istniejących} ramek: oceniano geometrię tego, co GroundingDINO już narysował. Nie liczono tam owocników bez ramki. $mAP$, $P$ i $R$ YOLO11m na teście są względem etykiet auto-anotacji, nie względem pełnego ręcznego spisu. To ograniczenie ground truth. Uzupełnieniem jest stratyfikowana próba $250$ kadrów izolowanego testu (\texttt{seed}=42; $100$ toksycznych, $80$ jadalnych, $70$ niejadalnych według rejestru z dnia losowania, sprzed korekty z 31 sierpnia 2026~r.; \texttt{sample\_fn\_audit.py}, \texttt{fn\_audit\_gui.py}): na każdym zdjęciu liczono owocniki widoczne okiem, ramki leżące na grzybie, pominięcia (FN) i ramki puste (FP). Wyniki są w \texttt{fn\_audit\_sheet.csv} (tab.~\ref{tab:audyt_fn}).

Wynik inspekcji geometrii: $1\,011$ z $1\,065$ ramek ($94{,}93\%$) oceniono jako poprawne. W $54$ przypadkach ($5{,}07\%$) lokalizacja była wadliwa. Tych $54$ obrazów nie poprawiano ręcznym rysowaniem ramek -- etykiety całego zbioru, w tym nadmiarowe obrysy, poddano następnie automatycznemu czyszczeniu \texttt{clean\_duplicate\_boxes\_v2} (kolejny podrozdział). Audyt nie zastępuje pełnej ręcznej anotacji zbioru. Zestawienie podsumowano w tab.~\ref{tab:audyt_grounding}.

\begin{table}[htbp]
 \centering
 \caption{Wyniki ręcznego audytu lokalizacji GroundingDINO na próbie $1\,065$ auto-anotowanych obrazów. Źródło: opracowanie własne na podstawie \texttt{manual\_audit\_summary.json}.}
 \label{tab:audyt_grounding}
 \begin{tabular}{|l|c|}
 \hline
 \textbf{Wielkość} & \textbf{Wartość} \\
 \hline \hline
 Obrazy poddane audytowi & $1\,065$ \\
 \hline
 Ramki poprawne & $1\,011$ \quad ($94{,}93\%$) \\
 \hline
 Ramki niepoprawne & $54$ \quad ($5{,}07\%$) \\
 \hline
 Obrazy klas śmiertelnie trujących & $175$ \\
 \hline
 Obrazy klas trujących & $285$ \\
 \hline
 Obrazy klas jadalnych i niejadalnych & $605$ \\
 \hline
 Pominięcia owocników (FN nauczyciela) & próba $250$ kadrów, tab.~\ref{tab:audyt_fn} \\
 \hline
 \end{tabular}
\end{table}

W próbie $250$ kadrów doliczono $441$ owocników. Nauczyciel pominął $34$ z nich ($7{,}71\%$; $22$ kadry, $8{,}80\%$ próby). $17$ ramek nie leżało na owocniku. $471$ ramek uznano za leżące na grzybie -- więcej niż $441$ widocznych, bo na części okazów stało kilka niemal identycznych obrysów (artefakt z następnego podrozdziału). W czterech kadrach suma ``ramka na grzybie'' + FN różniła się od liczby widocznych o więcej niż $2$; pozostawiono ocenę audytora (kępy kontra zdublowane ramki). Odsetek FN jest wyższy na klasach toksycznych ($19/200$, $9{,}50\%$) niż na jadalnych ($6/117$, $5{,}13\%$). $mAP$ YOLO11m na teście pozostaje pomiarem względem etykiet auto-anotacji: owocniki pominięte przez nauczyciela nie wchodzą do mianownika czułości detektora.

\begin{table}[htbp]
 \centering
 \caption{Ręczny spis owocników i ramek GroundingDINO na stratyfikowanej próbie izolowanego testu ($N=250$, \texttt{seed}=42). FN -- widoczny owocnik bez ramki; FP -- ramka bez owocnika. Procent FN liczony względem liczby widocznych owocników w kolumnie. Źródło: \texttt{fn\_audit\_sheet.csv}.}
 \label{tab:audyt_fn}
 \begin{tabular}{|l|c|c|c|c|}
 \hline
 \textbf{Wielkość} & \textbf{Jadalne} & \textbf{Niejadalne} & \textbf{Toksyczne} & \textbf{Razem} \\
 \hline \hline
 Obrazy & $80$ & $70$ & $100$ & $250$ \\
 \hline
 Widoczne owocniki & $117$ & $124$ & $200$ & $441$ \\
 \hline
 Ramki na owocniku & $135$ & $136$ & $200$ & $471$ \\
 \hline
 FN (owocnik bez ramki) & $6$ ($5{,}13\%$) & $9$ ($7{,}26\%$) & $19$ ($9{,}50\%$) & $34$ ($7{,}71\%$) \\
 \hline
 FP (ramka bez owocnika) & $7$ & $5$ & $5$ & $17$ \\
 \hline
 Kadry z $\ge 1$ FN & $5$ & $7$ & $10$ & $22$ \\
 \hline
 Kadry z $\ge 1$ FP & $7$ & $4$ & $5$ & $16$ \\
 \hline
 \end{tabular}
\end{table}

\paragraph{Kontaminacja wielogatunkowa.}
Obecny potok anotacji przypisuje wszystkim ramkom wykrytym na danym obrazie etykietę klasy wynikającą z~nazwy folderu iNaturalist. Jeżeli na fotografii gatunku~$A$ widoczny jest również owocnik gatunku~$B$, obie ramki otrzymają etykietę~$A$ (ang. \textit{label noise}). Zjawisko to jest częściowo ograniczone przez charakter danych iNaturalist (dominujący jeden obiekt na zdjęciu, \texttt{quality\_grade=research}), lecz jego ilościowa skala nie została zbadana. Stanowi to ograniczenie metodologiczne; potencjalnym rozwiązaniem byłoby dwuetapowe filtrowanie: wstępna detekcja ramek, a~następnie klasyfikacja każdej wyciętej ramki odrębnym modelem wizerunkowym.

\subsection{Dwukryterialne czyszczenie duplikatów ramek}
Surowe etykiety GroundingDINO zawierają charakterystyczny artefakt: kilka niemal identycznych ramek na tym samym owocniku. Klasyczny NMS z łagodnym progiem IoU zawodzi tu z konkretnego powodu – naturalne kępy opieńki (\textit{Armillaria mellea}) mają wysokie przestrzenne nakładanie sąsiednich osobników. Agresywne tłumienie usunęłoby prawdziwe detekcje w gęstych koloniach. Skrypt \texttt{clean\_duplicate\_boxes\_v2.py} rozwiązuje ten konflikt koniunkcją dwóch kryteriów. Dla pary ramek tej samej klasy w formacie YOLO ze środkami $(c_x, c_y)$ i wymiarami $(w, h)$ definiuje się odległość środków:
\begin{equation}
\mathrm{Center\_Distance}(B_i, B_j)
=
\sqrt{\bigl(c_{x,i}-c_{x,j}\bigr)^{2}+\bigl(c_{y,i}-c_{y,j}\bigr)^{2}},
\end{equation}
a także współczynnik przecięcia względem mniejszej ramki (ang. \textit{Intersection over Smallest}):
\begin{equation}
\mathrm{IoS}(B_i, B_j)
=
\frac{\lvert B_i \cap B_j \rvert}{\min\bigl(\lvert B_i \rvert,\,\lvert B_j \rvert\bigr)}.
\end{equation}
Ramkę uznawano za fantomowy duplikat wyłącznie wtedy, gdy jednocześnie spełnione były dwa warunki: wysokie pokrycie przestrzenne $\bigl(\mathrm{IoU} > 0{,}75$ albo $\mathrm{IoS} > 0{,}85\bigr)$ oraz bliskość geometrycznych środków $\bigl(\mathrm{Center\_Distance} < 0{,}06\bigr)$. Koniunkcja kryterium nakładania z kryterium bliskości środków pozwalała usuwać niemal identyczne kopie tej samej detekcji, jednocześnie zachowując sąsiadujące, lecz odrębne owocniki w naturalnych kępach. Geometryczną różnicę pomiędzy obiema konfiguracjami zilustrowano na rys.~\ref{fig:ios_geometria}.

\begin{figure}[htbp]
 \centering
 \begin{tikzpicture}[
 ramka/.style={draw, thick},
 srodek/.style={circle, fill=EETangerine, inner sep=1.3pt},
 opis/.style={font=\footnotesize, align=center}
 ]
 % Przypadek A: duplikat fantomowy (zagniezdzenie, wspolny srodek)
 \draw[ramka, EEUltramarine] (0,0) rectangle (3.0,2.4);
 \draw[ramka, EEBlue] (0.45,0.4) rectangle (2.6,2.05);
 \node[srodek] at (1.5,1.2) {};
 \node[srodek] at (1.525,1.225) {};
 \node[opis] at (1.5,-0.55) {(a) duplikat fantomowy\\$\mathrm{IoU}\approx0{,}49$, $\mathrm{IoS}\approx0{,}99$\\$\mathrm{Center\_Distance}\approx0{,}00$};
 \node[opis, EETangerine] at (1.5,-1.75) {\textbf{usunięty}};

 % Przypadek B: sasiadujace owocniki w kepie
 \draw[ramka, EEUltramarine] (6.0,0) rectangle (8.4,2.4);
 \draw[ramka, EEBlue] (7.6,0.15) rectangle (10.0,2.25);
 \node[srodek] at (7.2,1.2) {};
 \node[srodek] at (8.8,1.2) {};
 \node[opis] at (8.0,-0.55) {(b) sąsiadujące owocniki\\$\mathrm{IoU}\approx0{,}17$, $\mathrm{IoS}\approx0{,}30$\\$\mathrm{Center\_Distance}$ znaczna};
 \node[opis, EEUltramarine] at (8.0,-1.75) {\textbf{zachowane}};
 \end{tikzpicture}
 \caption{Rozróżnienie duplikatu fantomowego od sąsiadujących owocników przez koniunkcję kryteriów $\mathrm{IoS}$ i $\mathrm{Center\_Distance}$. Źródło: opracowanie własne.}
 \label{fig:ios_geometria}
\end{figure} Procedurę uruchomiono na podziałach \texttt{train} i \texttt{val} oraz na pierwotnym, później wycofanym podziale testowym; wynik: $22\,327$ usuniętych fantomowych ramek w $19\,670$ plikach etykiet. Izolowany zbiór testowy (\texttt{label\_clean\_test.py}) tego czyszczenia nie przeszedł, co tłumaczy zdublowane ramki w audycie pominięć (tab.~\ref{tab:audyt_fn}). Implementację obu miar i warunku koniunkcyjnego pokazuje listing \ref{lst:nms_cleaner}.

\begin{lstlisting}[language=Python,caption={Dwukryterialna eliminacja fantomowych duplikatów ramek (\texttt{scripts/utils/clean\_duplicate\_boxes\_v2.py})},label={lst:nms_cleaner}]
IOU_THRESHOLD = 0.75
IOS_THRESHOLD = 0.85
CENTER_DIST_THRESHOLD = 0.06

def calculate_box_metrics(box1, box2):
 # IoU, IoS, and center distance for two YOLO boxes [cls, cx, cy, w, h]
 b1_x1, b1_x2 = box1[1] - box1[3] / 2, box1[1] + box1[3] / 2
 b1_y1, b1_y2 = box1[2] - box1[4] / 2, box1[2] + box1[4] / 2
 b2_x1, b2_x2 = box2[1] - box2[3] / 2, box2[1] + box2[3] / 2
 b2_y1, b2_y2 = box2[2] - box2[4] / 2, box2[2] + box2[4] / 2

 center_dist = math.sqrt((box1[1] - box2[1]) ** 2 + (box1[2] - box2[2]) ** 2)

 inter_x1, inter_y1 = max(b1_x1, b2_x1), max(b1_y1, b2_y1)
 inter_x2, inter_y2 = min(b1_x2, b2_x2), min(b1_y2, b2_y2)
 inter_area = max(0, inter_x2 - inter_x1) * max(0, inter_y2 - inter_y1)

 b1_area, b2_area = box1[3] * box1[4], box2[3] * box2[4]
 min_area = min(b1_area, b2_area)
 union_area = b1_area + b2_area - inter_area + 1e-6

 iou = inter_area / union_area
 ios = inter_area / (min_area + 1e-6)
 return iou, ios, center_dist

def filter_duplicates(boxes):
 filtered_boxes = []
 removed_in_file = 0

 for box_a in boxes:
 is_duplicate = False
 for box_b in filtered_boxes:
 if box_a[0] == box_b[0]:
 iou, ios, center_dist = calculate_box_metrics(box_a, box_b)
 if center_dist < CENTER_DIST_THRESHOLD and (
 iou > IOU_THRESHOLD or ios > IOS_THRESHOLD
):
 is_duplicate = True
 removed_in_file += 1
 break
 if not is_duplicate:
 filtered_boxes.append(box_a)

 return filtered_boxes, removed_in_file
\end{lstlisting}

Linie 7--10 przeliczają reprezentację YOLO, opisującą ramkę środkiem i wymiarami $(c_x, c_y, w, h)$, na współrzędne narożnikowe wymagane do wyznaczenia pola przecięcia. Przeliczenie wykonywane jest w przestrzeni znormalizowanej $[0,1]$, dzięki czemu progi decyzyjne pozostają niezależne od rozdzielczości obrazu źródłowego – ta sama wartość $\mathrm{Center\_Distance} < 0{,}06$ oznacza to samo kryterium bliskości dla zdjęcia o boku $1024$ i $4096$ pikseli.

Linia 12 realizuje bezpośrednio wzór na odległość środków. Linie 14--16 wyznaczają pole prostokąta przecięcia, przy czym funkcja \texttt{max(0, \ldots)} zeruje wynik dla ramek rozłącznych, zapobiegając ujemnym wartościom pola. Kluczowa dla przyjętego rozwiązania jest linia 23, wyznaczająca współczynnik $\mathrm{IoS}$ jako iloraz pola przecięcia i pola \emph{mniejszej} z ramek. W przeciwieństwie do klasycznego $\mathrm{IoU}$, miara ta przyjmuje wartość bliską jedności również wtedy, gdy mała ramka jest niemal całkowicie zawarta w większej. A taka właśnie konfiguracja charakteryzuje zagnieżdżone duplikaty generowane przez model nauczyciela, dla których $\mathrm{IoU}$ pozostaje niskie z powodu istotnej różnicy pól. Stała $10^{-6}$ w mianownikach (linie 20 i 23) zabezpiecza przed dzieleniem przez zero dla ramek zdegenerowanych.

Właściwa decyzja zapada w warunku z linii 34--36. Koniunkcja kryterium bliskości środków z alternatywą dwóch miar pokrycia oznacza, że wysokie nakładanie samo w sobie nie wystarcza do usunięcia ramki – konieczne jest dodatkowo, aby oba owocniki miały praktycznie wspólny środek geometryczny. Dzięki temu sąsiadujące, silnie stykające się osobniki w kępie opieńki zostają zachowane, mimo że ich wzajemne $\mathrm{IoS}$ bywa wysokie. Porównanie ograniczone jest do ramek tej samej klasy (linia 32), a algorytm ma charakter zachłanny: każda nowa ramka zestawiana jest wyłącznie ze zbiorem już zaakceptowanym, co daje złożoność $O(n^2)$ względem liczby ramek w pliku, akceptowalną przy typowej liczbie kilku--kilkunastu anotacji na obraz.

Wyniki zastosowania procedury do poszczególnych podziałów zbioru zestawiono w tab.~\ref{tab:nms_cleanup}.

\begin{table}[htbp]
 \centering
 \caption{Rezultaty dwukryterialnego czyszczenia duplikatów ramek w podziałach zbioru. Źródło: opracowanie własne.}
 \label{tab:nms_cleanup}
 \begin{tabular}{|p{5.2cm}|c|p{4.8cm}|}
 \hline
 \textbf{Kryterium / wynik} & \textbf{Wartość} & \textbf{Rola} \\
 \hline \hline
 Próg $\mathrm{IoU}$ & $0{,}75$ & pokrycie względem sumy pól \\
 \hline
 Próg $\mathrm{IoS}$ & $0{,}85$ & pokrycie względem mniejszej ramki \\
 \hline
 Próg $\mathrm{Center\_Distance}$ & $0{,}06$ & bliskość środków geometrycznych \\
 \hline
 Warunek usunięcia & \multicolumn{2}{p{9.5cm}|}{$\bigl(\mathrm{IoU}>0{,}75 \;\lor\; \mathrm{IoS}>0{,}85\bigr) \;\land\; \mathrm{Center\_Distance}<0{,}06$} \\
 \hline \hline
 Zmodyfikowane pliki etykiet & $19\,670$ & podziały \texttt{train}, \texttt{val}, \texttt{test} \\
 \hline
 Usunięte ramki fantomowe & $22\,327$ & duplikaty tego samego owocnika \\
 \hline
 Średnio na zmodyfikowany plik & $1{,}13$ & średnia liczba usuniętych ramek \\
 \hline
 \end{tabular}
\end{table}

\section{Proces treningu i optymalizacja modelu}
Po przygotowaniu zbioru danych model YOLO11m trenowano w środowisku PyTorch, a następnie konwertowano do pliku \texttt{.tflite} przeznaczonego do lokalnej inferencji na telefonie.

\subsection{Środowisko uczenia i parametry treningu YOLO11m}
Stacja robocza: AMD Ryzen~5~3600, $32$\,GB DDR4, NVIDIA RTX~4070 ($12$\,GB VRAM), CUDA. Stos programowy: Python $3.11+$, PyTorch $2.x$, CUDA~$12.x$, Ultralytics \cite{pytorch, ultralytics2024yolo11}. Skrypt \texttt{train.py} inicjalizuje model, wczytuje \texttt{data.yaml} i uruchamia trening na GPU bez ręcznej interwencji. Oba etapy uczenia korzystają z tego samego skryptu; różnią się wagami startowymi, liczbą epok, rozmiarem wsadu i liczbą procesów ładujących dane (tab.~\ref{tab:parametry_treningu}).

Architektura bazowa: YOLO11 Medium, $20\,143\,369$ parametrów i złożoność $68{,}3$\,GFLOPs według podsumowania modelu Ultralytics dla $147$ klas \cite{ultralytics2024yolo11}. Wariant Nano nie radzi sobie z $147$ klasami, w tym z gatunkami bliźniaczymi (ablacja w rozdz.~\ref{chap:badania}). Wariantu Large nie uruchamiano: nie ma dla niego pomiaru RAM ani opóźnienia, a większy artefakt wykraczałby poza przyjęty budżet pliku. Medium mieści się w przyjętym budżecie rozmiaru pliku FP16. Wagi startują z COCO (\texttt{yolo11m.pt}) – uczenie transferowe \cite{lin2014coco}: niskopoziomowe filtry krawędzi i tekstur wyuczone na ogólnym zbiorze COCO stanowią punkt startowy. Nie stosowano zamrażania warstw (parametr \texttt{freeze} pozostawiono domyślny, \texttt{None}), więc dostrajane są wszystkie warstwy sieci.

Rozdzielczość wejściowa: $640\times640$ pikseli w obu etapach. Biblioteka Ultralytics domyślnie aktywuje mieszaną precyzję (AMP/FP16), co na rdzeniach Tensor Core układu RTX~4070 skraca czas epoki bez destabilizacji gradientów. Domyślny harmonogram uczenia dba o stopniowe wygaszanie współczynnika w końcowej fazie zbieżności.

Trening podzielono na dwa etapy odpowiadające przyrostowej budowie rejestru klas.

Etap podstawowy (przebieg \texttt{train-6}): pierwsze $105$ klas rejestru (identyfikatory $0$--$104$), \texttt{batch=16}, \texttt{workers=4}, wczesne zatrzymanie z \texttt{patience=40} na metryce walidacyjnej. Brak poprawy w ciągu $40$ epok kończy etap i zachowuje najlepsze wagi.

Etap dostrajania (przebieg \texttt{fine\_147classes\_v2}): pełny zestaw $147$ klas, inicjalizacja wagami \texttt{best.pt} etapu podstawowego. Warstwy o zgodnych wymiarach przejmują wagi etapu podstawowego, natomiast końcowe warstwy klasyfikacyjne głowicy, których wymiar zależy od liczby klas, inicjalizowane są od nowa. Wsad zmniejszono do \texttt{batch=8}, a liczbę procesów ładujących do \texttt{workers=2}, aby zmieścić przebieg w zasobach stacji ($12$\,GB VRAM, $32$\,GB RAM). Wszystkie warstwy pozostają odmrożone. Czas trwania: $15$ epok (\texttt{epochs=15}), bez wczesnego zatrzymania (\texttt{patience=0}). Powtarzalność przebiegu zapewniają domyślne ustawienia \texttt{seed=0} i \texttt{deterministic=True}.

Konfigurację wywołania treningu dostrajającego, wraz z rozstrzyganiem urządzenia obliczeniowego, przedstawiono na listingu \ref{lst:train}.

\begin{lstlisting}[language=Python,caption={Wspólny skrypt treningowy \texttt{scripts/train.py}; wartości etapów podaje CLI (\texttt{-{}-model}, \texttt{-{}-epochs}, \texttt{-{}-batch}, \texttt{-{}-workers}, \texttt{-{}-patience}), a nie domyślne stałe skryptu; zestawienie jest w tab.~\ref{tab:parametry_treningu}},label={lst:train}]
DEFAULT_MODEL = "weights/yolo11m.pt"

def resolve_device(requested: str) -> str:
    cuda_available = torch.cuda.is_available()

    if requested == "auto":
        if cuda_available:
            print(f"[device] GPU: {torch.cuda.get_device_name(0)} -> '0'")
            return "0"
        print("[device] No CUDA, using CPU")
        return "cpu"

    if requested == "cpu":
        print("[device] Forced CPU")
        return "cpu"

    if not cuda_available:
        print(f"[device] Requested '{requested}' but CUDA unavailable, using CPU")
        return "cpu"

    device = "0" if requested == "cuda" else requested
    print(f"[device] Using '{device}'")
    return device

def train(args) -> str:
    device = resolve_device(args.device)
    project = str(Path(args.project).resolve())
    model = YOLO(args.model)

    results = model.train(
        data=args.data,
        epochs=args.epochs,
        imgsz=args.imgsz,
        batch=args.batch,
        device=device,
        workers=args.workers,
        patience=args.patience,
        project=project,
        name=args.name,
        exist_ok=args.exist_ok,
        resume=args.resume,
        plots=True,
        show=False,
        verbose=True,
    )

    save_dir = getattr(results, "save_dir", None)
    if save_dir is None:
        save_dir = os.path.join(project, args.name)
    best_path = os.path.join(str(save_dir), "weights", "best.pt")
    return best_path
\end{lstlisting}

\texttt{resolve\_device} (linie 3--23) mapuje argument CLI na identyfikator urządzenia Ultralytics. Tryb \texttt{auto} bez CUDA (linie 6--11), a także jawne żądanie GPU na maszynie bez CUDA (linie 17--19), przechodzi na CPU zamiast zgłaszać wyjątek – ten sam skrypt działa na stacji treningowej i na maszynie bez akceleratora do testów poprawności potoku. Listing pomija komunikaty diagnostyczne funkcji \texttt{train}.

Wczytanie wag realizuje \texttt{model = YOLO(args.model)}; stała \texttt{DEFAULT\_MODEL = weights/yolo11m.pt} jest jedynie wartością domyślną argumentu, wykorzystywaną w etapie podstawowym. W etapie dostrajania argument \texttt{-{}-model} wskazuje na \texttt{best.pt} z etapu podstawowego (zgodnie z tab.~\ref{tab:parametry_treningu}), a nie na wagi COCO \cite{lin2014coco}. \texttt{model.train} uruchamia uczenie transferowe. Kluczowe hiperparametry:

\texttt{epochs=15} (linia 32) -- stała długość etapu dostrajania. \texttt{batch=8} (linia 34) i \texttt{workers=2} (linia 36) -- wartości dobrane do zasobów stacji. \texttt{patience=0} (linia 37) wyłącza wczesne zatrzymanie.

Pozostałe hiperparametry przyjmują wartości domyślne biblioteki Ultralytics, ponieważ \texttt{train.py} ich nie przekazuje. Automatycznie uruchamiana jest mieszana precyzja (AMP): przebieg w przód i wsteczny liczony jest w FP16 na rdzeniach Tensor Core, a wagi główne i krok optymalizatora pozostają w FP32 (skalowanie gradientu). Przy wsadzie $8$ gradient akumulowany jest z $8$ kolejnych kroków do nominalnego wsadu \texttt{nbs=64}. Ustawienie \texttt{optimizer=auto} pozostawia wybór optymalizatora bibliotece: przy liczbie iteracji przekraczającej $10\,000$, jak w obu etapach, wybierany jest optymalizator z rodziny SGD ze współczynnikiem uczenia $0{,}01$ i momentem $0{,}9$. Harmonogram liniowy (\texttt{cos\_lr=False}) z trzema epokami rozgrzewki wygasza współczynnik uczenia do $\texttt{lr0}\cdot\texttt{lrf}$, a augmentacja mozaikowa jest wyłączana w ostatnich $10$ epokach (\texttt{close\_mosaic=10}).

Zestawienie hiperparametrów obu etapów uczenia przedstawiono w tab.~\ref{tab:parametry_treningu}. Domyślne wartości \texttt{argparse} w \texttt{train.py} (\texttt{epochs=200}, \texttt{batch=32}, \texttt{workers=8}) nie są wartościami wykonanych przebiegów. Przebiegi dostały argumenty z tabeli: wsad $16$ i $4$ procesy na etapie podstawowym oraz wsad $8$, $2$ procesy i $15$ epok na dostrajaniu.

\begin{table}[htbp]
 \centering
 \caption{Hiperparametry dwuetapowego procesu uczenia modelu YOLO11m. Źródło: opracowanie własne.}
 \label{tab:parametry_treningu}
 \begin{tabular}{|p{4.5cm}|c|c|}
 \hline
 \textbf{Parametr} & \textbf{Etap podstawowy} & \textbf{Dostrajanie ($147$ klas)} \\
 \hline \hline
 Wagi inicjalizujące & \texttt{yolo11m.pt} (COCO) & \texttt{best.pt} etapu podstawowego \\
 \hline
 Rozdzielczość wejściowa & $640\times640$ & $640\times640$ \\
 \hline
 Rozmiar wsadu (\texttt{batch}) & $16$ & $8$ \\
 \hline
 Procesy wczytujące (\texttt{workers}) & $4$ & $2$ \\
 \hline
 Zamrożone warstwy (\texttt{freeze}) & brak (domyślnie) & brak (domyślnie) \\
 \hline
 Liczba epok & do $200$ & $15$ (stała) \\
 \hline
 Wczesne zatrzymanie & \texttt{patience=40} & \texttt{patience=0} (wyłączone) \\
 \hline
 Optymalizator & \texttt{auto} (domyślnie) & \texttt{auto} (domyślnie) \\
 \hline
 Precyzja obliczeń & AMP FP16 & AMP FP16 \\
 \hline
 Harmonogram uczenia & liniowy, $3$ epoki rozgrzewki & liniowy, $3$ epoki rozgrzewki \\
 \hline
 $\texttt{lr0}$ / $\texttt{lrf}$ & $0{,}01$ / $0{,}01$ (domyślne) & $0{,}01$ / $0{,}01$ (domyślne) \\
 \hline
 Nominalny wsad (\texttt{nbs}) & $64$ (domyślnie) & $64$ (domyślnie) \\
 \hline
 \texttt{close\_mosaic} / \texttt{seed} & $10$ / $0$ (domyślne) & $10$ / $0$ (domyślne) \\
 \hline
 \end{tabular}
\end{table}

\subsection{Konwersja do formatu TensorFlow Lite}
Najlepszy punkt kontrolny wag (\texttt{best.pt}) musi trafić na telefon jako plik \texttt{.tflite}. Tu pojawia się przeszkoda platformowa: bezpośredni eksport \texttt{format='tflite'} w Ultralytics $8.4+$ korzysta z LiteRT, które kompiluje tylko na Linux~x86\_64 i macOS. Na stacji Windows wywołanie kończy się asercją – bez artefaktu. Rozwiązaniem jest dwuetapowy potok w skrypcie \texttt{scripts/export\_tflite.py}. W pierwszym etapie Ultralytics eksportuje graf do formatu TensorFlow SavedModel, obsługiwanego na wszystkich platformach. W drugim etapie natywny \texttt{tf.lite.TFLiteConverter} kompiluje SavedModel do binarnego pliku \texttt{.tflite} z żądaną precyzją. Cały proces jest powtarzalny na stacji Windows, bez przenoszenia wag na inny system. Oba etapy konwersji pokazuje listing \ref{lst:export}.

\begin{lstlisting}[language=Python,caption={Dwuetapowy eksport PyTorch $\rightarrow$ SavedModel $\rightarrow$ TFLite FP16 (\texttt{scripts/export\_tflite.py})},label={lst:export}]
def _run_export(weights: Path, quant: str, imgsz: int, data: str) -> Path:
    import tensorflow as tf

    model = YOLO(str(weights))

    export_kwargs = {
        "format": "saved_model",
        "imgsz": imgsz,
    }
    if quant == "int8":
        export_kwargs["int8"] = True
        export_kwargs["data"] = data

    saved_model_path = Path(str(model.export(**export_kwargs)))

    converter = tf.lite.TFLiteConverter.from_saved_model(str(saved_model_path))

    if quant == "int8":
        converter.optimizations = [tf.lite.Optimize.DEFAULT]
    elif quant == "fp16":
        converter.optimizations = [tf.lite.Optimize.DEFAULT]
        converter.target_spec.supported_types = [tf.float16]
    # fp32: converter defaults

    tflite_model = converter.convert()

    tflite_path = saved_model_path.parent / f"{saved_model_path.name}_{quant}.tflite"
    tflite_path.write_bytes(tflite_model)

    return tflite_path
\end{lstlisting}

Linia 14: \texttt{model.export} z \texttt{format='saved\_model'} omija LiteRT i generuje graf TensorFlow SavedModel, obsługiwany na każdej platformie. \texttt{imgsz} utrwala kształt tensora wejściowego $640\times640$ statycznie, dzięki czemu interpreter mobilny nie realokuje buforów przy każdej klatce.

Linia 16 tworzy konwerter ze ścieżki SavedModel. Rozgałęzienie w liniach 18--23 konfiguruje precyzję. FP16 (linie 21--22) wymaga dwóch ustawień jednocześnie: \texttt{Optimize.DEFAULT} uruchamia kwantyzację wag, a \texttt{target\_spec.supported\_types = [tf.float16]} wymusza połowiczną precyzję zamiast domyślnej INT8. Gałąź INT8 (linie 18--19) ustawia jedynie \texttt{Optimize.DEFAULT} bez reprezentatywnego zbioru danych, co daje kwantyzację wag z dynamicznym zakresem, a nie pełną, kalibrowaną kwantyzację INT8; argument \texttt{data} (linie 10--12) trafia wyłącznie do eksportu Ultralytics. Wariant INT8 odrzucono z powodów opisanych wcześniej.

\texttt{converter.convert()} (linia 25) kompiluje graf do formatu FlatBuffer, który linie 27--28 zapisują na dysk. Artefakt: \texttt{best\_saved\_model\_fp16.tflite}, $38{,}58$\,MiB – mieści się w budżecie $40$\,MiB (tab.~\ref{tab:budzety}).

Strukturę tensora wyjściowego skompilowanego modelu przedstawiono na rys.~\ref{fig:tensor_geometria}, a odwzorowanie kanałów na wielkości fizyczne zestawiono w tab.~\ref{tab:tensor}.

\begin{figure}[htbp]
 \centering
 \begin{tikzpicture}[
 opis/.style={font=\footnotesize, align=center},
 strz/.style={-latex, semithick}
 ]
 % Kanaly (os pionowa)
 \fill[EETangerine!35] (0,0) rectangle (10.5,0.55);
 \draw (0,0) rectangle (10.5,0.55);
 \node[opis] at (5.25,0.275) {kanały $0$--$3$: współrzędne ramki $(\mathit{cx},\,\mathit{cy},\,w,\,h)$};

 \fill[EEUltramarine!22] (0,-2.35) rectangle (10.5,-0.05);
 \draw (0,-2.35) rectangle (10.5,-0.05);
 \node[opis] at (5.25,-1.2) {kanały $4$--$150$: $147$ wyników klas};

 % Os pozioma - punkty predykcyjne
 \draw[strz] (0,-2.75) -- (10.5,-2.75);
 \node[opis] at (5.25,-3.2) {$8400$ punktów predykcyjnych na mapach cech};

 \draw[dashed, black!45] (5.0,-2.65) -- (5.0,0.65);
 \draw[dashed, black!45] (7.5,-2.65) -- (7.5,0.65);
 \node[opis] at (2.5,0.95) {siatka $80\times80$\\($6400$)};
 \node[opis] at (6.25,0.95) {$40\times40$\\($1600$)};
 \node[opis] at (9.0,0.95) {$20\times20$\\($400$)};

 % Wymiar pionowy
 \draw[latex-latex, semithick] (-0.45,-2.35) -- (-0.45,0.55);
 \node[opis, rotate=90] at (-0.85,-0.9) {$151$ kanałów};
 \end{tikzpicture}
 \caption{Geometria tensora wyjściowego $[1,\,151,\,8400]$ z podziałem punktów predykcyjnych na trzy skale map cech (architektura \textit{anchor-free}). Źródło: opracowanie własne.}
 \label{fig:tensor_geometria}
\end{figure}

\begin{table}[htbp]
 \centering
 \caption{Geometria tensora wyjściowego modelu TFLite FP16 o kształcie $[1,\,151,\,8400]$. Źródło: opracowanie własne.}
 \label{tab:tensor}
 \begin{tabular}{|c|c|p{7.3cm}|}
 \hline
 \textbf{Wymiar} & \textbf{Rozmiar} & \textbf{Interpretacja} \\
 \hline \hline
 $0$ & $1$ & rozmiar wsadu (pojedyncza klatka) \\
 \hline
 $1$ & $151$ & $4$ współrzędne ramki $(\mathit{cx},\mathit{cy},w,h)$ oraz $147$ wyników klas (sigmoid) \\
 \hline
 $2$ & $8400$ & punkty predykcyjne na mapach cech: $80^2 + 40^2 + 20^2 = 6400+1600+400$ \\
 \hline \hline
 \multicolumn{3}{|l|}{Indeksowanie w dekoderze: $\mathit{idx} = \text{kanał} \cdot 8400 + \text{punkt}$} \\
 \hline
 \multicolumn{3}{|l|}{Rozmiar bufora wyjściowego: $151 \cdot 8400 \cdot 4$\,B $\approx 4{,}84$\,MiB (Float32)} \\
 \hline
 \end{tabular}
\end{table}

Pełnej kwantyzacji INT8 nie wdrożono. Artefakt produkcyjny to \texttt{best\_saved\_model\_fp16.tflite}, $38{,}58$\,MiB. INT8 schodziłby do około $20$\,MiB i potencjalnie szybciej na NPU, lecz stały kwant $S$ spłaszcza gradienty jasności na blaszkach, pierścieniu i łuskach kapelusza -- cechach, które odróżniają jadalną czubajkę od trującej czubajeczki oraz borowika od goryczaka. W pracy nie mierzono $mAP$ wariantu INT8; rezygnacja wynika z własności reprezentacji (rozdz.~\ref{chap:przeglad_rozwiazan}), nie z nieprzeprowadzonej ablacji. FP16 daje niemal dwukrotną redukcję wobec FP32 przy zachowaniu struktury zmiennoprzecinkowej. Decyzja produkcyjna: tylko FP16.

Tensor wyjściowy modelu ma kształt $[1,\;151,\;8400]$: $8400$ punktów predykcyjnych na trzech skalowych mapach cech ($80{\times}80$, $40{\times}40$, $20{\times}20$), $151$ kanałów. Cztery współrzędne ramki i $147$ wyników klas. Geometria jest zgodna z dekoderem po stronie Androida. Indeksy klas mapuje na gatunki statyczna tablica \texttt{MushroomRegistry.speciesMap}, której kolejność odpowiada rejestrowi \texttt{class\_registry.txt} z etapu treningu. Skompilowany model FP16 trafia do aktywów aplikacji mobilnej jako \texttt{best\_float16.tflite}.

\section{Implementacja aplikacji mobilnej Android}
Aplikacja mobilna to natywny Kotlin na Androida. Trzy warstwy: deklaratywny interfejs w Jetpack Compose, akwizycja klatek przez CameraX i warstwa ML z interpreterem TFLite, logiką śledzenia i filtrem bezpieczeństwa. Granice między warstwami są ostre – każdą można testować niezależnie. Razem realizują jeden cel: lokalna inferencja na żywo, bez serwera, w lesie bez LTE.

\subsection{Architektura interfejsu w Jetpack Compose i CameraX}
Ekran główny to Compose: \texttt{MushroomCameraScreen} trzyma uprawnienia, podgląd i nakładkę; \texttt{DetectionOverlay} rysuje ramki osobno od kamery \cite{jetpackcompose}. CameraX dostarcza przypadki użycia z rys.~\ref{fig:camerax_web}: \texttt{Preview} zasila \texttt{PreviewView} przez \texttt{setSurfaceProvider(...)}, a \texttt{ImageAnalysis} podłącza analizator przez \texttt{setAnalyzer(...)} \cite{camerax}. Klatki do TFLite dostarcza \texttt{ImageAnalysis}. \texttt{ImageCapture} służy wyłącznie ukrytemu Trybowi Badawczemu do zapisu pełnorozdzielczego zdjęcia (podrozdz.~\ref{sec:field_quant}); \texttt{VideoCapture} nie jest używany.

\begin{figure}[htbp]
 \centering
 \includegraphics[width=0.9\textwidth]{camerax_arch_official}
 \caption{Architektura biblioteki CameraX z modułem ImageAnalysis. Źródło: \cite{camerax}.}
 \label{fig:camerax_web}
\end{figure}

\texttt{ImageAnalysis.Builder} ustawia strategię \texttt{STRATEGY\_KEEP\_ONLY\_LATEST} (w buforze zostaje tylko najnowsza klatka) oraz format wyjściowy \texttt{RGBA\_8888}. Podgląd działa w trybie \texttt{FILL\_CENTER}. Analizator pracuje na własnym wątku: każdą klatkę zamienia na mapę bitową obróconą do pionu i przekazuje do \texttt{MushroomViewModel.onFrame}. Tam \texttt{FrameCropper.letterbox} wycina obszar widoczny na ekranie i wpisuje go z zachowaniem proporcji w kwadrat $640\times640$, uzupełniając marginesy szarym kolorem (\texttt{Color.GRAY}, wartość $136$).

Analiza każdej z około $30$ klatek na sekundę dostarczanych przez kamerę nie jest potrzebna i niepotrzebnie nagrzewałaby układ SoC. W \texttt{MushroomViewModel} inferencja uruchamiana jest więc co najwyżej co $250$\,ms ($4$\,FPS), a \texttt{Mutex.tryLock()} przepuszcza co najwyżej jedną inferencję naraz. Gdy wszystkie potwierdzone tory są zamrożone (\texttt{isSceneLocked()}), interwał rośnie do $1{,}5$\,s (\textit{heartbeat}), ponieważ obraz tych samych owocników nie wymaga ciągłej analizy. Listing \ref{lst:throttling} pokazuje obie bramki.

\begin{lstlisting}[language=Kotlin,caption={Ograniczanie częstotliwości inferencji i tryb ekonomiczny (\texttt{MushroomViewModel.kt}; pominięto pomiar czasu dla profilera i estymację FPS kamery)},label={lst:throttling}]
private val inferenceMutex = Mutex()
private var lastInferenceAt = 0L

// Written by the inference coroutine; the tracker itself is not thread-safe.
@Volatile
private var lastSceneLocked = false

fun onFrame(
    upright: Bitmap,
    arrivedNs: Long = SystemClock.elapsedRealtimeNanos(),
    sensorTimestampNs: Long = 0L
) {
    val activeDetector = detector
    if (activeDetector == null) {
        upright.recycle()
        return
    }

    val now = SystemClock.elapsedRealtime()
    val sceneLocked = lastSceneLocked
    val interval = if (sceneLocked && !InferenceProfiler.bypassSceneLock) {
        LOCK_HEARTBEAT_MS
    } else {
        INFERENCE_INTERVAL_MS
    }
    if (now - lastInferenceAt < interval) {
        upright.recycle()
        return
    }
    // Skip if a previous inference is still running (keep only the latest frame).
    if (!inferenceMutex.tryLock()) {
        upright.recycle()
        return
    }
    lastInferenceAt = now

    val aspect = viewAspectRatio
    viewModelScope.launch(Dispatchers.Default) {
        try {
            val lb = FrameCropper.letterbox(upright, aspect, activeDetector.inputResolution)
            val rawDetections = activeDetector.detect(lb.bitmap)
            val detections = rawDetections
                .map { it.copy(boundingBox = FrameCropper.toVisibleNormalized(it.boundingBox, lb)) }
                // Asymmetric floor: toxic @ 0.18, other classes @ 0.20.
                .filter { it.confidence >= MushroomRegistry.displayThresholdFor(it.species.classId) }
            lb.bitmap.recycle()

            val tracked = tracker.update(detections, now)
            val locked = tracker.isSceneLocked()
            lastSceneLocked = locked
            currentDetections = tracked

            synchronized(frameLock) {
                latestFrame?.let { if (it !== upright) it.recycle() }
                latestFrame = upright
            }

            _uiState.value = MushroomUiState.Scanning(tracked, locked)
        } catch (t: Throwable) {
            Log.e(TAG, "Inference failed for frame", t)
            upright.recycle()
        } finally {
            inferenceMutex.unlock()
        }
    }
}

companion object {
    /** Inference cadence: 4 FPS (one frame every 250 ms). */
    private const val INFERENCE_INTERVAL_MS = 250L

    /** Slow heartbeat used while the scene is locked, to conserve battery. */
    private const val LOCK_HEARTBEAT_MS = 1_500L
}
\end{lstlisting}

Linie 20--25 realizują adaptacyjny dobór interwału. Flaga \texttt{lastSceneLocked} przechowuje wynik \texttt{tracker.isSceneLocked()} z ostatniej inferencji (linie 49--50); odczyt flagi zamiast bezpośredniego wywołania trackera z wątku analizatora eliminuje współbieżny dostęp do struktury, która nie jest bezpieczna wątkowo. Scena jest zablokowana, gdy wszystkie potwierdzone tory są dopasowywane nieprzerwanie od co najmniej $5$\,s (\texttt{FREEZE\_MS}), co w praktyce odpowiada kamerze utrzymywanej nad tymi samymi owocnikami. Interwał wydłuża się wtedy z $250$\,ms do $1500$\,ms, sześciokrotnie zmniejszając liczbę wywołań interpretera. W kompilacji debug, używanej do profilowania (rozdz.~\ref{chap:badania}), flaga \texttt{InferenceProfiler.bypassSceneLock} wyłącza ten tryb, aby pomiar odbywał się przy stałej kadencji $4$\,FPS. Wykorzystano zegar \texttt{SystemClock.elapsedRealtime()} (linia 19), monotoniczny i odporny na zmiany czasu systemowego, w przeciwieństwie do \texttt{System.currentTimeMillis()}.

Linie 26--29 stanowią pierwszą bramkę: klatka nadchodząca przed upływem interwału jest natychmiast zwalniana wywołaniem \texttt{recycle()}. Jawne zwalnianie map bitowych jest tu konieczne, ponieważ CameraX dostarcza klatki z częstotliwością rzędu $30$\,FPS, a przy $4$\,FPS inferencji około $87\%$ z nich zostaje odrzuconych. Konwersja klatki do mapy bitowej następuje jeszcze w analizatorze, przed tą bramką, więc bramka oszczędza letterboxing i inferencję, ale nie samą konwersję.

Linie 31--34 realizują drugą bramkę, opartą na wzajemnym wykluczaniu. Metoda \texttt{tryLock()} obiektu \texttt{Mutex} z biblioteki \texttt{kotlinx.coroutines} nie zawiesza wywołującego: przy zajętej blokadzie natychmiast zwraca \texttt{false}. Metoda \texttt{onFrame} wykonywana jest na wątku analizatora CameraX, więc oczekiwanie na zwolnienie blokady opóźniałoby kolejne klatki. W danej chwili wykonywana jest co najwyżej jedna inferencja, a klatka, która nadejdzie w jej trakcie, jest porzucana.

Zwolnienie blokady umieszczono w bloku \texttt{finally} (linie 62--64), dzięki czemu wyjątek zgłoszony podczas inferencji nie blokuje potoku na stałe. Po udanej inferencji klatka nie jest zwalniana, lecz zastępuje poprzednią jako \texttt{latestFrame} (linie 53--56), z której korzysta ręczny zapis znaleziska. Filtr w linii 45 powtarza progi dekodera ($0{,}18$/$0{,}20$); pozostawiono go jako zabezpieczenie na wypadek zmiany progów w samym dekoderze.

Wyniki trafiają na płótno Jetpack Compose (\texttt{Canvas}) w komponencie \texttt{DetectionOverlay}. Współrzędne ramek, znormalizowane po \texttt{FrameCropper.toVisibleNormalized} względem obszaru widocznego na ekranie (linia 43), mapują się na piksele podglądu w geometrii \texttt{FILL\_CENTER}, dzięki czemu nakładka pokrywa się z obrazem kamery niezależnie od proporcji ekranu. Kolor obrysu i etykiety wyznacza \texttt{MushroomRegistry.displayStatus}: pomarańczowy dla gatunków trujących, czerwony dla śmiertelnie trujących, zielony dla jadalnych (tylko po konsensusie) oraz neutralny szary dla wszystkich pozostałych przypadków, w tym gatunków niejadalnych. Gesty dotykowe obsługuje \texttt{modifier.pointerInput} z \texttt{detectTapGestures}: po dotknięciu hit-test przestrzenny wskazuje aktywny prostokąt \texttt{Rect} i uruchamia ręczny zapis wybranej detekcji do historii znalezisk.

\subsection{Dekodowanie tensorów i dwustopniowy filtr bezpieczeństwa}
Interpreter TFLite zwraca tensor $[1,\;151,\;8400]$ w transponowanym układzie Ultralytics \cite{ultralytics2024yolo11}: $8400$ punktów predykcyjnych na mapach cech, $151$ kanałów. Cztery współrzędne ramki i $147$ wyników klas. Głowica YOLO11 jest \textit{anchor-free}: te $8400$ lokacji to komórki siatek $80{\times}80$, $40{\times}40$ i $20{\times}20$, a nie klasyczne \textit{anchor boxes}. Dekoder przegląda wszystkie punkty, wyznacza klasę o najwyższym wyniku i odrzuca kandydatów poniżej progu – globalnego i klasowo-zależnego – zanim utworzy dla nich jakikolwiek obiekt.

Do NMS trafia co najwyżej $50$ najwyżej ocenionych kandydatów (\texttt{MAX\_NMS\_INPUT=50}). Zachłanny NMS usuwa nakładające się ramki powyżej progu IoU \cite{redmon2016yolo}. Górny pułap wejść ogranicza kwadratowy koszt NMS przy gęstych scenach. Interpreter korzysta z delegata GPU, jeżeli \texttt{CompatibilityList} uznaje urządzenie za zgodne; w przeciwnym razie działa na CPU w czterech wątkach. Model jest mapowany do pamięci bezpośrednio z aktywów aplikacji (TensorFlow Lite $2.14$). Dekoder i NMS implementuje listing \ref{lst:decoder}.

\begin{lstlisting}[language=Kotlin,caption={Dekodowanie tensora \texttt{[1, 151, 8400]} z szybką eliminacją i tłumieniem niemaksymalnym (\texttt{MushroomDetector.kt}; pominięto zbieranie pięciu najlepszych klas dla Trybu Badawczego)},label={lst:decoder}]
private fun decode(data: FloatArray): List<RawDetection> {
    val detections = ArrayList<RawDetection>()
    val globalMin = MushroomRegistry.minSafetyThreshold

    for (anchor in 0 until numAnchors) {
        var bestClass = -1
        var bestScore = 0f

        for (c in 0 until numClasses) {
            val score = data[(BOX_COORDS + c) * numAnchors + anchor]
            if (score > bestScore) {
                bestScore = score
                bestClass = c
            }
        }

        // Cheap global gate first, then the per-class safety threshold.
        if (bestClass < 0 || bestScore < globalMin) continue
        if (bestScore < MushroomRegistry.thresholdFor(bestClass)) continue

        val cx = data[0 * numAnchors + anchor]
        val cy = data[1 * numAnchors + anchor]
        val w = data[2 * numAnchors + anchor]
        val h = data[3 * numAnchors + anchor]
        val half = 0.5f

        detections.add(
            RawDetection(
                x1 = cx - w * half,
                y1 = cy - h * half,
                x2 = cx + w * half,
                y2 = cy + h * half,
                score = bestScore,
                classIndex = bestClass
            )
        )
    }
    return detections
}

private fun nonMaxSuppression(input: List<RawDetection>): List<RawDetection> {
    if (input.isEmpty()) return emptyList()

    val sorted = input.sortedByDescending { it.score }
        .take(MAX_NMS_INPUT)
        .toMutableList()
    val kept = ArrayList<RawDetection>()

    while (sorted.isNotEmpty()) {
        val best = sorted.removeAt(0)
        kept.add(best)
        if (kept.size >= MAX_DETECTIONS) break

        val iterator = sorted.iterator()
        while (iterator.hasNext()) {
            if (iou(best, iterator.next()) > IOU_THRESHOLD) iterator.remove()
        }
    }
    return kept
}

private const val IOU_THRESHOLD = 0.45f
private const val MAX_NMS_INPUT = 50
private const val MAX_DETECTIONS = 20
\end{lstlisting}

Pętla zewnętrzna (linia 5) przebiega kolejno wszystkie $8400$ punktów predykcyjnych na mapach cech (w kodzie indeks \texttt{anchor}), natomiast pętla wewnętrzna (linia 9) wyszukuje wśród $147$ kanałów klas wynik maksymalny. Adresowanie w linii 10 odzwierciedla transponowany układ tensora: element o indeksie kanału $c$ i punktu $a$ znajduje się pod pozycją $(4+c)\cdot 8400 + a$, ponieważ dane zapisane są kanałami, a nie punktami przestrzennymi. Bezpośrednia praca na spłaszczonej tablicy \texttt{FloatArray} zamiast na strukturach wielowymiarowych eliminuje pośrednictwo obiektów opakowujących w ścieżce wykonywanej $8400 \cdot 147 \approx 1{,}23$ mln razy na klatkę.

Kluczowa dla wydajności jest kolejność bramek w liniach 18--19. Najpierw wykonywane jest porównanie z progiem globalnym \texttt{globalMin} (równym $0{,}18$), będące pojedynczym porównaniem zmiennoprzecinkowym; dopiero kandydat, który je przejdzie, podlega kosztowniejszemu odpytaniu rejestru o próg właściwy dla swojej klasy. Odwrócenie tej kolejności oznaczałoby odczyt z mapy dla każdego z $8400$ punktów, podczas gdy w praktyce warunek z linii 18 odrzuca zdecydowaną większość z nich. Obiekt \texttt{RawDetection} (linia 28) powstaje dopiero po przejściu obu bramek, więc klatka bez owocników nie tworzy obiektów kandydatów. Stałym kosztem każdej klatki pozostają natomiast kopia tensora wyjściowego do tablicy \texttt{FloatArray} ($\approx4{,}84$\,MiB) oraz mapy bitowe wejścia.

Metoda \texttt{nonMaxSuppression} (linie 41--60) realizuje szybką eliminację oraz właściwe tłumienie. Wywołanie \texttt{take(MAX\_NMS\_INPUT)} w linii 45 ogranicza wejście algorytmu do $50$ najwyżej ocenionych kandydatów. Ponieważ złożoność zachłannego NMS jest kwadratowa względem liczby wejść, ograniczenie to zamienia potencjalnie duży koszt na górny pułap rzędu $50^2/2 = 1250$ porównań, także wtedy, gdy tekstura tła generuje setki słabych detekcji; samo sortowanie kandydatów nadal obejmuje wszystkich kandydatów, którzy przeszli bramki. Pętla główna (linie 49--58) w każdej iteracji przenosi kandydata o najwyższym wyniku do zbioru zaakceptowanego, a następnie usuwa z listy roboczej wszystkie ramki wykazujące z nim $\mathrm{IoU} > 0{,}45$. Iterator z metodą \texttt{remove()} (linie 54--57) pozwala usuwać elementy w trakcie przechodzenia kolekcji bez tworzenia list pośrednich. Ogranicznik \texttt{MAX\_DETECTIONS} (linia 52) przerywa pętlę po $20$ zachowanych detekcjach, co odpowiada praktycznemu maksimum liczby owocników czytelnie prezentowanych jednocześnie na ekranie telefonu.

Zastosowane tłumienie niemaksymalne jest klasowo-niezależne (ang. \textit{class-agnostic}): porównanie $\mathrm{IoU}$ w linii 56 zestawia wyłącznie geometrię ramek, bez uwzględniania indeksu klasy gatunkowej. Oznacza to, że ramka gatunku $B$ o~niższym wyniku zostanie usunięta, jeżeli jej $\mathrm{IoU}$ z~ramką gatunku $A$ przekroczy próg $0{,}45$. W~typowym scenariuszu terenowym, gdzie owocniki różnych gatunków leżą obok siebie na stole lub w~koszu i~ich projekcje na obraz nie nachodzą na siebie znacząco, sytuacja ta nie występuje. Może się jednak zdarzyć w~przypadku gęstych kęp, w~których owocniki różnych gatunków częściowo przysłaniają się nawzajem. Wariant klasowo-zależny (ang. \textit{per-class NMS}), w~którym tłumienie dotyczy wyłącznie ramek tej samej klasy, wyeliminowałby to ograniczenie, ale pozostawiałby na tym samym owocniku nakładające się ramki konkurujących gatunków. Wariant klasowo-niezależny celowo pokazuje dla jednego owocnika tylko najwyżej ocenioną hipotezę.

Detekcje przeżywające etap NMS podlegają następnie asymetrycznej polityce bezpieczeństwa mykologicznego, scentralizowanej w metodzie \texttt{MushroomRegistry.displayStatus} oraz w buforze konsensusu klatek klasy \texttt{DetectionTracker}. Filtr ten stanowi jedyne źródło prawdy dla etykiet i kolorów prezentowanych zarówno na nakładce na żywo, jak i w historii zapisów. Zanim kandydat trafi do warstwy decyzyjnej, podlega on jeszcze absolutnemu progowi odrzucenia na poziomie dekodera: wynik poniżej $0{,}18$ dla klas \texttt{POISONOUS} i \texttt{DEADLY} lub poniżej $0{,}20$ dla pozostałych klas jest odrzucany jeszcze przed alokacją struktur wynikowych, co eliminuje fałszywe alarmy tła, takie jak tekstury biurek czy klawiatury. Powyżej tego progu obowiązują trzy reguły prezentacji:
\begin{enumerate}
 \item Wczesne ostrzeżenie toksyczne ($\ge 0{,}18$ dla klas \texttt{POISONOUS} i \texttt{DEADLY}): wyświetlenie pomarańczowej etykiety „Trujący” albo czerwonej „Śmiertelnie trujący” bez wymogu konsensusu $4$ z $5$. Tor musi jednak zostać potwierdzony przez dwie kolejne zgodne klatki, zanim ramka w ogóle trafi na nakładkę.
 \item Status jadalności ($\ge 0{,}60$ dla gatunku klasy \texttt{EDIBLE}): wyświetlenie zielonej etykiety „Jadalny” wyłącznie wtedy, gdy dodatkowo spełniony jest warunek \texttt{voteConfirmed}, to znaczy dany gatunek stanowił najwyżej ocenioną predykcję w co najmniej $4$ z ostatnich $5$ klatek śledzonych w danym torze przez \texttt{DetectionTracker}.
 \item Status neutralny (pasmo pośrednie): gdy detekcja nie spełnia żadnego z powyższych warunków, to znaczy dla pewności poniżej progu $0{,}60$ (dla klas jadalnych z przedziału $0{,}20$--$0{,}60$, a dla niejadalnych powyżej $0{,}20$), a także dla detekcji gatunku jadalnego o pewności $\ge 0{,}60$, która nie została jeszcze potwierdzona przez bufor konsensusu klatek, wyświetlana jest szara etykieta „Obiekt nierozpoznany / Zbliż kamerę”. W tym paśmie zielony status jadalności nie jest prezentowany, aby ograniczyć ryzyko błędnej asercji jadalności.
\end{enumerate}
Rozdzielenie warunku pewności od warunku czasowego bufora konsensusu wynika z założenia projektowego, że jednoklatkowy skok pewności powyżej progu $0{,}60$ bywa chwilowym artefaktem szumu inferencji. Bufor ma ograniczać liczbę fałszywie potwierdzonych statusów jadalności na strumieniu kamery; nie przeprowadzono osobnego eksperymentu A/B z buforem i bez niego na sekwencjach wideo. Dwustopniowość rozwiązania polega zatem na rozdzieleniu filtracji geometryczno-statystycznej (dekodowanie tensora, szybka eliminacja, NMS) od filtracji semantycznej (asymetryczna polityka progów oraz, niezależnie, konsensus czasowy). Pierwszy etap ogranicza złożoność obliczeniową i liczbę fałszywych detekcji przestrzennych. Drugi priorytetyzuje czułość dla klas toksycznych: ostrzeżenie nie wymaga konsensusu $4$ z $5$, a asercja jadalności wymaga jednocześnie wyższej pewności. Oba statusy pojawiają się dopiero po dwóch zgodnych klatkach. Bufor $4$ z $5$ klatek jest mechanizmem stabilizacji czasowej nakładki, a nie składnikiem hipotezy $H_2$. Rozkład progów decyzyjnych na osi pewności, wraz z zaznaczeniem pasma neutralnego oraz asymetrii pomiędzy ścieżką ostrzeżenia a ścieżką asercji jadalności, przedstawiono na rys.~\ref{fig:progi_asymetryczne}.

\begin{figure}[htbp]
 \centering
 \begin{tikzpicture}[
 opis/.style={font=\footnotesize, align=center},
 strz/.style={-latex, semithick}
 ]
 % Os pewnosci
 \draw[strz] (0,0) -- (12.4,0);
 \node[opis] at (12.9,0) {$\mathrm{conf}$};
 \foreach \x/\l in {0/0{,}00, 2.16/0{,}18, 2.4/{}, 7.2/0{,}60, 12/1{,}00}{
 \draw (\x,-0.1) -- (\x,0.1);
 }
 \node[opis] at (0,-0.4) {$0{,}00$};
 \node[opis] at (2.16,-0.4) {$0{,}18$};
 \node[opis] at (2.75,-0.85) {$0{,}20$};
 \node[opis] at (7.35,-0.85) {$0{,}60$};
 \node[opis] at (12,-0.4) {$1{,}00$};

 % Sciezka toksyczna (gora)
 \fill[black!12] (0,0.35) rectangle (2.16,1.05);
 \fill[EETangerine!45] (2.16,0.35) rectangle (12,1.05);
 \draw (0,0.35) rectangle (12,1.05);
 \node[opis] at (1.08,0.7) {odrzucone};
 \node[opis] at (7.1,0.7) {\textbf{ostrzeżenie toksyczne} – bez konsensusu $4$/$5$};
 \node[opis, anchor=east] at (-0.2,0.7) {klasy\\trujące};

 % Sciezka jadalna (dol)
 \fill[black!12] (0,-2.2) rectangle (2.4,-1.5);
 \fill[black!25] (2.4,-2.2) rectangle (7.2,-1.5);
 \fill[EEBlue!35] (7.2,-2.2) rectangle (12,-1.5);
 \draw (0,-2.2) rectangle (12,-1.5);
 \node[opis] at (1.2,-1.85) {odrzucone};
 \node[opis] at (4.8,-1.85) {status neutralny (szary)};
 \node[opis] at (9.6,-1.85) {„Jadalny”\\wymaga konsensusu $4$/$5$};
 \node[opis, anchor=east] at (-0.2,-1.85) {klasy\\jadalne};
 \end{tikzpicture}
 \caption{Asymetryczna polityka progów decyzyjnych: ścieżka ostrzeżenia toksycznego (próg $0{,}18$) wobec ścieżki asercji jadalności (próg $0{,}60$ i konsensus czasowy). Źródło: opracowanie własne.}
 \label{fig:progi_asymetryczne}
\end{figure}

Scentralizowaną implementację tej polityki przedstawiono na listingu \ref{lst:safety}.

\begin{lstlisting}[language=Kotlin,caption={Asymetryczny filtr bezpieczeństwa mykologicznego (\texttt{MushroomRegistry.kt})},label={lst:safety}]
// Asymmetric thresholds: toxic warns early, edible needs a high bar + vote buffer.
const val UNSAFE_THRESHOLD = 0.18f
const val SAFE_THRESHOLD = 0.60f

// Mid band: detection exists but app asserts neither SAFE nor UNSAFE.
const val NEUTRAL_MIN_CONFIDENCE = 0.20f

// SAFE badge: species must win 4 of the last 5 frames.
const val VOTE_REQUIRED = 4
const val VOTE_WINDOW = 5

// Decode / pre-tracker gate: toxic from 0.18, everything else from 0.20.
fun thresholdFor(classId: Int): Float =
    if (speciesMap[classId]?.toxicity?.isDangerous == true) UNSAFE_THRESHOLD
    else NEUTRAL_MIN_CONFIDENCE

// Badge policy:
// - toxic/deadly @ >= 0.18 -> warning colour, no vote needed
// - edible @ >= 0.60 AND voteConfirmed -> green "Jadalny"
// - otherwise -> neutral gray ("Obiekt nierozpoznany / Zbliż kamerę")
fun displayStatus(
    toxicity: Toxicity,
    confidence: Float,
    voteConfirmed: Boolean
): DisplayStatus =
    when {
        toxicity.isDangerous && confidence >= UNSAFE_THRESHOLD ->
            DisplayStatus(toxicity.labelPl, toxicity.color)

        toxicity == Toxicity.EDIBLE &&
            confidence >= SAFE_THRESHOLD &&
            voteConfirmed ->
            DisplayStatus(toxicity.labelPl, toxicity.color)

        else ->
            DisplayStatus(NEUTRAL_LABEL, NEUTRAL_COLOR)
    }
\end{lstlisting}

Linie 2--3 definiują dwa niesymetryczne progi decyzyjne hipotezy $H_2$. Ponad trzykrotna różnica pomiędzy progiem ostrzeżenia ($0{,}18$) a progiem asercji jadalności ($0{,}60$) odzwierciedla nierówny koszt obu typów błędu: fałszywe ostrzeżenie skutkuje pozostawieniem jadalnego owocnika w lesie, natomiast fałszywe potwierdzenie jadalności może prowadzić do spożycia gatunku toksycznego. Warunek \texttt{voteConfirmed} w dalszej części listingu należy do bufora konsensusu i nie jest składnikiem $H_2$.

Funkcja \texttt{thresholdFor} (linie 13--15) stanowi bramkę stosowaną jeszcze na poziomie dekodera, przed przekazaniem detekcji do trackera. Zwraca ona $0{,}18$ dla gatunków oznaczonych predykatem \texttt{isDangerous}, obejmującym kategorie \texttt{POISONOUS} oraz \texttt{DEADLY}, natomiast $0{,}20$ dla pozostałych. Dolny próg $0{,}20$ dla klas nietoksycznych pełni funkcję filtra szumu tła i ogranicza przypadkowe aktywacje na teksturach niebędących owocnikami.

Właściwa polityka prezentacji zawarta jest w wyrażeniu \texttt{when} (linie 26--37), którego kolejność gałęzi ma znaczenie semantyczne. Gałąź pierwsza (linie 27--28) obsługuje gatunki niebezpieczne i jest ewaluowana jako pierwsza, wobec czego ostrzeżenie ma pierwszeństwo przed jakimkolwiek innym statusem. Nie występuje w niej odwołanie do \texttt{voteConfirmed}: ostrzeżenie nie czeka na konsensus $4$ z $5$. Sam tor pojawia się jednak dopiero po dwóch kolejnych zgodnych klatkach (\texttt{CONFIRM\_STREAK}~$=2$, podrozdział o trackerze), czyli najwcześniej po około $250$\,ms przy kadencji $4$\,FPS.

Gałąź druga (linie 30--33) formułuje warunek koniunkcyjny dla statusu jadalności: wymagana jest jednocześnie przynależność do kategorii \texttt{EDIBLE}, pewność nie mniejsza niż $0{,}60$ oraz pozytywny wynik bufora konsensusu. Niespełnienie któregokolwiek z trzech warunków przenosi detekcję do gałęzi domyślnej (linie 35--36), zwracającej neutralną, szarą etykietę. Parametr \texttt{voteConfirmed} nie ma wartości domyślnej: nakładka, ekran historii i Tryb Badawczy przekazują go jawnie, więc zielony status nie może powstać przez pominięcie argumentu.

Zasadę działania okna przesuwnego, wraz z przykładowym przebiegiem prowadzącym do potwierdzenia i do jego cofnięcia, zilustrowano na rys.~\ref{fig:bufor_konsensusu}.

\begin{figure}[htbp]
 \centering
 \begin{tikzpicture}[
 klatka/.style={draw, minimum width=1.15cm, minimum height=0.72cm,
 font=\footnotesize, inner sep=1pt},
 zgodna/.style={klatka, fill=EEBlue!28},
 odstajaca/.style={klatka, fill=EETangerine!45},
 opis/.style={font=\footnotesize, align=left}
 ]
 % Naglowek osi klatek
 \foreach \i/\x in {1/0, 2/1.25, 3/2.5, 4/3.75, 5/5.0}{
 \node[font=\scriptsize] at (\x,0.75) {$k_{\i}$};
 }

 % Przebieg A: 4 z 5 -> potwierdzone
 \node[zgodna] at (0,0) {A};
 \node[zgodna] at (1.25,0) {A};
 \node[odstajaca] at (2.5,0) {B};
 \node[zgodna] at (3.75,0) {A};
 \node[zgodna] at (5.0,0) {A};
 \draw[EEUltramarine, thick, rounded corners=3pt] (-0.65,-0.5) rectangle (5.65,0.5);
 \node[opis, anchor=west, EEUltramarine] at (6.1,0) {$4$ z $5$ dla klasy A\\$\Rightarrow$ \texttt{voteConfirmed} $=$ \texttt{true}};

 % Przebieg B: 3 z 5 -> niepotwierdzone
 \node[zgodna] at (0,-1.55) {A};
 \node[odstajaca] at (1.25,-1.55) {B};
 \node[zgodna] at (2.5,-1.55) {A};
 \node[odstajaca] at (3.75,-1.55) {C};
 \node[zgodna] at (5.0,-1.55) {A};
 \draw[EETangerine, thick, rounded corners=3pt] (-0.65,-2.05) rectangle (5.65,-1.05);
 \node[opis, anchor=west, EETangerine] at (6.1,-1.55) {$3$ z $5$ dla klasy A\\$\Rightarrow$ status neutralny};
 \end{tikzpicture}
 \caption{Bufor konsensusu klatek: okno pięcioelementowe dopuszcza jedną klatkę odstającą, lecz blokuje status jadalności przy dwóch lub więcej niezgodnościach. Źródło: opracowanie własne.}
 \label{fig:bufor_konsensusu}
\end{figure}

Klasę \texttt{DetectionTracker} realizującą ten bufor prezentuje listing \ref{lst:voting}.

\begin{lstlisting}[language=Kotlin,caption={Bufor konsensusu klatek – potwierdzenie gatunku w $4$ z ostatnich $5$ klatek (\texttt{DetectionTracker.kt})},label={lst:voting}]
private class Track(val id: Int, seed: DetectionResult, nowMs: Long) {
    var species: MushroomSpecies = seed.species
    var confidence: Float = seed.confidence
    var box: NormalizedBox = seed.boundingBox
    val firstSeenMs: Long = nowMs
    var lastSeenMs: Long = nowMs
    var frozen: Boolean = false

    // Number of consecutive frames this track kept the same species class.
    var confirmCount: Int = 1

    // Rolling top-class ids for SAFE vote: need VOTE_REQUIRED hits in VOTE_WINDOW.
    private val recentClassIds = ArrayDeque<Int>(MushroomRegistry.VOTE_WINDOW)

    init {
        recentClassIds.addLast(seed.species.classId)
    }

    fun absorb(detection: DetectionResult, nowMs: Long) {
        lastSeenMs = nowMs
        confidence = detection.confidence

        confirmCount = if (detection.species.classId == species.classId) {
            confirmCount + 1
        } else {
            1
        }
        species = detection.species

        recentClassIds.addLast(detection.species.classId)
        if (recentClassIds.size > MushroomRegistry.VOTE_WINDOW) {
            recentClassIds.removeFirst()
        }

        if (nowMs - firstSeenMs >= FREEZE_MS) {
            frozen = true
            return
        }
        box = smooth(box, detection.boundingBox)
    }

    // True when current species won >= VOTE_REQUIRED of the recent window.
    fun voteConfirmed(): Boolean =
        recentClassIds.count { it == species.classId } >= MushroomRegistry.VOTE_REQUIRED

    fun toTracked(): TrackedDetection =
        TrackedDetection(id, species, confidence, box, frozen, voteConfirmed())
}
\end{lstlisting}

Strukturą nośną bufora jest kolejka dwustronna \texttt{ArrayDeque} (linia 13), zainicjalizowana pojemnością równą rozmiarowi okna. W każdej klatce wykonywane jest jedno dopisanie na końcu oraz co najwyżej jedno usunięcie z początku, a obie operacje realizowane są w czasie stałym.

Metoda \texttt{absorb} (linie 19--40) wywoływana jest po dopasowaniu toru do detekcji z bieżącej klatki przy $\mathrm{IoU}>0{,}35$. Linie 30--33 realizują okno przesuwne: nowa etykieta klasowa dopisywana jest na końcu kolejki, a po przekroczeniu rozmiaru okna najstarszy element jest usuwany. Kolejka utrzymuje zatem co najwyżej pięć ostatnich rozstrzygnięć klasowych danego toru. Linie 35--38 zamrażają ramkę po $5$\,s nieprzerwanego śledzenia (\texttt{FREEZE\_MS}); listing pomija przypisanie diagnostycznych wyników top-$5$, używanych wyłącznie w Trybie Badawczym.

Metoda \texttt{voteConfirmed} (linie 43--44) zlicza w oknie wystąpienia identyfikatora klasy aktualnie przypisanej do toru i porównuje wynik z progiem $4$. Licznik \texttt{confirmCount} (linie 23--27) pełni inną rolę: zlicza wyłącznie kolejne klatki o tej samej klasie i jest zerowany przy zmianie etykiety. Tor trafia na nakładkę dopiero przy \texttt{confirmCount}~$\ge 2$, więc zmiana klasy na jedną klatkę chowa ramkę. Kryterium głosowania jest łagodniejsze i dopuszcza pojedynczą klatkę odstającą wewnątrz okna: chwilowe migotanie predykcji pomiędzy dwoma zbliżonymi gatunkami nie unieważnia stabilnego rozpoznania, natomiast dwie lub więcej klatek odstających w oknie pięcioelementowym blokują zielony status.

Metoda \texttt{toTracked} (linie 46--47) przekazuje wynik głosowania do warstwy prezentacji jako pole \texttt{voteConfirmed} struktury \texttt{TrackedDetection}, skąd odczytuje je \texttt{displayStatus}. Tracker odpowiada za stabilność czasową predykcji, a rejestr za interpretację bezpieczeństwa. Tor niewidziany przez $1{,}2$\,s (\texttt{STALE\_MS}) jest usuwany. W trybie heartbeat ($1{,}5$\,s) pojedyncze pominięcie na klatce kontrolnej może więc zakończyć tor; jest to ograniczenie obecnej implementacji. Cykl życia toru przedstawiono na rys.~\ref{fig:cykl_toru}; rysunek pokazuje konsensus $4$ z $5$, a nie dwuklatkowy warunek publikacji ani zamrożenie po $5$\,s.

\begin{figure}[htbp]
 \centering
 \begin{tikzpicture}[
 stan/.style={draw, rounded corners=4pt, align=center, font=\scriptsize,
 text width=2.5cm, minimum height=9mm, inner sep=2.5pt},
 nowy/.style={stan, fill=black!7},
 akt/.style={stan, fill=EEBlue!14},
 pot/.style={stan, fill=EETangerine!25},
 kon/.style={stan, fill=black!12},
 p/.style={-latex, semithick, font=\scriptsize}
 ]
 \node[nowy] (n) at (0,0) {\textsc{Nowy}\\okno: $1$ klatka};
 \node[akt] (a) at (4.6,0) {\textsc{Śledzony}\\dopasowanie IoU};
 \node[pot] (p) at (9.2,0) {\textsc{Potwierdzony}\\$4$/$5$ zgodnych};
 \node[kon] (k) at (4.6,-3.4) {\textsc{Wygasły}\\brak dopasowania};

 \draw[p] (-1.7,0) -- node[above] {detekcja} (n);
 \draw[p] (n) -- node[above, align=center] {kolejna\\klatka} (a);
 \draw[p] (a) to[bend left=22] node[above, align=center] {konsensus\\osiągnięty} (p);
 \draw[p] (p) to[bend left=22] node[below, align=center] {zmiana klasy\\w oknie} (a);
 \draw[p] (a) -- node[right, align=left, xshift=2pt] {przekroczony\\limit czasu} (k);
 \draw[p] (p.south) |- node[pos=0.75, above] {utrata obiektu} (k.east);
 \end{tikzpicture}
 \caption{Cykl życia pojedynczego toru w module \texttt{DetectionTracker}: przejście do stanu potwierdzonego wymaga konsensusu $4$ z $5$ klatek i jest odwracalne. Źródło: opracowanie własne.}
 \label{fig:cykl_toru}
\end{figure}

\subsection{Lokalna baza danych Room i zapis historii znalezisk}
Historia znalezisk przechowuje się między uruchomieniami aplikacji. Persystencję realizuje Room nad SQLite \cite{androidroom}. Ekran historii czyta strumień przez \texttt{HistoryViewModel}, a zapis idzie przez \texttt{FindsRepository} do DAO. Baza ma schemat w wersji $2$; przy zmianie schematu ustawiono \texttt{fallbackToDestructiveMigration()}, więc aktualizacja schematu kasuje dotychczasowe znaleziska. W prototypie jest to świadome uproszczenie. Schemat biblioteki ilustruje rys.~\ref{fig:room_db_web}.

\begin{figure}[htbp]
 \centering
 \includegraphics[width=0.85\textwidth]{room_db_arch_official}
 \caption{Struktura lokalnej bazy danych Room. Źródło: \cite{androidroom}.}
 \label{fig:room_db_web}
\end{figure}

Encja \texttt{MushroomFind} ma wyłącznie pola \texttt{val} i mapuje tabelę \texttt{mushroom\_finds}: klucz \texttt{id}, nazwa łacińska \texttt{speciesName}, etykieta toksyczności \texttt{toxicityLabel}, \texttt{confidence}, flaga \texttt{voteConfirmed} (domyślnie \texttt{false} w kodzie; kolumna ma \texttt{defaultValue="1"} dla wierszy sprzed migracji), \texttt{timestamp} oraz \texttt{imagePath} do JPEG w \texttt{filesDir}. DAO to \texttt{MushroomFindDao}; \texttt{FindsRepository} woła je z \texttt{Dispatchers.IO}, żeby dysk nie blokował Compose. Lista historii to \texttt{Flow}, więc ekran odświeża się sam po każdym zapisie.

Interfejs DAO i warstwę repozytorium pokazano na listingu \ref{lst:room}.

\begin{lstlisting}[language=Kotlin,caption={Asynchroniczna warstwa persystencji Room z reaktywnym strumieniem \texttt{Flow} (\texttt{MushroomFindDao.kt}, \texttt{FindsRepository.kt})},label={lst:room}]
@Dao
interface MushroomFindDao {

    /** Insert a new find and return its generated row id. */
    @Insert
    suspend fun insert(find: MushroomFind): Long

    /** Observe all finds, newest first. */
    @Query("SELECT * FROM mushroom_finds ORDER BY timestamp DESC")
    fun observeAll(): Flow<List<MushroomFind>>

    /** Delete a find by id (does not remove the backing image file). */
    @Query("DELETE FROM mushroom_finds WHERE id = :id")
    suspend fun deleteById(id: Long)
}

class FindsRepository(context: Context) {

    private val appContext = context.applicationContext
    private val dao = MushroomDatabase.get(appContext).findDao()

    /** All saved finds, newest first, as a reactive stream. */
    val finds: Flow<List<MushroomFind>> = dao.observeAll()

    suspend fun save(
        speciesName: String,
        toxicityLabel: String,
        confidence: Float,
        voteConfirmed: Boolean,
        crop: Bitmap
    ): MushroomFind = withContext(Dispatchers.IO) {
        val dir = File(appContext.filesDir, FINDS_DIR).apply { mkdirs() }
        val file = File(dir, "find_${UUID.randomUUID()}.jpg")

        FileOutputStream(file).use { out ->
            crop.compress(Bitmap.CompressFormat.JPEG, JPEG_QUALITY, out)
        }

        val find = MushroomFind(
            speciesName = speciesName,
            toxicityLabel = toxicityLabel,
            confidence = confidence,
            voteConfirmed = voteConfirmed,
            timestamp = System.currentTimeMillis(),
            imagePath = file.absolutePath
        )
        val id = dao.insert(find)
        find.copy(id = id)
    }
}
\end{lstlisting}

Adnotacje \texttt{@Dao}, \texttt{@Insert} oraz \texttt{@Query} (linie 1--15) przetwarzane są przez procesor adnotacji Room na etapie kompilacji, który generuje implementację interfejsu. Zapytania SQL są weryfikowane względem schematu bazy w czasie budowy projektu, więc literówka w nazwie kolumny kończy się błędem kompilacji, a nie wyjątkiem w czasie działania.

Metody \texttt{insert} (linia 6) oraz \texttt{deleteById} (linia 14) zadeklarowano jako \texttt{suspend}, więc wywołanie spoza korutyny odrzuca kompilator. Metoda \texttt{observeAll} (linia 10) zwraca strumień \texttt{Flow}: Room emituje nową listę przy każdej zmianie tabeli, a ekran historii odświeża się bez osobnego odczytu.

Metoda \texttt{save} (linie 25--49) łączy zapis pliku JPEG ze wstawieniem rekordu. Całość opakowano w \texttt{withContext(Dispatchers.IO)} (linia 31). Katalog powstaje wewnątrz \texttt{context.filesDir} (linia 32), w prywatnym katalogu aplikacji, i znika wraz z odinstalowaniem. Nazwę pliku generuje \texttt{UUID} (linia 33). Listing pomija wpis do dziennika, obecny w kodzie po wstawieniu rekordu.

Konstrukcja \texttt{FileOutputStream(file).use} (linie 35--37) zamyka strumień także przy wyjątku kompresji. Obraz zapisywany jest przed \texttt{dao.insert} (linia 47), więc w bazie nie powstaje wiersz wskazujący na nieistniejący plik. \texttt{find.copy(id = id)} (linia 48) zwraca nową encję z wygenerowanym kluczem, ponieważ pola \texttt{MushroomFind} są niemodyfikowalne.

Ręczny zapis uruchamia \texttt{FrameCropper}: aplikacja bierze ostatnią zorientowaną pionowo klatkę z toru \texttt{ImageAnalysis} i wycina region znormalizowanej ramki. Wycinek trafia do JPEG w prywatnym katalogu piaskownicy (\texttt{context.filesDir}) – poza publiczną galerią urządzenia. Dane użytkownika pozostają izolowane i znikają automatycznie przy odinstalowaniu.

Ekran historii (\texttt{HistoryScreen}) to deklaratywna lista \texttt{LazyColumn}. Biblioteka Coil (\texttt{coil-compose}, \texttt{AsyncImage}) ładuje miniatury leniwie. Pamięć nie rośnie liniowo wraz z liczbą zapisów \cite{coil}. Status wizualny wyznacza \texttt{MushroomRegistry.displayStatus} na podstawie zapisanej pewności i flagi \texttt{voteConfirmed}. Zielone „Jadalny” pojawia się tylko dla rekordów, które w chwili zapisu miały pewność $\ge 0{,}60$ i domknięty konsensus $4$ z $5$. W pozostałych przypadkach wyświetlana jest szara etykieta neutralna, także dla gatunków niejadalnych. Stały czerwony baner na górze ekranu przypomina: „Ostrzeżenie: Aplikacja ma charakter pomocniczy. Nigdy nie spożywaj grzybów bez pewności eksperta”.

\chapter{Badanie i analiza skuteczności systemu}\label{chap:badania}

\section{Metodyka testowania i miary oceny}
Do końcowej ewaluacji użyto niezależnego zbioru testowego (\texttt{images/test}, \texttt{labels/test}), złożonego z obserwacji iNaturalist, których \texttt{obs\_id} nie występuje w \texttt{train} ani \texttt{val}. Zapobiega to wyciekowi informacji ze zbioru uczącego do etapu oceny: dwa zdjęcia tego samego owocnika nie mogą jednocześnie uczestniczyć w uczeniu i w pomiarze. Zgodnie z tab.~\ref{tab:podzial_zbioru} zbiór liczy $32\,727$ obrazów oraz $73\,272$ anotowanych instancji owocników ($147$ klas; przecięcie \texttt{obs\_id} z korpusem uczącym: $0$). Wszystkie raportowane w tym rozdziale wartości $mAP$, $P$ i $R$ uzyskano wyłącznie na tym izolowanym podzbiorze.

Jakościowe dopasowanie ramki predykowanej $B_p$ do ramki referencyjnej $B_{gt}$ mierzono współczynnikiem przecięcia nad sumą (IoU), zgodnie z definicją wprowadzoną wcześniej:
\begin{equation}
\mathrm{IoU}(B_p, B_{gt}) = \frac{\lvert B_p \cap B_{gt} \rvert}{\lvert B_p \cup B_{gt} \rvert}.
\end{equation}
Interpretację geometryczną licznika i mianownika powyższego ilorazu, wraz z wpływem wartości progu na kwalifikację detekcji, przedstawiono na rys.~\ref{fig:iou_definicja}.

\begin{figure}[htbp]
 \centering
 \begin{tikzpicture}[scale=0.92]
 % – (a) przeciecie
 \begin{scope}[shift={(0,0)}]
 \fill[EEBlue!30] (0,0) rectangle (2.4,1.8);
 \fill[EETangerine!35] (1.1,0.6) rectangle (3.5,2.4);
 \fill[EEUltramarine!55] (1.1,0.6) rectangle (2.4,1.8);
 \draw[EEUltramarine, very thick] (0,0) rectangle (2.4,1.8);
 \draw[EETangerine, very thick] (1.1,0.6) rectangle (3.5,2.4);
 \node[font=\scriptsize, anchor=north] at (1.75,-0.25) {(a) $\lvert B_p \cap B_{gt}\rvert$};
 \end{scope}
 %, (b) suma
 \begin{scope}[shift={(5.0,0)}]
 \fill[black!20] (0,0) rectangle (2.4,1.8);
 \fill[black!20] (1.1,0.6) rectangle (3.5,2.4);
 \draw[EEUltramarine, very thick] (0,0) rectangle (2.4,1.8);
 \draw[EETangerine, very thick] (1.1,0.6) rectangle (3.5,2.4);
 \node[font=\scriptsize, anchor=north] at (1.75,-0.25) {(b) $\lvert B_p \cup B_{gt}\rvert$};
 \end{scope}
 % legenda
 \draw[EEUltramarine, very thick] (9.4,1.9) -- (10.0,1.9);
 \node[font=\scriptsize, anchor=west] at (10.1,1.9) {$B_{gt}$ – ramka referencyjna};
 \draw[EETangerine, very thick] (9.4,1.35) -- (10.0,1.35);
 \node[font=\scriptsize, anchor=west] at (10.1,1.35) {$B_p$. Ramka predykowana};
 \node[font=\scriptsize, anchor=west, align=left] at (9.4,0.45)
 {$\mathrm{IoU}\ge0{,}50 \Rightarrow TP$\\[1pt] $\mathrm{IoU}<0{,}50 \Rightarrow FP$ oraz $FN$};
 \end{tikzpicture}
 \caption{Geometryczna interpretacja współczynnika IoU oraz reguła kwalifikacji detekcji przy progu $0{,}50$. Źródło: opracowanie własne.}
 \label{fig:iou_definicja}
\end{figure}

Na podstawie progu IoU oraz zgodności etykiet klasowych klasyfikowano detekcje jako prawdziwie dodatnie ($TP$), fałszywie dodatnie ($FP$) lub fałszywie ujemne ($FN$). Precyzję $P$ oraz czułość (ang. \textit{recall}) $R$ zdefiniowano odpowiednio jako:
\begin{equation}
P = \frac{TP}{TP + FP}, \qquad R = \frac{TP}{TP + FN}.
\end{equation}
Średnią precyzję dla pojedynczej klasy ($AP$) obliczano jako pole pod krzywą precyzja--czułość:
\begin{equation}
AP = \int_{0}^{1} p(r)\,dr,
\end{equation}
natomiast średnią precyzję dla wielu klas ($mAP$) jako średnią arytmetyczną wartości $AP$ po $K$ gatunkach:
\begin{equation}
mAP = \frac{1}{K}\sum_{i=1}^{K} AP_i.
\end{equation}

W praktyce raportowano dwa warianty tej miary, zgodne ze standardem protokołu COCO \cite{lin2014coco}:
\begin{enumerate}
 \item $mAP_{50}$ (mAP@$0{,}50$) -- ewaluacja przy łagodnym progu przestrzennym $\mathrm{IoU}=0{,}50$, akcentująca poprawność klasyfikacji przy umiarkowanych wymaganiach lokalizacyjnych.
 \item $mAP_{50\text{--}95}$ (mAP@$0{,}50{:}0{,}95$) -- średnia z dziesięciu progów IoU od $0{,}50$ do $0{,}95$ z krokiem $0{,}05$, mierząca ścisłość lokalizacji ramek.
\end{enumerate}
Pierwsza miara jest bardziej odporna na drobne przesunięcia geometrii, druga natomiast surowiej karze niedokładne dopasowanie przestrzenne.

W kontekście identyfikacji grzybów profil ryzyka jest silnie asymetryczny, przy czym należy rozróżnić dwa poziomy błędu. Na poziomie detekcji fałszywie dodatni wynik ($FP$) to zbędna lub źle zlokalizowana ramka (np.\ detekcja na kamieniu albo pomyłka między dwoma gatunkami jadalnymi) -- w~izolacji jest to głównie niedogodność użytkowa. Fałszywie ujemny wynik ($FN$) to przeoczenie owocnika obecnego w~kadrze. Z~perspektywy bezpieczeństwa konsumenta najpoważniejszy jest natomiast błąd innego rodzaju: przypisanie gatunkowi \emph{rzeczywiście toksycznemu} asercji „Jadalny”. Nie jest to detekcyjny $FN$, lecz krytyczny wynik „fałszywie bezpieczny”, mierzony osobno wskaźnikiem CFPR i~macierzą ryzyka (rozdz.~\ref{sec:safety_matrix}). To właśnie minimalizacja tego błędu -- nawet kosztem precyzji i~liczby fałszywych alarmów na gatunkach nietoksycznych -- jest nadrzędnym kryterium projektowym, realizowanym przez asymetryczną parę progów opisaną w~rozdziale \ref{chap:implementacja}. W~konsekwencji dla klas toksycznych priorytetem pozostaje wysoka czułość ($R$).

\section{Analiza wyników detekcji na zbiorze testowym}
Model wytrenowany na architekturze YOLO11m poddano ewaluacji na izolowanym zbiorze testowym (tab.~\ref{tab:podzial_zbioru}). Obserwacje tego podzbioru nie brały udziału w uczeniu ani w strojeniu hiperparametrów: nazwy plików zachowują konwencję \texttt{\{obs\_id\}\_\{photo\_id\}.jpg}, a zbiór odrzucono i odbudowano tak, by żaden \texttt{obs\_id} z \texttt{train}/\texttt{val} nie wrócił do testu. Pomiarów dokonano standardowym potokiem walidacyjnym biblioteki Ultralytics, zapewniającym porównywalność wyników z protokołem raportowania mAP w detekcji obiektów \cite{ultralytics2024yolo11, lin2014coco}.

Globalne metryki na całym zbiorze testowym ($147$ klas) wyniosły: $mAP_{50}=0{,}8275$ ($82{,}75\%$), $mAP_{50\text{--}95}=0{,}7823$ ($78{,}23\%$), precyzja $P=0{,}8090$ ($80{,}90\%$) oraz czułość $R=0{,}7532$ ($75{,}32\%$). Są to wyniki wag PyTorch (\texttt{best.pt}) w protokole Ultralytics \texttt{val} (\texttt{split=test}). Artefakt produkcyjny na telefonie to TFLite FP16; próbę porównania $mAP$ obu formatów opisano w podrozdz.~\ref{sec:tflite_parity}. Wartość $mAP_{50}$ potwierdza, że model lokalizuje i rozpoznaje owocniki na niewidzianych obserwacjach, w zróżnicowanych warunkach oświetlenia naturalnego oraz na tle leśnym. Poza wynikiem zagregowanym po wszystkich klasach obliczono również metryki odrębnie dla dwóch podzbiorów gatunkowych zdefiniowanych kategorią toksykologiczną: gatunków jadalnych oraz gatunków trujących i śmiertelnie trujących (średnia arytmetyczna $AP$ po klasach danego podzbioru, zgodnie z definicją $mAP$). Charakterystykę zbioru testowego wraz z globalnymi metrykami zagregowanymi po wszystkich klasach przedstawiono w tab.~\ref{tab:wyniki_detekcji}, natomiast rozbicie wyników na kategorie toksykologiczne – w tab.~\ref{tab:wyniki_kategorie}. Asymetrię ryzyka konsumpcyjnego oceniono odrębnie w macierzy bezpieczeństwa (rozdz.~\ref{sec:safety_matrix}).

\begin{table}[h!]
 \centering
 \caption{Charakterystyka zbioru testowego oraz globalne metryki detekcji YOLO11m (wagi PyTorch, \texttt{split=test}). Źródło: opracowanie własne.}
 \label{tab:wyniki_detekcji}
 \begin{tabular}{|l|c|}
 \hline
 \textbf{Wielkość} & \textbf{Wartość} \\
 \hline \hline
 Liczba obrazów testowych & $32\,727$ \quad (tab.~\ref{tab:podzial_zbioru}) \\
 \hline
 Liczba anotowanych instancji & $73\,272$ \\
 \hline
 Izolacja po \texttt{obs\_id} & $0$ wspólnych z \texttt{train}/\texttt{val} \\
 \hline
 Liczba klas (odrębnych gatunków) & $147$ ($145$) \\
 \hline \hline
 $mAP_{50}$ & $0{,}8275$ \quad ($82{,}75\%$) \\
 \hline
 $mAP_{50\text{--}95}$ & $0{,}7823$ \quad ($78{,}23\%$) \\
 \hline
 Precyzja $P$ & $0{,}8090$ \quad ($80{,}90\%$) \\
 \hline
 Czułość $R$ & $0{,}7532$ \quad ($75{,}32\%$) \\
 \hline
 \end{tabular}
\end{table}

\begin{table}[htbp]
 \centering
 \caption{Wyniki detekcji na izolowanym zbiorze testowym w podziale na kategorie toksykologiczne (średnia po klasach danej kategorii, rejestr z tab.~\ref{tab:toksycznosc}). Źródło: opracowanie własne na podstawie \texttt{category\_test\_metrics.json}.}
 \label{tab:wyniki_kategorie}
 \begin{tabular}{|l|c|c|c|c|}
 \hline
 \textbf{Zakres ewaluacji} & $\boldsymbol{mAP_{50}}$ & $\boldsymbol{mAP_{50\text{--}95}}$ & $\boldsymbol{P}$ & $\boldsymbol{R}$ \\
 \hline \hline
 Cały zbiór testowy ($147$ klas) & $0{,}8275$ & $0{,}7823$ & $0{,}8090$ & $0{,}7532$ \\
 \hline
 Gatunki jadalne ($69$ klas) & $0{,}8134$ & $0{,}7723$ & $0{,}7991$ & $0{,}7487$ \\
 \hline
 Gatunki niejadalne ($44$ klasy) & $0{,}8506$ & $0{,}8004$ & $0{,}8235$ & $0{,}7672$ \\
 \hline
 Gatunki trujące i śmiertelnie trujące ($34$ klasy) & $0{,}8262$ & $0{,}7791$ & $0{,}8102$ & $0{,}7444$ \\
 \hline
 \end{tabular}
\end{table}

Uzyskany poziom $mAP_{50}=82{,}75\%$ przy precyzji $80{,}90\%$ i czułości $75{,}32\%$ stanowi bazę operacyjną detektora na obserwacjach, których model nie widział w uczeniu. Rozbicie na podzbiory toksykologiczne (\texttt{test\_category\_map.py}) pokazuje wyrównane wyniki między kategoriami: klasy toksyczne osiągają $mAP_{50}=82{,}62\%$ przy czułości $74{,}44\%$, nieco wyżej niż klasy jadalne ($81{,}34\%$, $R=74{,}87\%$), a najwyższy wynik mają klasy niejadalne ($85{,}06\%$). Model nie jest zatem systematycznie słabszy na gatunkach trujących, co jest pożądane ze względów bezpieczeństwa. Średnia po wszystkich $147$ klasach jest średnią ważoną liczbą klas trzech wierszy kategorii. Redukcja ryzyka konsumpcyjnego nie wynika z samej głowicy klasyfikacyjnej, lecz z pary progów $0{,}18$/$0{,}60$ oraz z macierzy ryzyka.

Zaobserwowano także systematyczną lukę pomiędzy $mAP_{50}$ a $mAP_{50\text{--}95}$, widoczną we wszystkich trzech zakresach ewaluacji. Niższy wynik przy surowych progach IoU ($0{,}50$--$0{,}95$) wynika z naturalnej zmienności morfologicznej kapeluszy oraz częściowego przesłonięcia owocników ściółką liściową, które utrudniają idealne dopasowanie geometryczne ramek. Różnica ta nie podważa użyteczności modelu w scenariuszu mobilnym, gdzie nadrzędne jest prawidłowe wykrycie i klasyfikacja gatunku, a nie milimetrowa precyzja obrysu.

\subsection{Porównanie PyTorch i TFLite FP16}\label{sec:tflite_parity}
Hipoteza $H_1$ mówi o modelu FP16 na telefonie, podczas gdy tablice $mAP$ powyżej pochodzą z wag \texttt{best.pt}. Na izolowanym teście wylosowano $2\,500$ obrazów (\texttt{seed}=42; $5\,568$ instancji). Dla \texttt{.pt} uruchomiono Ultralytics \texttt{val} na RTX~4070. Dla \texttt{best\_saved\_model\_fp16.tflite} nie użyto LiteRT \texttt{val}: ten dekoder nie parował tensora $[1,\;151,\;8400]$ z protokołem COCO. Zamiast tego Interpreter TFLite czyta wyjście w układzie aplikacji (\texttt{MushroomDetector.kt}: indeks ${}= \text{kanał}\times 8400+\text{punkt}$), letterbox $640$ z tłem $114$ (w aplikacji tło ma wartość $136$), NMS zachłanny niezależny od klasy przy jednej etykiecie na kotwicę ($\mathrm{conf}=0{,}001$, $\mathrm{IoU}=0{,}70$, $\mathrm{max\_det}=300$; \texttt{val} tłumi ramki osobno w każdej klasie), a $mAP$ liczy \texttt{ap\_per\_class} Ultralytics (\texttt{test\_tflite\_parity.py}).

Na tej samej próbie PyTorch dał $mAP_{50}=0{,}8467$, $mAP_{50\text{--}95}=0{,}8042$, $P=0{,}8139$, $R=0{,}7671$. TFLite FP16: $mAP_{50}=0{,}7914$, $mAP_{50\text{--}95}=0{,}7559$, $P=0{,}8169$, $R=0{,}7454$, po NMS $16\,917$ predykcji, średnio $177{,}2$\,ms/kadr na CPU (XNNPACK). Różnica $mAP_{50}$ wynosi $-0{,}0553$ (spadek względny $6{,}5\%$); $mAP_{50\text{--}95}$ spada o $0{,}0483$, czułość o $0{,}0217$, precyzja jest nieznacznie wyższa ($+0{,}0030$). Konwersja nie zeruje detekcji, ale \textbf{nie} jest ilościowym parytetem z \texttt{.pt}. Metryki pełnego testu w tab.~\ref{tab:wyniki_detekcji} pozostają pomiarem PyTorch; tab.~\ref{tab:tflite_parity} dotyczy wyłącznie tej próby. Czas $177{,}2$\,ms to Interpreter na CPU stacji, nie \texttt{detect()} na S23 Ultra.

\begin{table}[htbp]
 \centering
 \caption{Ta sama próba $2\,500$ obrazów izolowanego testu (\texttt{seed}=42). PyTorch: Ultralytics \texttt{val} (RTX~4070). TFLite: Interpreter CPU + dekoder Androida. Źródło: \texttt{tflite\_parity.json}.}
 \label{tab:tflite_parity}
 \begin{tabular}{|l|c|c|}
 \hline
 \textbf{Wielkość} & \textbf{PyTorch \texttt{best.pt}} & \textbf{TFLite FP16 (dekoder Android)} \\
 \hline \hline
 Obrazy / instancje & \multicolumn{2}{c|}{$2\,500$ / $5\,568$} \\
 \hline
 $mAP_{50}$ & $0{,}8467$ & $0{,}7914$ \\
 \hline
 $mAP_{50\text{--}95}$ & $0{,}8042$ & $0{,}7559$ \\
 \hline
 Precyzja $P$ & $0{,}8139$ & $0{,}8169$ \\
 \hline
 Czułość $R$ & $0{,}7671$ & $0{,}7454$ \\
 \hline
 Inferencja / obraz & GPU (RTX~4070) & $177{,}2$\,ms (CPU XNNPACK) \\
 \hline
 \end{tabular}
\end{table}

Należy podkreślić metodyczną różnicę w postprocessingu: eksperyment parzystości w tab.~\ref{tab:tflite_parity} posłużył do oceny jakości surowego tensora według standardowego protokołu mAP Ultralytics ($\mathrm{conf}=0{,}001$, $\mathrm{IoU}=0{,}70$, $\mathrm{max\_det}=300$). W docelowej aplikacji mobilnej (\texttt{MushroomDetector.kt}) stosowany jest odmienny potok operacyjny: wstępna selekcja top-$50$ kandydatów, próg class-agnostic NMS $\mathrm{IoU}=0{,}45$ oraz bramki pewności $0{,}18$ i $0{,}20$. Pierwszy pomiar CFPR wykonano na predykcjach PyTorch (\texttt{test\_h2\_baseline.py}). Dlatego powtórzono go na artefakcie TFLite FP16 z postprocessingiem odtwarzającym potok aplikacji (podrozdz.~\ref{sec:safety_matrix}, tab.~\ref{tab:cfpr_porownanie}).

\section{Analiza macierzy ryzyka i asymetryczna polityka bezpieczeństwa}\label{sec:safety_matrix}
Szczególną kategorię błędów detekcji stanowią pomyłki pomiędzy gatunkami bliźniaczymi (ang. \textit{look-alikes}). Na zbiorze testowym oprócz $mAP$ liczono macierz ryzyka. Rejestr aplikacji (\texttt{MushroomRegistry}, $147$ klas) ma trzy etykiety biologiczne: jadalne (\texttt{EDIBLE}), niejadalne (\texttt{INEDIBLE}) oraz toksyczne (\texttt{POISONOUS} albo \texttt{DEADLY}). Warstwa wyświetlania (\texttt{displayStatus}) jest inna niż ta etykieta:
\begin{itemize}
 \item \textbf{asercja jadalności} (zielone „Jadalny”) -- tylko klasa \texttt{EDIBLE} i pewność $\ge 0{,}60$;
 \item \textbf{ostrzeżenie toksyczne} -- klasa toksyczna i pewność $\ge 0{,}18$;
 \item \textbf{brak asercji} (szary NEUTRAL) -- pozostałe przypadki, w tym \emph{wszystkie} predykcje niejadalne, niezależnie od pewności.
\end{itemize}
\texttt{CFPR} liczy się wyłącznie na instancjach toksycznych: $P(\text{asercja jadalności}\mid\text{gt toksyczny})$. Niejadalne nie są jadalne. Hipotezę $H_2$ rozstrzyga porównanie tej samej definicji przy parze $0{,}18$/$0{,}60$ oraz przy jednym progu $\tau=0{,}30$, na \emph{tym samym} izolowanym teście, w jednym przebiegu \texttt{predict} (\texttt{test\_h2\_baseline.py}, $32\,727$ obrazów, $73\,272$ instancji). Ten pomiar dotyczy wag PyTorch (\texttt{best.pt}) i postprocessingu Ultralytics (NMS w obrębie klasy, $\mathrm{IoU}=0{,}70$, $\mathrm{conf}\ge0{,}05$). Aplikacja uruchamia artefakt TFLite FP16 z własnym dekoderem, dlatego pomiar powtórzono skryptem \texttt{test\_tflite\_cfpr.py}, który odtwarza potok \texttt{MushroomDetector.kt}: letterbox z szarym tłem, progi wstępne $0{,}18$ (klasy toksyczne) i $0{,}20$ (pozostałe), $50$ najwyżej ocenionych kandydatów i NMS niezależny od klasy ($\mathrm{IoU}>0{,}45$, najwyżej $20$ detekcji). Bufora $4$ z $5$ klatek nie da się odtworzyć na pojedynczych zdjęciach: każdą asercję jadalności liczono tak, jakby konsensus był spełniony, więc CFPR TFLite jest górnym oszacowaniem dla nakładki na żywo. Wiersze macierzy to etykiety rejestru: jadalne / niejadalne / toksyczne. Kolumny to \texttt{displayStatus}: asercja jadalności / brak asercji / ostrzeżenie. Starsza macierz ze zwiniętym wierszem ``nietoksyczne'' (wiersz bezpieczny $61\,906$ instancji, wiersz niebezpieczny $11\,366$; CFPR $284/11\,366=2{,}50\%$) nie jest używana w weryfikacji $H_2$: składała jadalne z niejadalnymi i liczyła inną kolumnę ``SAFE''.

Na $15\,840$ instancjach toksycznych polityka asymetryczna dała $327$ asercje jadalności, a próg globalny $580$:
\begin{equation}
\mathrm{CFPR}_{0{,}18/0{,}60}
=
\frac{327}{15\,840}
=
0{,}0206
\quad (2{,}06\%),\qquad
\mathrm{CFPR}_{\tau=0{,}30}
=
\frac{580}{15\,840}
=
0{,}0366
\quad (3{,}66\%).
\end{equation}
Dla wag PyTorch para progów obniża CFPR względem baseline ($H_2$ spełniona). $15\,513$ obiektów toksycznych ($97{,}94\%$) nie dostało zielonej asercji przy $0{,}18$/$0{,}60$. FAR na klasach jadalnych, $P(\text{ostrzeżenie}\mid\text{gt jadalny})$, wyniósł $1{,}44\%$ ($478/33\,301$) przy polityce asymetrycznej i $1{,}17\%$ ($388/33\,301$) przy $\tau=0{,}30$: niższy CFPR kupuje się niewielkim wzrostem ostrzeżeń na jadalnych. Na artefakcie TFLite FP16 z dekoderem aplikacji CFPR wyniósł $2{,}11\%$ ($334/15\,840$) wobec $3{,}64\%$ ($577/15\,840$) przy $\tau=0{,}30$, a FAR $1{,}38\%$ wobec $1{,}17\%$; $H_2$ pozostaje spełniona także na artefakcie wdrożeniowym. Tab.~\ref{tab:cfpr_porownanie}--\ref{tab:macierz_ryzyka_global}, rys.~\ref{fig:safety_matrix_h2}.

\begin{table}[htbp]
 \centering
 \caption{CFPR i FAR na izolowanym teście: para $0{,}18$/$0{,}60$ wobec jednego progu $\tau=0{,}30$, dla wag PyTorch oraz artefaktu TFLite FP16 z dekoderem aplikacji. Źródło: opracowanie własne na podstawie \texttt{safety\_cfpr\_comparison.csv} i \texttt{tflite\_cfpr.json}.}
 \label{tab:cfpr_porownanie}
 \begin{tabular}{|l|l|c|c|c|c|}
 \hline
 \textbf{Model} & \textbf{Polityka} & \textbf{CFP} & \textbf{CFPR} & \textbf{ostrzeżenia$\mid$jadalne} & \textbf{FAR jadalne} \\
 \hline \hline
 PyTorch & $0{,}18$ / $0{,}60$ & $327/15\,840$ & $2{,}06\%$ & $478/33\,301$ & $1{,}44\%$ \\
 \hline
 PyTorch & $\tau=0{,}30$ & $580/15\,840$ & $3{,}66\%$ & $388/33\,301$ & $1{,}17\%$ \\
 \hline \hline
 TFLite FP16 & $0{,}18$ / $0{,}60$ & $334/15\,840$ & $2{,}11\%$ & $461/33\,301$ & $1{,}38\%$ \\
 \hline
 TFLite FP16 & $\tau=0{,}30$ & $577/15\,840$ & $3{,}64\%$ & $388/33\,301$ & $1{,}17\%$ \\
 \hline
 \end{tabular}
\end{table}

\begin{table}[htbp]
 \centering
 \caption{Macierz ryzyka przy progach $0{,}18$/$0{,}60$ (wagi PyTorch, \texttt{displayStatus}, bez bufora $4$ z $5$). Źródło: opracowanie własne na podstawie \texttt{safety\_matrix\_asymmetric.csv}.}
 \label{tab:macierz_ryzyka}
 \begin{tabular}{|l|c|c|c|c|}
 \hline
 \multirow{2}{*}{\textbf{Rejestr (gt)}} & \multicolumn{3}{c|}{\textbf{Predykcja}} & \multirow{2}{*}{\textbf{Razem}} \\
 \cline{2-4}
 & Jadalny & NEUTRAL & ostrzeżenie & \\
 \hline \hline
 jadalne & $21\,228$ & $11\,595$ & $478$ & $33\,301$ \\
 \hline
 niejadalne & $159$ & $23\,776$ & $196$ & $24\,131$ \\
 \hline
 toksyczne & $\mathbf{327}$ & $3\,108$ & $12\,405$ & $15\,840$ \\
 \hline
 \end{tabular}
\end{table}

\begin{table}[htbp]
 \centering
 \caption{Macierz ryzyka przy jednym progu $\tau=0{,}30$ (wagi PyTorch, ten sam test i to samo dopasowanie IoU). Źródło: opracowanie własne na podstawie \texttt{safety\_matrix\_global.csv}.}
 \label{tab:macierz_ryzyka_global}
 \begin{tabular}{|l|c|c|c|c|}
 \hline
 \multirow{2}{*}{\textbf{Rejestr (gt)}} & \multicolumn{3}{c|}{\textbf{Predykcja}} & \multirow{2}{*}{\textbf{Razem}} \\
 \cline{2-4}
 & Jadalny & NEUTRAL & ostrzeżenie & \\
 \hline \hline
 jadalne & $25\,746$ & $7\,167$ & $388$ & $33\,301$ \\
 \hline
 niejadalne & $324$ & $23\,641$ & $166$ & $24\,131$ \\
 \hline
 toksyczne & $\mathbf{580}$ & $3\,547$ & $11\,713$ & $15\,840$ \\
 \hline
 \end{tabular}
\end{table}

\begin{figure}[htbp]
 \centering
 \begin{subfigure}[t]{0.48\textwidth}
  \centering
  \includegraphics[width=\linewidth]{wyniki/safety_matrix_asymmetric}
  \caption{Para $0{,}18$/$0{,}60$, $\mathrm{CFPR}=2{,}06\%$.}
  \label{fig:safety_matrix_asymm}
 \end{subfigure}
 \hfill
 \begin{subfigure}[t]{0.48\textwidth}
  \centering
  \includegraphics[width=\linewidth]{wyniki/safety_matrix_global}
  \caption{Próg $\tau=0{,}30$, $\mathrm{CFPR}=3{,}66\%$.}
  \label{fig:safety_matrix_glob}
 \end{subfigure}
 \caption{Macierze ryzyka na izolowanym teście (wagi PyTorch). Wiersze: EDIBLE / INEDIBLE / TOXIC. Kolumny: JADALNY / NEUTRAL / OSTRZEZENIE. Źródło: opracowanie własne.}
 \label{fig:safety_matrix_h2}
\end{figure}

Uzyskany CFPR $2{,}06\%$ (PyTorch) należy interpretować ostrożnie. $327$ z $15\,840$ obiektów toksycznych nadal dostało asercję jadalności: warstwa progów redukuje to ryzyko względem $\tau=0{,}30$ ($580$ przypadków), lecz go nie eliminuje i nie jest warunkiem wystarczającym do spożycia. Na przeciekającym podziale testowym ta sama polityka dawała $1{,}06\%$ przy innej definicji kubełków -- te liczby nie są porównywalne jeden do jednego z tab.~\ref{tab:cfpr_porownanie}. Bufor głosowania czasowego ($4$ z $5$ klatek) nie wchodzi w pomiar CFPR ani w treść $H_2$.

\textbf{Pochodzenie progów.} Para $(\mathrm{conf}_{\mathrm{toxic}},\,\mathrm{conf}_{\mathrm{edible}})=(0{,}18,\,0{,}60)$ została wpisana do rejestru aplikacji (\texttt{MushroomRegistry.kt}) 11 sierpnia 2026~r., czyli przed zbudowaniem izolowanego zbioru testowego (skrypty \texttt{harvest\_clean\_test.py} i \texttt{label\_clean\_test.py}, 18--19 sierpnia 2026~r.). Zastąpiła ona wcześniejsze wartości robocze, w tym jeden próg globalny $0{,}304$ opisany niżej. Wybór opierał się na asymetrycznej przesłance bezpieczeństwa: możliwie wczesne ostrzeganie o~gatunkach toksycznych (niski próg $0{,}18$) i~konserwatywna asercja jadalności (wysoki próg $0{,}60$). Przy ustalaniu progów dostępne były jedynie wyniki z~poprzedniego podziału testowego, który -- jak wykazano później -- dzielił obserwacje (\texttt{obs\_id}) z~korpusem uczącym. Nie można więc wykluczyć, że wartości $0{,}18$ i $0{,}60$ były pośrednio inspirowane krzywymi z~tamtego, przeciekającego podziału. Wyniki izolowanego testu -- przeciągnięcie progu z~tab.~\ref{tab:threshold_sweep} i~macierze z~tab.~\ref{tab:cfpr_porownanie} -- powstały po zamrożeniu progów i~nie posłużyły do ich zmiany.

Przeciągnięcie progu (ang. \textit{confidence threshold sweep}) w~zakresie $0{,}05$--$0{,}95$ (skrypt \texttt{test3\_threshold\_sweep.py}) charakteryzuje zatem kompromis już przyjętych progów, a~nie jest procedurą ich doboru. Para $0{,}18$/$0{,}60$ jest przy tym \emph{nieoptymalna} względem $F_1$ ($F_1=0{,}7801$ utrzymuje się na całym przedziale $\mathrm{conf}=0{,}05$--$0{,}20$). W~zakresie $\mathrm{conf}\le 0{,}20$ raportowane przez Ultralytics $P$, $R$ i~$F_1$ pozostają praktycznie stałe (tab.~\ref{tab:threshold_sweep}), podczas gdy $mAP_{50}$ maleje monotonicznie; z~samej tabeli nie wynika więc, że próg $0{,}05$ generuje więcej fałszywych alarmów niż $0{,}20$.

We wcześniejszej iteracji polityki progów, poprzedzającej wprowadzenie podziału asymetrycznego, aplikacja korzystała z pojedynczego, globalnego progu operacyjnego $\mathrm{conf}=0{,}304$. Najbliższy zarejestrowany węzeł przeciągnięcia na izolowanym teście ($\mathrm{conf}=0{,}30$) daje: precyzja $P=80{,}41\%$, czułość $R=75{,}18\%$, miara $F_1=77{,}71\%$ oraz $mAP_{50}=71{,}01\%$. Pojedynczy próg globalny tego rzędu równoważy $P$ i $R$ w ujęciu zagregowanym, nie różnicuje jednak ryzyka konsumpcyjnego pomiędzy gatunkami jadalnymi a trującymi. Dlatego zastąpiono go parą $(\mathrm{conf}_{\mathrm{toxic}},\,\mathrm{conf}_{\mathrm{edible}})=(0{,}18,\,0{,}60)$. Jako docelowy punkt asercji jadalności przyjęto $\mathrm{conf}_{\mathrm{edible}}=0{,}60$, natomiast próg ostrzeżeń toksycznych ustalono na $\mathrm{conf}_{\mathrm{toxic}}=0{,}18$. Wybrane węzły przeciągnięcia zestawiono w tab.~\ref{tab:threshold_sweep}. Przebieg $F_1$ przedstawia rys.~\ref{fig:threshold_sweep_f1}.

\begin{table}[htbp]
 \centering
 \caption{Wybrane węzły przeciągnięcia progu pewności detekcji w zakresie $0{,}05$--$0{,}95$. Źródło: opracowanie własne na podstawie \texttt{threshold\_sweep.csv}.}
 \label{tab:threshold_sweep}
 \begin{tabular}{|c|c|c|c|c|p{4.4cm}|}
 \hline
 $\mathbf{conf}$ & $\boldsymbol{P}$ & $\boldsymbol{R}$ & $\boldsymbol{F_1}$ & $\boldsymbol{mAP_{50}}$ & \textbf{Uwagi} \\
 \hline \hline
 $0{,}05$ & $0{,}8090$ & $0{,}7532$ & $\mathbf{0{,}7801}$ & $0{,}7926$ & początek plateau $F_1$ \\
 \hline
 $0{,}15$ & $0{,}8090$ & $0{,}7532$ & $0{,}7801$ & $0{,}7538$ & węzeł poniżej progu toksycznego $0{,}18$ \\
 \hline
 $0{,}20$ & $0{,}8090$ & $0{,}7532$ & $0{,}7801$ & $0{,}7390$ & dolna granica pasma neutralnego \\
 \hline
 $0{,}30$ & $0{,}8041$ & $0{,}7518$ & $0{,}7771$ & $0{,}7101$ & dawny próg globalny \\
 \hline
 $0{,}45$ & $0{,}8488$ & $0{,}6988$ & $0{,}7665$ & $0{,}6668$ & --- \\
 \hline
 $0{,}55$ & $0{,}8726$ & $0{,}6582$ & $0{,}7504$ & $0{,}6325$ & węzeł przed progiem jadalnym $0{,}60$ \\
 \hline
 $0{,}60$ & $0{,}8827$ & $0{,}6352$ & $0{,}7388$ & $0{,}6129$ & \textbf{próg asercji jadalności} \\
 \hline
 $0{,}75$ & $0{,}9125$ & $0{,}5448$ & $0{,}6822$ & $0{,}5326$ & --- \\
 \hline
 $0{,}90$ & $0{,}9354$ & $0{,}3604$ & $0{,}5204$ & $0{,}3570$ & załamanie czułości \\
 \hline
 \end{tabular}
\end{table}

Zestawienie ilustruje kompromis, jaki niesie przyjęty wcześniej próg asercji jadalności: przejście od $\mathrm{conf}=0{,}30$ do $\mathrm{conf}=0{,}60$ podnosi precyzję z $80{,}41\%$ do $88{,}27\%$ kosztem spadku czułości z $75{,}18\%$ do $63{,}52\%$. W zastosowaniu ogólnym byłby to kompromis niekorzystny, ponieważ miara $F_1$ maleje z $0{,}7771$ do $0{,}7388$. W rozważanym zadaniu utrata czułości oznacza jednak wyłącznie brak zielonej etykiety dla części owocników rzeczywiście jadalnych (niedogodność, nie zagrożenie), podczas gdy przyrost precyzji bezpośrednio redukuje liczbę błędnych asercji jadalności. Symetryczna miara $F_1$ nie jest zatem właściwym kryterium oceny progu w warstwie decyzyjnej tego systemu.

\begin{figure}[htbp]
 \centering
 \includegraphics[width=0.85\textwidth]{wyniki/threshold_sweep_f1}
 \caption{Zależność miary $F_1$ od progu pewności detekcji w zakresie $0{,}05$--$0{,}95$; punkt asercji jadalności aplikacji $\mathrm{conf}_{\mathrm{edible}}=0{,}60$. Źródło: opracowanie własne.}
 \label{fig:threshold_sweep_f1}
\end{figure}

Skuteczność w przypadkach bliźniaczych nie sprowadza się wyłącznie do zdolności sieci do rozróżniania klas: zależy od sprzężenia detektora z polityką progów. Tab.~\ref{tab:cfpr_porownanie} pokazuje, że para $0{,}18$/$0{,}60$ obniża CFPR z $3{,}66\%$ do $2{,}06\%$ względem $\tau=0{,}30$ na tym samym teście. $H_2$ w tym ujęciu jest spełniona. Pozostałe $2{,}06\%$ przypomina, że system wspomaga identyfikację, a nie zapewnia bezpieczeństwa konsumpcyjnego. Ponieważ progi ustalono przed zbudowaniem izolowanego testu, a~jego wyniki nie posłużyły do ich zmiany, porównanie $H_2$ nie jest obciążone selekcją progów na danych testowych. Nie wyklucza to wpływu wcześniejszych wyników z~przeciekającego podziału testowego na wybór wartości $0{,}18$ i~$0{,}60$. Wynik $H_2$ należy więc traktować jako potwierdzenie na niezależnych danych polityki dobranej heurystycznie, a~nie jako dowód jej optymalności. Czystszym rozwiązaniem byłby osobny zbiór walidacyjny przeznaczony wyłącznie do kalibracji warstwy decyzyjnej.

\subsection{Pary gatunków bliźniaczych}\label{sec:lookalikes}
We wstępie wymieniono pary o wysokim koszcie pomyłki. Na izolowanym teście \texttt{test\_h2\_baseline.py} zliczał oba kierunki przy $\mathrm{IoU}\ge 0{,}50$ (tab.~\ref{tab:lookalike_pary}). Licznik nie jest CFPR: tu chodzi o konkretną parę taksonów, nie o asercję jadalności na wszystkich klasach toksycznych.

Najczęstsza pomyłka spośród tych par to \textit{Kuehneromyces mutabilis} $\to$ \textit{Galerina marginata} ($42/1\,018$, $4{,}13\%$) oraz odwrotnie ($13/808$, $1{,}61\%$). Muchomor sromotnikowy nie został pomylony z pieczarką polną ($0/390$ i $0/420$); z muchomorem cytrynowym: $7/390$ ($1{,}79\%$) i $4/305$ ($1{,}31\%$). Kurki nie sklasyfikowano jako lisówki ($0/527$); dwie instancje kurki dostały etykietę \textit{Omphalotus olearius} ($2/527$, $0{,}38\%$), w drugą stronę $0/482$. Gatunki występujące w dwóch parach (\textit{A.\ phalloides}, \textit{A.\ campestris}, \textit{C.\ cibarius}) liczono raz na instancję.

\begin{table}[htbp]
 \centering
 \caption{Pomyłki pomiędzy parami bliźniaczymi na izolowanym teście. $n_{\mathrm{gt}}$ to instancje gatunku referencyjnego po dopasowaniu IoU. Źródło: opracowanie własne na podstawie \texttt{lookalike\_confusions.csv}.}
 \label{tab:lookalike_pary}
 \begin{tabular}{|p{4.2cm}|p{4.2cm}|r|r|r|}
 \hline
 \textbf{gt} & \textbf{pred} & $\boldsymbol{n_{\mathrm{gt}}}$ & \textbf{pomyłki} & \textbf{udział} \\
 \hline \hline
 \textit{A.\ phalloides} & \textit{A.\ campestris} & $390$ & $0$ & $0{,}00\%$ \\
 \hline
 \textit{A.\ campestris} & \textit{A.\ phalloides} & $420$ & $0$ & $0{,}00\%$ \\
 \hline
 \textit{A.\ phalloides} & \textit{A.\ citrina} & $390$ & $7$ & $1{,}79\%$ \\
 \hline
 \textit{A.\ citrina} & \textit{A.\ phalloides} & $305$ & $4$ & $1{,}31\%$ \\
 \hline
 \textit{A.\ virosa} & \textit{A.\ campestris} & $127$ & $1$ & $0{,}79\%$ \\
 \hline
 \textit{A.\ campestris} & \textit{A.\ virosa} & $420$ & $0$ & $0{,}00\%$ \\
 \hline
 \textit{O.\ olearius} & \textit{C.\ cibarius} & $482$ & $0$ & $0{,}00\%$ \\
 \hline
 \textit{C.\ cibarius} & \textit{O.\ olearius} & $527$ & $2$ & $0{,}38\%$ \\
 \hline
 \textit{H.\ aurantiaca} & \textit{C.\ cibarius} & $464$ & $0$ & $0{,}00\%$ \\
 \hline
 \textit{C.\ cibarius} & \textit{H.\ aurantiaca} & $527$ & $0$ & $0{,}00\%$ \\
 \hline
 \textit{G.\ marginata} & \textit{K.\ mutabilis} & $808$ & $13$ & $1{,}61\%$ \\
 \hline
 \textit{K.\ mutabilis} & \textit{G.\ marginata} & $1\,018$ & $42$ & $4{,}13\%$ \\
 \hline
 \end{tabular}
\end{table}

\section{Test odporności na zakłócenia leśne}
W warunkach terenowych jakość obrazu z kamery smartfonu jest systematycznie degradowana przez czynniki optyczne i środowiskowe: obniżone oświetlenie podszytu, rozmycie wynikające z ruchu ręki oraz szum elektroniczny matrycy przy podwyższonej czułości ISO. W celu ilościowej oceny odporności detektora na takie zaburzenia przeprowadzono test obciążeniowy (ang. \textit{stress test}) zaimplementowany w skrypcie \texttt{test2\_robustness\_noise.py}. Na izolowanym zbiorze testowym zasymulowano jednoczesne działanie trzech czynników: ściemnienia obrazu ze współczynnikiem jasności $0{,}6$ (redukcja o $40\%$), rozmycia gaussowskiego z jądrem $5\times5$, przybliżającego nieostrość kadru oraz addytywnego szumu Gaussa o odchyleniu standardowym $\sigma=8{,}0$, odpowiadającego typowemu szumowi matrycy przy podwyższonym ISO. Parametry symulacji oraz uzyskane wyniki zestawiono w tab.~\ref{tab:odpornosc}.

\begin{table}[htbp]
 \centering
 \caption{Parametry symulacji zakłóceń leśnych oraz degradacja skuteczności detekcji. Źródło: opracowanie własne na podstawie \texttt{robustness\_noise\_summary.json}.}
 \label{tab:odpornosc}
 \begin{tabular}{|p{5.0cm}|c|p{5.0cm}|}
 \hline
 \textbf{Czynnik zakłócający} & \textbf{Parametr} & \textbf{Odwzorowywane zjawisko} \\
 \hline \hline
 Ściemnienie obrazu & $0{,}6$ & zacienienie podszytu leśnego \\
 \hline
 Rozmycie gaussowskie & jądro $5\times5$ & przybliżenie nieostrości kadru \\
 \hline
 Addytywny szum Gaussa & $\sigma=8{,}0$ & szum matrycy przy wysokim ISO \\
 \hline \hline
 \textbf{Wynik} & \textbf{Wartość} & \textbf{Komentarz} \\
 \hline \hline
 $mAP_{50}$ – obrazy czyste & $0{,}8275$ & wynik odniesienia \\
 \hline
 $mAP_{50}$. Obrazy zakłócone & $0{,}6145$ & wszystkie trzy czynniki łącznie \\
 \hline
 Degradacja bezwzględna & $21{,}30$\,pp & różnica wartości $mAP_{50}$ \\
 \hline
 Degradacja względna & $25{,}74\%$ & odniesiona do wyniku czystego \\
 \hline
 \end{tabular}
\end{table}

\begin{figure}[htbp]
 \centering
 \begin{tikzpicture}
 \def\sk{5.6} % cm na jednostke mAP
 % os pozioma
 \draw[semithick] (0,-0.55) -- (1.0*\sk,-0.55);
 \foreach \x/\e in {0/0{,}0, 0.2/0{,}2, 0.4/0{,}4, 0.6/0{,}6, 0.8/0{,}8, 1.0/1{,}0}{
 \draw[semithick] (\x*\sk,-0.55) -- (\x*\sk,-0.68) node[below, font=\scriptsize] {$\e$};
 }
 \node[font=\scriptsize, anchor=north] at (0.5*\sk,-1.1) {$mAP_{50}$};
 % slupki
 \node[font=\scriptsize, anchor=east] at (-0.15,0.85) {obrazy czyste};
 \fill[EEBlue!45, draw=black!50] (0,0.55) rectangle (0.8275*\sk,1.15);
 \node[font=\scriptsize, anchor=west] at (0.8275*\sk+0.12,0.85) {$0{,}8275$};

 \node[font=\scriptsize, anchor=east] at (-0.15,0.0) {obrazy zakłócone};
 \fill[EETangerine!60, draw=black!50] (0,-0.3) rectangle (0.6145*\sk,0.3);
 \node[font=\scriptsize, anchor=west] at (0.6145*\sk+0.12,0.0) {$0{,}6145$};

 % zaznaczenie degradacji
 \draw[latex-latex, semithick, red!70]
 (0.6145*\sk,1.55) -- node[above, font=\scriptsize, red!70] {$-21{,}30$\,pp} (0.8275*\sk,1.55);
 \draw[dashed, red!50] (0.6145*\sk,0.3) -- (0.6145*\sk,1.55);
 \draw[dashed, red!50] (0.8275*\sk,1.15) -- (0.8275*\sk,1.55);
 \end{tikzpicture}
 \caption{Degradacja skuteczności detekcji pod wpływem skumulowanych zakłóceń leśnych. Źródło: opracowanie własne na podstawie \texttt{robustness\_noise\_summary.json}.}
 \label{fig:odpornosc_wykres}
\end{figure}

Na obrazach niezakłóconych model YOLO11m osiągnął $mAP_{50}=0{,}8275$ ($82{,}75\%$). Po wprowadzeniu skumulowanych zakłóceń leśnych wartość ta spadła do $mAP_{50}=0{,}6145$ ($61{,}45\%$), co zilustrowano na rys.~\ref{fig:odpornosc_wykres}. Absolutna degradacja wyniosła zatem $21{,}30$ punktu procentowego, a degradacja względna – $25{,}74\%$ względem wyniku bazowego. Mimo istotnego spadku dokładności detektor zachował użyteczność operacyjną na poziomie $mAP_{50}>61\%$, co wskazuje, że reprezentacje cech wyuczone na zbiorze iNaturalist generalizują się do warunków o obniżonej jakości obrazu. Należy jednak zastrzec, że eksperyment przeprowadzono wyłącznie ze skumulowanym pakietem zakłóceń leśnych; wpływ każdego czynnika w izolacji (np. samego ściemnienia, samego rozmycia czy samego szumu Gaussa) nie był badany osobno, co stanowi ograniczenie wnioskowania o odporności w specyficznych warunkach pogodowych. Wynik ten uzasadnia jednocześnie konserwatywną politykę progów w aplikacji mobilnej: w warunkach leśnych część predykcji naturalnie traci pewność, a warstwa \texttt{displayStatus} zapobiega prezentowaniu niepewnych wskazań jadalności jako zielonego statusu „Jadalny”.

\section{Analiza porównawcza architektur (studium ablacyjne)}
Wybór wariantu YOLO11m zamiast lżejszego YOLO11n wymagał empirycznego uzasadnienia, ponieważ na urządzeniach mobilnych naturalną tendencją projektową jest minimalizacja liczby parametrów i rozmiaru artefaktu TFLite. Przeprowadzono zatem studium ablacyjne (ang. \textit{ablation study}), w którym obie architektury wytrenowano na tym samym zadaniu rozpoznawania $147$ klas, a następnie zmierzono na identycznym izolowanym teście po \texttt{obs\_id} (tab.~\ref{tab:podzial_zbioru}). Metryki w tab.~\ref{tab:ablation_yolo11} pochodzą z tego samego protokołu Ultralytics \texttt{val} (\texttt{split=test}, \texttt{imgsz=640}). Nano osiągnął $mAP_{50}=48{,}54\%$ oraz $mAP_{50\text{--}95}=42{,}82\%$; Medium -- $82{,}75\%$ / $78{,}23\%$. Różnica $34{,}21$ punktu procentowego na $mAP_{50}$ jest luką pojemności reprezentacyjnej, a nie artefaktem protokołu ewaluacji.

\begin{table}[htbp]
 \centering
 \caption{Porównanie ablacyjne YOLO11m i YOLO11n na izolowanym teście po \texttt{obs\_id} ($32\,727$ obrazów). Źródło: opracowanie własne.}
 \label{tab:ablation_yolo11}
 \begin{tabular}{|l|c|c|}
 \hline
 \textbf{Parametr} & \textbf{YOLO11m (wybrany)} & \textbf{YOLO11n (bazowy)} \\
 \hline \hline
 Liczba parametrów & $20{,}14$\,M & $2{,}66$\,M \\
 \hline
 Złożoność obliczeniowa & $68{,}3$\,GFLOPs & $6{,}7$\,GFLOPs \\
 \hline
 Rozmiar TFLite FP16 & $38{,}58$\,MiB & $5{,}60$\,MiB \\
 \hline
 $mAP_{50}$ & $82{,}75\%$ & $48{,}54\%$ \\
 \hline
 $mAP_{50\text{--}95}$ & $78{,}23\%$ & $42{,}82\%$ \\
 \hline
 Precyzja $P$ & $80{,}90\%$ & $51{,}49\%$ \\
 \hline
 Czułość $R$ & $75{,}32\%$ & $48{,}71\%$ \\
 \hline
 Opóźnienie inferencji (GPU) & $8{,}62$\,ms & $1{,}10$\,ms \\
 \hline
 \end{tabular}
\end{table}

\begin{figure}[htbp]
 \centering
 \begin{tikzpicture}
 \def\sk{0.058} % cm na punkt procentowy
 \def\gw{0.42} % wysokosc slupka
 % osie
 \draw[semithick] (0,-1.75) -- (100*\sk,-1.75);
 \foreach \x in {0,20,40,60,80,100}{
 \draw[semithick] (\x*\sk,-1.75) -- (\x*\sk,-1.88) node[below, font=\scriptsize] {\x};
 }
 \node[font=\scriptsize, anchor=north] at (50*\sk,-2.3) {wartość metryki [\%]};

 % mAP50
 \node[font=\scriptsize, anchor=east] at (-0.15,0.15) {$mAP_{50}$};
 \fill[EEBlue!45, draw=black!50] (0,0.15) rectangle (82.75*\sk,0.15+\gw);
 \node[font=\scriptsize, anchor=west] at (82.75*\sk+0.1,0.15+\gw/2) {$82{,}75$};
 \fill[EETangerine!60, draw=black!50] (0,0.15-\gw) rectangle (48.54*\sk,0.15);
 \node[font=\scriptsize, anchor=west] at (48.54*\sk+0.1,0.15-\gw/2) {$48{,}54$};

 % mAP50-95
 \node[font=\scriptsize, anchor=east] at (-0.15,-1.0) {$mAP_{50\text{--}95}$};
 \fill[EEBlue!45, draw=black!50] (0,-1.0) rectangle (78.23*\sk,-1.0+\gw);
 \node[font=\scriptsize, anchor=west] at (78.23*\sk+0.1,-1.0+\gw/2) {$78{,}23$};
 \fill[EETangerine!60, draw=black!50] (0,-1.0-\gw) rectangle (42.82*\sk,-1.0);
 \node[font=\scriptsize, anchor=west] at (42.82*\sk+0.1,-1.0-\gw/2) {$42{,}82$};

 % legenda
 \fill[EEBlue!45, draw=black!50] (0,1.35) rectangle (0.42,1.62);
 \node[font=\scriptsize, anchor=west] at (0.52,1.485) {YOLO11m (wybrany)};
 \fill[EETangerine!60, draw=black!50] (3.4,1.35) rectangle (3.82,1.62);
 \node[font=\scriptsize, anchor=west] at (3.92,1.485) {YOLO11n (bazowy)};
 \end{tikzpicture}
 \caption{Porównanie $mAP$ wariantów YOLO11m i YOLO11n na izolowanym teście po \texttt{obs\_id}. Źródło: opracowanie własne.}
 \label{fig:ablacja_wykres}
\end{figure}

Wariant Nano, mimo ponad siedmiokrotnie mniejszej liczby parametrów i prawie siedmiokrotnie mniejszego artefaktu FP16, osiągnął $48{,}54\%$ $mAP_{50}$ na tym samym izolowanym teście (rys.~\ref{fig:ablacja_wykres}). YOLO11m ma $82{,}75\%$. Precyzja i czułość Nano ($P=51{,}49\%$, $R=48{,}71\%$) pozostają wyraźnie poniżej Medium ($P=80{,}90\%$, $R=75{,}32\%$). Nano nie modeluje $147$ klas, w tym gatunków bliźniaczych, z akceptowalną rozdzielczością semantyczną (blaszki, pierścień, łuski, barwa trzonu). Opóźnienie GPU wariantu Medium ($8{,}62$\,ms na RTX~4070) nie jest pomiarem telefonu; docelowa kadencja analizy to $250$\,ms, a kryterium $H_1$ to czas reakcji poniżej $1$\,s (tab.~\ref{tab:latency_e2e}). Zysk prędkości Nano na GPU nie rekompensuje utraty jakości. Dlatego YOLO11m przyjęto jako wariant produkcyjny.

\section{Profilowanie opóźnień potoku inferencji}
W celu rozdzielenia kosztów obliczeniowych poszczególnych etapów potoku przeprowadzono profilowanie na akceleratorze NVIDIA GeForce RTX~4070 przy wejściu $640{\times}640$ pikseli. Pomiar obejmował $100$ przebiegów i rozróżniał trzy fazy: wstępne przetwarzanie obrazu, inferencję sieci oraz postprocessing z tłumieniem niemaksymalnym (NMS). Średni czas preprocessingu wyniósł $1{,}64$\,ms (mediana $p_{50}=1{,}68$\,ms, percentyl $p_{95}=1{,}76$\,ms), średni czas inferencji $8{,}62$\,ms ($p_{50}=8{,}62$\,ms, $p_{95}=8{,}75$\,ms), a średni czas NMS $0{,}50$\,ms ($p_{50}=0{,}50$\,ms, $p_{95}=0{,}52$\,ms). Pełny rozkład statystyczny czasów poszczególnych faz zestawiono w tab.~\ref{tab:latency}.

\begin{table}[htbp]
 \centering
 \caption{Profilowanie opóźnień potoku inferencji na akceleratorze NVIDIA GeForce RTX~4070 ($100$ przebiegów, wejście $640\times640$). Wszystkie wartości w milisekundach. Źródło: opracowanie własne na podstawie \texttt{latency\_breakdown.json}.}
 \label{tab:latency}
 \begin{tabular}{|l|c|c|c|c|c|}
 \hline
 \textbf{Faza potoku} & \textbf{Średnia} & $\boldsymbol{\sigma}$ & $\boldsymbol{p_{50}}$ & $\boldsymbol{p_{95}}$ & \textbf{Maksimum} \\
 \hline \hline
 Wstępne przetwarzanie & $1{,}64$ & $0{,}08$ & $1{,}68$ & $1{,}76$ & $1{,}80$ \\
 \hline
 Inferencja sieci & $8{,}62$ & $0{,}08$ & $8{,}62$ & $8{,}75$ & $8{,}96$ \\
 \hline
 Postprocessing (NMS) & $0{,}50$ & $0{,}01$ & $0{,}50$ & $0{,}52$ & $0{,}54$ \\
 \hline \hline
 \textbf{Całkowity czas e2e} & $\mathbf{10{,}94}$ & $0{,}16$ & $10{,}98$ & $11{,}15$ & $11{,}45$ \\
 \hline \hline
 \multicolumn{6}{|l|}{Przepustowość równoważna: $1000/10{,}94 = 91{,}4$\,FPS} \\
 \hline
 \multicolumn{6}{|l|}{Wynik na karcie PC; pomiar na S23 Ultra: tab.~\ref{tab:latency_s23}} \\
 \hline
 \end{tabular}
\end{table}

\begin{figure}[htbp]
 \centering
 \begin{tikzpicture}
 \def\sk{0.70} % cm na milisekunde
 \def\gw{0.62}
 % skladowe czasu sredniego
 \fill[EEBlue!35, draw=black!50] (0,0) rectangle (1.64*\sk,\gw);
 \fill[EEUltramarine!45, draw=black!50] (1.64*\sk,0) rectangle (10.26*\sk,\gw);
 \fill[EETangerine!55, draw=black!50] (10.26*\sk,0) rectangle (10.76*\sk,\gw);
 \node[font=\scriptsize, anchor=east] at (-0.15,\gw/2) {czas średni};

 % etykiety skladowych
 \node[font=\scriptsize] at (0.82*\sk,\gw+0.32) {$1{,}64$};
 \node[font=\scriptsize] at (5.95*\sk,\gw+0.32) {$8{,}62$};
 \node[font=\scriptsize] at (10.51*\sk,\gw+0.72) {$0{,}50$};
 \draw[gray, thin] (10.51*\sk,\gw+0.58) -- (10.51*\sk,\gw+0.08);

 % os
 \draw[semithick] (0,-0.3) -- (12*\sk,-0.3);
 \foreach \x in {0,2,4,6,8,10,12}{
 \draw[semithick] (\x*\sk,-0.3) -- (\x*\sk,-0.43) node[below, font=\scriptsize] {\x};
 }
 \node[font=\scriptsize, anchor=north] at (6*\sk,-0.88) {czas [ms]};

 % legenda
 \fill[EEBlue!35, draw=black!50] (0,-2.15) rectangle (0.4,-1.88);
 \node[font=\scriptsize, anchor=west] at (0.5,-2.015) {wstępne przetwarzanie};
 \fill[EEUltramarine!45, draw=black!50] (4.0,-2.15) rectangle (4.4,-1.88);
 \node[font=\scriptsize, anchor=west] at (4.5,-2.015) {inferencja sieci};
 \fill[EETangerine!55, draw=black!50] (7.3,-2.15) rectangle (7.7,-1.88);
 \node[font=\scriptsize, anchor=west] at (7.8,-2.015) {NMS};
 \end{tikzpicture}
 \caption{Udział poszczególnych faz w średnim opóźnieniu potoku inferencji na akceleratorze NVIDIA GeForce RTX~4070. Źródło: opracowanie własne na podstawie \texttt{latency\_breakdown.json}.}
 \label{fig:latency_wykres}
\end{figure}

Całkowite opóźnienie ściany zegara (ang. \textit{end-to-end wall clock}) na stacji roboczej wyniosło średnio $10{,}94$\,ms ($p_{50}=10{,}98$\,ms, $p_{95}=11{,}15$\,ms), co odpowiada przepustowości rzędu $91{,}4$\,FPS. Dominującym składnikiem pozostaje inferencja sieci ($\approx 79\%$ średniego czasu e2e), natomiast preprocessing i NMS wnoszą łącznie ok.\ $2{,}1$\,ms w typowym przebiegu; proporcje te zilustrowano na rys.~\ref{fig:latency_wykres}. Wynik RTX~4070 jest górnym ograniczeniem wydajności potoku algorytmicznego. Hipotezę $H_1$ rozstrzyga jednak pomiar na telefonie, a nie na karcie graficznej PC.

Dlatego powtórzono profilowanie na urządzeniu produkcyjnym: Samsung Galaxy S23 Ultra (SM-S918U1, Snapdragon~8~Gen~2, Android~16). W pierwszej serii mierzono wyłącznie wywołanie \texttt{MushroomDetector.detect()} zegarem monotonicznym \texttt{SystemClock.elapsedRealtimeNanos}. Wejście było identyczne z produkcją ($640\times640$ po letterboxingu CameraX). Pierwsze $8$ klatek odrzucono jako rozgrzewkę interpretera TFLite; zapisano $1\,899$ próbek przy wymuszonej kadencji $4$\,FPS (heartbeat $1{,}5$\,s wyłączono, aby SoC nie ostygł na statycznej scenie). Seria zimna to pierwsze $150$ klatek po rozgrzewce, w temperaturze pokojowej (bateria $31{,}2^{\circ}\mathrm{C}$). Seria gorąca to ostatnie $150$ klatek po ok.\ $8$--$9$\,min ciągłej inferencji (bateria $36{,}3^{\circ}\mathrm{C}$). Rozkład przedstawiono w tab.~\ref{tab:latency_s23}.

\begin{table}[htbp]
 \centering
 \caption{Czas wywołania \texttt{detect()} na Samsung Galaxy S23 Ultra. Wartości w milisekundach. Źródło: opracowanie własne na podstawie \texttt{inference\_profile.csv}.}
 \label{tab:latency_s23}
 \begin{tabular}{|l|c|c|c|}
 \hline
 \textbf{Wielkość} & \textbf{Zimny SoC} & \textbf{Po obciążeniu} & \textbf{Cała seria} \\
 \hline
  & ($n=150$) & ($n=150$) & ($n=1\,899$) \\
 \hline \hline
 Średnia & $156{,}5$ & $245{,}9$ & $216{,}4$ \\
 \hline
 $\sigma$ & $11{,}7$ & $36{,}9$ & $43{,}8$ \\
 \hline
 $p_{50}$ & $159{,}3$ & $262{,}5$ & $232{,}0$ \\
 \hline
 $p_{95}$ & $171{,}5$ & $278{,}3$ & $272{,}3$ \\
 \hline
 Minimum & $129{,}1$ & $145{,}3$ & $129{,}1$ \\
 \hline
 Maksimum & $177{,}2$ & $287{,}4$ & $406{,}4$ \\
 \hline
 Udział klatek $>250$\,ms & $0\%$ & $64\%$ & $27\%$ \\
 \hline
 Temperatura baterii & $31{,}2^{\circ}\mathrm{C}$ & $36{,}3^{\circ}\mathrm{C}$ & $31{,}2$--$36{,}3^{\circ}\mathrm{C}$ \\
 \hline
 \end{tabular}
\end{table}

Na zimnym SoC średnie $156{,}5$\,ms ($p_{95}=171{,}5$\,ms) mieści się w interwale throttlingu $250$\,ms: żadna z $150$ klatek nie przekroczyła tego progu. Po kilkuminutowym obciążeniu średnia rośnie do $245{,}9$\,ms, a $p_{95}$ do $278{,}3$\,ms; $64\%$ klatek gorącej serii przekracza $250$\,ms. Te liczby dotyczą wyłącznie \texttt{detect()}; pełny cykl klatki zmierzono osobno (podrozdz.~\ref{sec:latency_e2e}). Przyrost o ok.\ $90$\,ms i wzrost temperatury baterii o $5{,}1$\,K to thermal throttling Snapdragona bez aktywnego chłodzenia (rozdz.~\ref{chap:przeglad_rozwiazan}). W produkcji heartbeat $1{,}5$\,s po zablokowaniu sceny ogranicza ten reżim: ciągłe $4$\,FPS przez ok.\ $8$--$9$\,min jest scenariuszem obciążeniowym, nie typową sesją zbieracza. Sesja leśna z~Trybu Badawczego (podrozdz.~\ref{sec:field_quant}) jest bliższa temu drugiemu: bateria $27{,}7$--$30{,}8^{\circ}\mathrm{C}$, delegat GPU włączony, \texttt{detect()} $136$--$196$\,ms. Przekroczenie $250$\,ms przez \texttt{detect()} obniża chwilową kadencję analizy, ale nie rozstrzyga o $H_1$: kryterium czasu rzeczywistego (podrozdz.~\ref{sec:wymagania_niefunkcjonalne}) dotyczy czasu reakcji poniżej $1$\,s, weryfikowanego w kolejnym podrozdziale.

\subsection{Opóźnienie end-to-end i czas reakcji na S23 Ultra}\label{sec:latency_e2e}
Pomiar \texttt{detect()} nie obejmuje konwersji klatki CameraX do bitmapy, letterboxingu, filtra progów, \texttt{DetectionTracker} ani publikacji stanu nakładki. Dlatego profiler debugowy (\texttt{InferenceProfiler}) rozszerzono o znaczniki czasu \texttt{SystemClock.elapsedRealtimeNanos} na granicach etapów potoku. Opóźnienie end-to-end liczone jest od wywołania analizatora CameraX z nową klatką do opublikowania nowego \texttt{MushroomUiState}; nie obejmuje ono czasu ekspozycji matrycy ani odświeżenia ekranu, które dodaje najwyżej jedną ramkę wyświetlacza ($8{,}3$\,ms przy $120$\,Hz). Czas reakcji liczony jest od nadejścia pierwszej klatki, na której owocnik został wykryty, do pierwszej publikacji jego ramki, czyli po spełnieniu warunku dwóch kolejnych zgodnych klatek (\texttt{CONFIRM\_STREAK}); osobno zapisano czas do pierwszej zielonej asercji „Jadalny”. Protokół był taki sam jak dla tab.~\ref{tab:latency_s23}: wydruki zdjęć owocników, wymuszona kadencja $4$\,FPS, odrzucona rozgrzewka $8$ klatek, seria zimna (SoC $\approx 28\,^\circ$C) i seria po kilkunastu minutach ciągłej inferencji (SoC $\approx 41\,^\circ$C). Czas reakcji mierzono, odwracając kamerę od zdjęcia na ponad $2$\,s i kierując ją z powrotem; zarejestrowano $138$ takich zdarzeń w całej sesji. Wyniki zestawiono w tab.~\ref{tab:latency_e2e}.

\begin{table}[htbp]
 \centering
 \caption{Pełny cykl klatki i czas reakcji na Samsung Galaxy S23 Ultra. Wartości w milisekundach. Źródło: opracowanie własne na podstawie \texttt{inference\_profile.csv} i \texttt{response\_profile.csv}.}
 \label{tab:latency_e2e}
 \begin{tabular}{|l|c|c|}
 \hline
 \textbf{Wielkość} & \textbf{Zimny SoC} & \textbf{Po obciążeniu} \\
 \hline \hline
 Konwersja klatki CameraX (średnia) & $8{,}1$ & $8{,}1$ \\
 \hline
 Letterbox $640\times640$ (średnia) & $1{,}3$ & $1{,}8$ \\
 \hline
 \texttt{detect()} (średnia) & $152{,}1$ & $278{,}2$ \\
 \hline
 Progi, tracker, stan UI (średnia) & $0{,}6$ & $0{,}6$ \\
 \hline
 \textbf{End-to-end klatki}: średnia / $p_{95}$ & $162{,}7$ / $197{,}3$ & $289{,}1$ / $298{,}5$ \\
 \hline \hline
 \textbf{Czas reakcji} (ramka lub ostrzeżenie): średnia / $p_{95}$ & \multicolumn{2}{c|}{$509{,}4$ / $618{,}0$ ($138$ zdarzeń)} \\
 \hline
 Czas do zielonej asercji „Jadalny” & \multicolumn{2}{c|}{brak zdarzeń ($0$ asercji)} \\
 \hline
 \end{tabular}
\end{table}

W obu seriach pełny cykl klatki pozostaje wyraźnie poniżej progu $1$\,s: średnio $162{,}7$\,ms na zimnym SoC i $289{,}1$\,ms po obciążeniu, przy $p_{95}$ odpowiednio $197{,}3$ i $298{,}5$\,ms. Dominującym składnikiem jest \texttt{detect()} (ok.\ $93\%$ cyklu na zimno i $96\%$ po obciążeniu); pozostałe etapy -- konwersja klatki, letterbox, filtr progów, \texttt{DetectionTracker} i publikacja stanu UI -- sumują się do ok.\ $10$\,ms i nie zależą istotnie od temperatury. Czas reakcji, liczony od pierwszej klatki z wykrytym owocnikiem do publikacji jego ramki, wyniósł średnio $509{,}4$\,ms ($p_{95}=618{,}0$\,ms); jedynie $2$ z $138$ zdarzeń przekroczyły $1$\,s. Wartość ta wynika głównie z wymogu dwóch kolejnych zgodnych klatek (\texttt{CONFIRM\_STREAK}) przy kadencji $250$\,ms, a nie z samego kosztu inferencji. Zielona asercja „Jadalny” nie wystąpiła w sesji terenowej (gatunki w kadrze były trujące lub niejadalne), więc osobnego pomiaru nie odnotowano. Kryterium czasu rzeczywistego rozumiane jako interaktywny czas reakcji poniżej $1$\,s zostało zatem spełnione.

\section{Jakościowa analiza detekcji w warunkach rzeczywistych (dane zewnętrzne)}\label{sec:ood_field}
Metryki z poprzednich podrozdziałów pochodzą z izolowanego zbioru testowego iNaturalist. To ewaluacja wewnątrz rozkładu (ang. \textit{in-distribution}): to samo źródło, podobna optyka i ten sam protokół anotacji co w uczeniu. Aplikacja ma jednak działać w lesie, na strumieniu CameraX, przy oświetleniu podszytu i na matrycy konkretnego smartfonu. Taki reżim jest z definicji poza rozkładem uczącym (ang. \textit{out-of-distribution}, OOD). Dlatego uzupełniono ocenę ilościową jakościową inspekcją w terenie.

Testy przeprowadzono bezpośrednio na urządzeniu Samsung Galaxy S23 Ultra, opisanym w załączniku ze specyfikacją sprzętową. Autor skanował owocniki w naturalnym biotopie leśnym kamerą telefonu (lipiec--sierpień 2026), bez podawania modelowi zdjęć ze zbioru iNaturalist. Poprawność gatunku ustalano wizualnie na podstawie cech makromorfologicznych oraz konsultacji z atlasami mykologicznymi (First Nature, Rogers Plants); identyfikacja in situ nie była weryfikowana laboratoryjnie (np. analizą zarodników pod mikroskopem lub sekwencjonowaniem DNA), co stanowi ograniczenie tej oceny terenowej. Zebrano kilkadziesiąt zrzutów ekranu nakładki CameraX + Jetpack Compose; w załączniku~\ref{app:ood_screens} pozostawiono tylko interfejs, przykład detekcji wieloobiektowej, aktywację progu $0{,}18$ oraz dwa błędy modelu. Zrzuty w załączniku pochodzą z 26.07--01.08.2026, czyli ze wcześniejszej kompilacji niż progi $0{,}18$/$0{,}60$ wpisane do kodu 11.08.2026: \textit{Amanita pantherina} miała wtedy status „Trujący”, a nie obecny „Śmiertelnie trujący”. Nie rozstrzygają one progu $0{,}18$ wobec $0{,}30$. Ponieważ kadry terenowe nie mają ramek referencyjnych, nie liczono $mAP$. Oceniano zgodność gatunku, poprawność kategorii toksykologicznej oraz czytelność interfejsu.

Zaobserwowane scenariusze, uporządkowane od najprostszych do najtrudniejszych, zestawiono w tab.~\ref{tab:ood_scenariusze}.

\begin{table}[htbp]
 \centering
 \small
 \caption{Scenariusze jakościowej ewaluacji aplikacji na danych zewnętrznych (kamera smartfonu, biotop leśny). Źródło: opracowanie własne.}
 \label{tab:ood_scenariusze}
 \begin{tabular}{|p{2.55cm}|p{4.15cm}|p{3.55cm}|p{3.15cm}|}
 \hline
 \textbf{Scenariusz} & \textbf{Przykładowe gatunki} & \textbf{Status UI} & \textbf{Obserwacja} \\
 \hline \hline
 Pojedynczy owocnik jadalny & \textit{Pluteus cervinus}, \textit{Gyroporus castaneus}, \textit{Xerocomus chrysenteron} & zielona ramka, status Jadalny & poprawna asercja przy pewnym ujęciu \\
 \hline
 Wiele owocników śmiertelnych & \textit{Amanita phalloides} ($3+$ w kadrze) & czerwona ramka, Śmiertelnie trujący & jednoczesne wykrycie wszystkich widocznych owocników \\
 \hline
 Sceny mieszane i kępy & \textit{Amanita pantherina} / \textit{Amanita phalloides}, \textit{Hypholoma fasciculare} & pomarańczowa ramka, status Trujący & wieloklasowa detekcja w jednym kadrze \\
 \hline
 Gatunki nadrzewne & \textit{Tremella mesenterica}, \textit{Schizophyllum commune} & szara ramka, status Niejadalny & status neutralny, bez asercji jadalności \\
 \hline
 Tryb oszczędzania baterii & dowolny ustabilizowany kadr & baner trybu oszczędzania baterii & przejście na \textit{heartbeat} $1{,}5$\,s \\
 \hline
 \end{tabular}
\end{table}

W ujęciach pojedynczych gatunków jadalnych (\textit{Pluteus cervinus}, \textit{Gyroporus castaneus}, \textit{Xerocomus chrysenteron}, a także \textit{Russula cyanoxantha}) model lokalizował kapelusz zieloną ramką i przypisywał status „Jadalny” przy pewnym, zbliżonym kadrze. Nie obserwowano w tych warunkach fałszywej eskalacji do ostrzeżenia toksycznego. Przykład takiego ujęcia, zestawiony z kadrem o najwyższym priorytecie bezpieczeństwa, przedstawia rys.~\ref{fig:ood_jadalny_toksyczny}.

W scenach z wieloma owocnikami muchomora sromotnikowego (\textit{Amanita phalloides}) detektor oznaczał jednocześnie trzy lub więcej owocników w jednym kadrze. Każdy otrzymywał natychmiast czerwoną ramkę i etykietę „Śmiertelnie trujący”, zgodnie z progiem $\mathrm{conf}_{\mathrm{toxic}}=0{,}18$. W kadrze mieszanym obok muchomora sromotnikowego pojawiał się muchomor plamisty (\textit{Amanita pantherina}) -- system rozróżniał obie klasy i oznaczał ramką pomarańczową „Trujący” dla plamistego, nie degradując statusu sromotnikowego. Gęste kępy maślanki wiązkowej (\textit{Hypholoma fasciculare}) oraz kadry wielogatunkowe z \textit{Mycena galericulata} potwierdziły, że wieloobiektowa głowica YOLO11m nie sprowadza całej sceny do jednej etykiety: sąsiadujące kępy otrzymują osobne ramki.

Gatunki nadrzewne -- trzęsak pomarańczowożółty (\textit{Tremella mesenterica}) na korze oraz rozszczepka pospolita (\textit{Schizophyllum commune}) na leżaninie -- otrzymywały szarą ramkę „Niejadalny”. Warstwa prezentacji nie eskalowała ich ani do zielonej asercji jadalności, ani do fałszywego alarmu toksycznego. Po ustabilizowaniu kadru nakładka wyświetlała baner „Zablokowano -- tryb oszczędzania baterii”: \texttt{isSceneLocked()} podnosiła interwał inferencji z $250$\,ms do \textit{heartbeat} $1{,}5$\,s, zgodnie z implementacją z rozdz.~\ref{chap:implementacja}. Oba zjawiska -- detekcję nadrzewną i blokadę energetyczną -- widać na rys.~\ref{fig:ood_battery_nadrzewne}.

\begin{figure}[H]
 \centering
 \begin{subfigure}[t]{0.48\textwidth}
  \centering
  \includegraphics[width=\linewidth,height=11.2cm,keepaspectratio]{screenshots/ood_01_pluteus_cervinus.jpg}
  \caption{\textit{Pluteus cervinus}, jadalny, zielona ramka.}
  \label{fig:ood_pluteus}
 \end{subfigure}
 \hfill
 \begin{subfigure}[t]{0.48\textwidth}
  \centering
  \includegraphics[width=\linewidth,height=11.2cm,keepaspectratio]{screenshots/ood_05_amanita_mixed.jpg}
  \caption{\textit{Amanita phalloides} (czerwone) i \textit{Amanita pantherina} (pomarańczowa) w jednym kadrze.}
  \label{fig:ood_amanita_mix}
 \end{subfigure}
 \caption{Detekcja w warunkach terenowych: pojedynczy owocnik jadalny oraz równoczesne oznaczenie gatunków śmiertelnych i trujących. Źródło: opracowanie własne (zrzuty aplikacji, CameraX + Jetpack Compose).}
 \label{fig:ood_jadalny_toksyczny}
\end{figure}

\begin{figure}[H]
 \centering
 \begin{subfigure}[t]{0.48\textwidth}
  \centering
  \includegraphics[width=\linewidth,height=11.2cm,keepaspectratio]{screenshots/ood_09_tremella_battery.jpg}
  \caption{\textit{Tremella mesenterica}, niejadalny; baner trybu oszczędzania baterii.}
  \label{fig:ood_tremella}
 \end{subfigure}
 \hfill
 \begin{subfigure}[t]{0.48\textwidth}
  \centering
  \includegraphics[width=\linewidth,height=11.2cm,keepaspectratio]{screenshots/ood_07_hypholoma_fasciculare.jpg}
  \caption{Kępy \textit{Hypholoma fasciculare}, trujący, pomarańczowe ramki.}
  \label{fig:ood_hypholoma}
 \end{subfigure}
 \caption{Praca aplikacji poza zbiorem iNaturalist: gatunek nadrzewny z blokadą energetyczną oraz wieloobiektowa detekcja kęp trujących. Źródło: opracowanie własne.}
 \label{fig:ood_battery_nadrzewne}
\end{figure}

Ograniczenia tej ewaluacji są jasne i nie należy ich ukrywać. Próba jest okazjonalna, sezonowa (lipiec--sierpień) i nie pokrywa wszystkich $147$ klas. Częściowe przesłonięcie ściółką, źdźbłami trawy lub gałązkami obniża pewność i bywa przyczyną pominięcia drugiego, mniejszego owocnika w kadrze; gęste kępy generują nachodzące etykiety nakładki. Zmienne oświetlenie podszytu nie uniemożliwiało jednak poprawnej klasyfikacji kategorii toksykologicznej w zebranych ujęciach poprawnych. Fałszywy pozytyw na kamieniu oraz pominięcie mniejszego skupiska \textit{Tremella} udokumentowano w załączniku~\ref{app:ood_screens}. W zebranym materiale terenowym nie zaobserwowano błędu typu „śmiertelnie trujący prezentowany jako jadalny”: czerwone i pomarańczowe ramki pojawiały się już od progu ostrzeżenia $0{,}18$, a zielony status wymagał pewności $\ge 0{,}60$. To zgodność z logiką \texttt{displayStatus} na niewielkiej próbie OOD, a nie pomiar CFPR. $H_2$ weryfikuje tab.~\ref{tab:cfpr_porownanie} na izolowanym teście iNaturalist. Bufor konsensusu $4$ z $5$ klatek, opisany w rozdz.~\ref{chap:implementacja}, działał w tle jako mechanizm stabilizacji nakładki; nie był przedmiotem osobnego porównania z wariantem bez bufora.

\section{Ilościowa ewaluacja terenowa z anotacją prawdy podstawowej}\label{sec:field_quant}
Jakościowa inspekcja z podrozdz.~\ref{sec:ood_field} nie pozwala policzyć skuteczności, ponieważ nie rejestrowała predykcji w formie umożliwiającej późniejszą weryfikację. Aby uzyskać liczby przy zachowaniu naturalnego reżimu OOD, do aplikacji dodano ukryty \textit{Tryb Badawczy} (ang. \textit{Research Mode}), aktywowany gestem w interfejsie. Każde naciśnięcie przycisku zapisywało pełny pakiet danych: pełnowymiarowe zdjęcie z toru \texttt{ImageCapture}, zrzut nakładki, metadane predykcji (pięć najlepszych klas z pewnościami) oraz telemetrię (bateria, temperatura, FPS, opóźnienie \texttt{detect()}). Zebrano $67$ kadrów w dwóch sesjach terenowych ($15$ września $2026$, Samsung Galaxy S23 Ultra, biotop leśny).

Prawdę podstawową (ang. \textit{ground truth}) ustalono hybrydowo: dla każdego kadru autor weryfikował gatunek na pełnowymiarowym zdjęciu, korzystając z cech makromorfologicznych i atlasów mykologicznych. Podobnie jak w podrozdz.~\ref{sec:ood_field}, identyfikacja nie była potwierdzana laboratoryjnie (mikroskopia zarodników, sekwencjonowanie DNA), co pozostaje ograniczeniem oceny. Agregację i obliczenie metryk wykonał skrypt \texttt{analyze\_field\_research.py} na podstawie zapisanych metadanych i ręcznie uzupełnionej tabeli anotacji.

Wszystkie miary w tym podrozdziale liczone są na poziomie kadru. Kluczową obserwacją metodyczną jest rozdzielenie dwóch źródeł błędu, których nie wolno mieszać w jednej liczbie: pominięcia na etapie \textbf{detekcji} (owocnik jest w kadrze, lecz model nie zwraca żadnej ramki) oraz pomyłki na etapie \textbf{klasyfikacji} (ramka jest, lecz przypisany gatunek jest błędny). Zestawienie trzech poziomów skuteczności przedstawia tab.~\ref{tab:field_accuracy}.

Jednostką analizy jest pojedynczy kadr (jedno naciśnięcie przycisku przechwytywania), a~nie pojedynczy owocnik. Każdemu kadrowi przypisano jeden docelowy gatunek -- owocnik w~centrum zainteresowania zbieracza; detekcję uznawano za udaną, gdy model zwrócił co najmniej jedną ramkę, a~trafienie top-1 -- gdy ramka o~najwyższej pewności wskazywała właściwy gatunek. Podawane wartości są więc odsetkami \emph{na kadr}, a~nie ścisłym recall na poziomie pojedynczego obiektu; w~kadrach z~wieloma owocnikami (oznaczonych w~danych jako $\times n$) oceniano owocnik dominujący. Ponadto co najmniej trzy pary kadrów to serie tego samego egzemplarza wykonane w~odstępie sekund ($24$/$25$, $45$/$46$, $66$/$67$; prawdopodobnie też $42$/$43$), zatem $67$ kadrów odpowiada ok.\ $63$ odrębnym owocnikom. Liczby należy traktować jako charakterystykę jakościową próby okazjonalnej, a~nie precyzyjny estymator populacyjny. Średnia pewność predykcji w~próbie wyniosła $0{,}736$, a~średni czas \texttt{detect()} -- $154{,}7$\,ms (zakres $136$--$196$\,ms), spójnie z~profilowaniem opóźnień na zimnym SoC (tab.~\ref{tab:latency_s23}).

W~tych samych $67$ pakietach Trybu Badawczego zapisano telemetrię telefonu, a~nie stacji roboczej. We wszystkich kadrach \texttt{gpuDelegate} był włączony, więc ta sesja leśna szła przez delegata GPU, a~nie przez rezerwową ścieżkę CPU. Temperatura baterii wynosiła $27{,}7$--$30{,}8^{\circ}\mathrm{C}$: to reżim zbliżony do serii zimnej, a~nie do wymuszonego nagrzewania przez $8$--$9$\,min. Żaden z~$67$ pomiarów \texttt{detect()} nie przekroczył $250$\,ms. Równolegle zapisany wskaźnik podglądu kamery wynosił średnio $33$\,FPS (od $30$ do $50$). Nie jest to przepustowość modelu: przy \texttt{detect()} rzędu $155$\,ms inferencja jest wolniejsza niż podgląd, a~nadmiarowe klatki odrzuca \texttt{STRATEGY\_KEEP\_ONLY\_LATEST}. W pakietach Trybu Badawczego nie rejestrowano czasu letterboxingu; pełny cykl klatki zmierzono w podrozdz.~\ref{sec:latency_e2e}. W tej sesji wąskim gardłem był \texttt{detect()}, nie dostarczanie klatek.

\begin{table}[htbp]
 \centering
 \small
 \caption{Trójpoziomowa skuteczność na kadrach terenowych, liczona na poziomie kadru, a nie pojedynczego owocnika: wszystkie $67$ kadrów oraz podzbiór $57$ kadrów z~gatunkami należącymi do rejestru $147$ klas (closed-set). Źródło: opracowanie własne.}
 \label{tab:field_accuracy}
 \begin{tabular}{|p{6.0cm}|c|c|}
 \hline
 \textbf{Miara} & \textbf{Wszystkie ($67$)} & \textbf{Closed-set ($57$)} \\
 \hline \hline
 Kadry bez żadnej ramki & $17{,}9\%$ ($12/67$) & $12{,}3\%$ ($7/57$) \\
 \hline
 Klasyfikacja top-1 (gdy jest detekcja) & $81{,}8\%$ ($45/55$) & $90{,}0\%$ ($45/50$) \\
 \hline
 Klasyfikacja top-5 (gdy jest detekcja) & $89{,}1\%$ ($49/55$) & $98{,}0\%$ ($49/50$) \\
 \hline
 Skuteczność łączna (detekcja i klasyfikacja) top-1 & $67{,}2\%$ ($45/67$) & $78{,}9\%$ ($45/57$) \\
 \hline
 Skuteczność łączna (detekcja i klasyfikacja) top-5 & $73{,}1\%$ ($49/67$) & $86{,}0\%$ ($49/57$) \\
 \hline
 \end{tabular}
\end{table}

Wynik należy czytać dwutorowo. Po pierwsze, $10$ z~$67$ kadrów przedstawia gatunki spoza rejestru $147$ klas (OOD), których model z~definicji nie może wskazać poprawnie; dlatego w~tab.~\ref{tab:field_accuracy} podano osobno wariant \emph{closed-set} ($57$ kadrów z~gatunkami należącymi do rejestru). Po drugie, w~wariancie closed-set oba etapy potoku ograniczają skuteczność w~zbliżonym stopniu: detektor zwraca co najmniej jedną ramkę w~$87{,}7\%$ kadrów ($50/57$), a~klasyfikator trafia w~$90{,}0\%$ zlokalizowanych przypadków ($45/50$). Nie można więc wskazać jednego dominującego wąskiego gardła -- obie składowe wnoszą porównywalny wkład do błędu łącznego, przy czym błędy detekcji koncentrują się na ujęciach nietypowych i~przesłoniętych, a~błędy klasyfikacji -- na gatunkach bliźniaczych (głównie borowikokształtnych). Wszystkie $12$ kadrów „bez detekcji” zawierało rzeczywisty owocnik (potwierdzone wizualnie), są to więc kadry fałszywie negatywne (na poziomie kadru), a~nie puste sceny. Przyczyny pominięć zestawia tab.~\ref{tab:field_fn}.

\begin{table}[htbp]
 \centering
 \small
 \caption{Przyczyny pominięć detekcji na $12$ fałszywie negatywnych kadrach terenowych. Źródło: opracowanie własne.}
 \label{tab:field_fn}
 \begin{tabular}{|p{7.5cm}|c|}
 \hline
 \textbf{Przyczyna pominięcia} & \textbf{Liczba kadrów} \\
 \hline \hline
 Częściowe przesłonięcie (trawa, liście, gałązki, światło kontrowe) & $5$ \\
 \hline
 Nietypowe ujęcie lub morfologia (od spodu, owocnik zdeformowany) & $3$ \\
 \hline
 Owocnik mały lub odległy w kadrze & $2$ \\
 \hline
 Owocnik trzymany w dłoni / ścięty & $2$ \\
 \hline
 \end{tabular}
\end{table}

Charakter tych pominięć jest spójny: model uczono na „katalogowych” zdjęciach iNaturalist, wykonywanych zwykle z góry na dobrze wyeksponowanym owocniku. W terenie zawodzi on przede wszystkim przy częściowym przesłonięciu ściółką lub roślinnością oraz na ujęciach nietypowych (od spodu, w dłoni, owocniki zdeformowane) -- jest to typowa luka dziedzinowa (ang. \textit{domain gap}), a nie wada architektury detektora. Przydział poszczególnych kadrów do kategorii ma charakter jakościowy.

Na $10$ z $55$ zlokalizowanych kadrów klasyfikacja była błędna. Cztery pomyłki mieściły się w obrębie grzybów rurkowych (borowik$\leftrightarrow$borowik, wewnątrz rodzaju \textit{Suillus}) -- morfologicznie zbliżonych i, co istotne, żadna z nich nie przekroczyła granicy między gatunkiem niebezpiecznym a jadalnym, pozostając neutralną z perspektywy bezpieczeństwa. W czterech przypadkach pomyłka była grubsza (mylenie odległych rodzajów), a w dwóch pozostałych (kadry z \textit{Phaeotremella foliacea}) model błędnie zgłosił śmiertelnie trującą \textit{Gyromitra esculenta} -- był to jednak fałszywy alarm (kierunek bezpieczny), a nie zaniżenie zagrożenia. W żadnym z $10$ błędów gatunek rzeczywiście niebezpieczny nie został przedstawiony jako jadalny.

W wymiarze bezpieczeństwa spośród $57$ kadrów o znanej kategorii toksykologicznej $7$ kadrów (odpowiadających $4$ gatunkom: \textit{Lactarius helvus}, \textit{Hypholoma fasciculare}, \textit{Russula foetens}, \textit{Amanita muscaria}; wszystkie klasy POISONOUS, bez gatunku śmiertelnie trującego w~tej próbie) należało do kategorii niebezpiecznej i wszystkie $7$ oznaczono jako niebezpieczne; żaden nie otrzymał asercji „Jadalny”. Przy $n=7$ i bez gatunku śmiertelnie trującego nie jest to estymator czułości na zagrożenie. Wynik ten należy jednak interpretować ostrożnie z dwóch powodów. Po pierwsze, próba jest mała. Po drugie, miara fałszywych alarmów jest zaniżona: skrypt pomija kadry, których gatunek prawdziwy nie występuje w tablicy toksyczności -- dlatego wspomniane fałszywe ostrzeżenie „\textit{Gyromitra esculenta}” na nieszkodliwym trzęsaku zostało wykluczone z bilansu, gdyż \textit{Phaeotremella foliacea} nie należy do $147$ klas modelu. W ogóle $10$ z $67$ kadrów przedstawiało gatunki spoza katalogu modelu (m.in. \textit{Phaeotremella foliacea}, \textit{Chalciporus piperatus}, \textit{Phaeolus schweinitzii}), co unaocznia zasadnicze ograniczenie: w rzeczywistym lesie występuje znacznie więcej gatunków, niż obejmuje $147$ klas modelu.

\chapter{Podsumowanie i wnioski}\label{chap:podsumowanie}
Zaprojektowano i wdrożono autonomiczny system na platformę Android. Aplikacja detekuje i klasyfikuje offline $147$ klas grzybów, odpowiadających $145$ odrębnym gatunkom biologicznym. Klatki analizowane są z kadencją $250$\,ms ($4$\,FPS); na gorącym SoC samo \texttt{detect()} często ją przekracza (tab.~\ref{tab:latency_s23}), co obniża chwilową kadencję. Kryterium czasu rzeczywistego -- czas reakcji poniżej $1$\,s -- zweryfikowano pomiarem end-to-end (tab.~\ref{tab:latency_e2e}): spełnione, przy średnim czasie reakcji $509{,}4$\,ms ($p_{95}=618{,}0$\,ms) i pełnym cyklu klatki $162{,}7$--$289{,}1$\,ms zależnie od temperatury SoC. Potok obejmował akwizycję danych, auto-anotację GroundingDINO, trening YOLO11m, konwersję TFLite FP16 oraz aplikację z asymetryczną warstwą progów, która priorytetyzuje czułość dla klas toksycznych.

\section{Wnioski z realizacji projektu}
Cel główny – autonomiczny, offline system identyfikacji grzybów na Androidzie – został osiągnięty. Rozwiązanie łączy detekcję wieloobiektową z nakładką kamery i trwałym zapisem historii znalezisk, spełniając założenia z rozdz.~\ref{chap:projekt_systemu}.

Potok półautomatycznej anotacji oparty na GroundingDINO zastąpił masowe ręczne znakowanie ramek auto-anotacją zero-shot, przy zachowaniu formatu YOLO ($\texttt{class\_id cx cy w h}$) dla ponad stu gatunków \cite{liu2023grounding}. Ręczny audyt $1\,065$ auto-anotowanych obrazów potwierdził poprawność lokalizacji dla $94{,}93\%$ istniejących ramek ($1\,011$ z $1\,065$). Na stratyfikowanej próbie $250$ kadrów izolowanego testu nauczyciel pominął $34$ z $441$ widocznych owocników ($7{,}71\%$; tab.~\ref{tab:audyt_fn}). Bez tego mechanizmu zbudowanie zbioru uczącego o takiej skali byłoby poza zasięgiem pracy inżynierskiej.

Model YOLO11m osiągnął na izolowanym zbiorze testowym wartości $mAP_{50}=82{,}75\%$ oraz $mAP_{50\text{--}95}=78{,}23\%$ (PyTorch, \texttt{split=test}, zero wspólnych \texttt{obs\_id} z korpusem uczącym). Na tym samym teście wariant Nano uzyskał $48{,}54\%$ / $42{,}82\%$, co potwierdza wybór Medium. Na próbie $2\,500$ kadrów (\texttt{seed}=42) artefakt TFLite FP16 z dekoderem Androida dał $mAP_{50}=0{,}7914$ wobec $0{,}8467$ dla \texttt{.pt} (różnica $-0{,}0553$; podrozdz.~\ref{sec:tflite_parity}). Na stacji roboczej (RTX~4070) potok algorytmiczny zamyka się w $10{,}94$\,ms; na S23 Ultra średni czas \texttt{detect()} wynosi $156{,}5$\,ms na zimnym SoC i $245{,}9$\,ms po kilkuminutowym obciążeniu. Pełny cykl klatki od analizatora CameraX do publikacji nakładki trwał średnio $162{,}7$\,ms na zimnym SoC i $289{,}1$\,ms po obciążeniu, a czas do pierwszej ramki lub ostrzeżenia $509{,}4$\,ms ($p_{95}$ $618{,}0$\,ms), czyli poniżej przyjętego progu interaktywności $1$\,s (tab.~\ref{tab:latency_e2e}). Wyniki te potwierdzają praktyczną wykonalność podejścia \textit{edge AI} w zadaniu identyfikacji mykologicznej, bez zależności od zasobów chmurowych \cite{ultralytics2024yolo11}.

Polityka \texttt{displayStatus} nie prezentuje zielonej etykiety „Jadalny” przy pewności poniżej $0{,}60$ ani dla predykcji niejadalnych. Na izolowanym teście CFPR pary $0{,}18$/$0{,}60$ wynosi dla wag PyTorch $2{,}06\%$ ($327/15\,840$), wobec $3{,}66\%$ ($580/15\,840$) przy $\tau=0{,}30$, a dla artefaktu TFLite FP16 z dekoderem aplikacji $2{,}11\%$ wobec $3{,}64\%$ -- $H_2$ spełniona na obu pomiarach. Progi ustalono przed zbudowaniem izolowanego testu (podrozdz.~\ref{sec:safety_matrix}). Jakościowa inspekcja w terenie (podrozdz.~\ref{sec:ood_field}) pokazała, że zielone ramki jadalne, czerwone i pomarańczowe ostrzeżenia oraz szary status niejadalny działają na strumieniu CameraX poza zbiorem iNaturalist. Nie zastępuje to oceny mykologa: około $2\%$ obiektów toksycznych nadal może otrzymać asercję jadalności.

Ilościowa ewaluacja terenowa na $67$ kadrach (podrozdz.~\ref{sec:field_quant}) potwierdziła te obserwacje liczbowo. Klasyfikator osiągnął $81{,}8\%$ top-1 dla zlokalizowanych owocników ($90{,}0\%$ w wariancie closed-set), przy czym wszystkie $7$ przypadków gatunków niebezpiecznych w~próbie oznaczono ostrzeżeniem, a żaden nie otrzymał statusu „Jadalny”. Zarazem ujawniła ona, że detekcja i~klasyfikacja ograniczają skuteczność terenową w~porównywalnym stopniu (w~wariancie closed-set detektor nie zwraca żadnej ramki w~ok.\ $12\%$ kadrów, głównie na ujęciach nietypowych (miara na poziomie kadru), a~klasyfikator myli ok.\ $10\%$ zlokalizowanych owocników na gatunkach bliźniaczych), co wyznacza dwa równoległe kierunki dalszych prac.

\section{Perspektywy rozwoju i dalsze prace}
Kolejnym krokiem może być integracja algorytmów śledzenia obiektów wideo (np. ByteTrack lub SORT) z potokiem CameraX. Stabilizacja ramek w czasie zmniejszyłaby chwilowe wahania predykcji pomiędzy klatkami i poprawiła płynność nakładki interfejsu użytkownika.

Drugim istotnym obszarem pozostaje optymalizacja rozmiaru modelu. Warto zbadać uczenie z uwzględnieniem kwantyzacji (ang. \textit{Quantization-Aware Training}, QAT) oraz skalibrowaną precyzję mieszaną (FP16/INT8), tak aby zredukować artefakt poniżej $20$\,MiB przy zachowaniu wrażliwości na subtelne cechy makromorfologiczne (blaszki, pierścień), kluczowe dla rozróżniania gatunków bliźniaczych.

Trzecią perspektywą jest fuzja osadzeń wizyjnych z metadanymi kontekstowymi: klasyfikacją biotopu na podstawie GPS, typem lasu, sezonowością oraz lokalnymi warunkami pogodowymi. Takie wzbogacenie sygnału wejściowego mogłoby poprawić rozróżnianie gatunków o zbliżonej morfologii, lecz odmiennych wymaganiach siedliskowych.

Ewaluacja terenowa (podrozdz.~\ref{sec:field_quant}) wskazała trzy kolejne, praktyczne kierunki rozwoju. Pierwszym jest \textbf{rozszerzenie katalogu gatunków}. Aż $10$ z $67$ owocników napotkanych w terenie należało do gatunków spoza $147$ klas modelu, co powodowało wymuszone błędne przypisania do klasy najbliższej wizualnie (w tym pojedyncze fałszywe ostrzeżenia toksyczne). Zwiększenie liczby klas -- ze szczególnym uwzględnieniem gatunków pospolitych w docelowym biotopie -- bezpośrednio ograniczyłoby tę kategorię błędu. Przy okazji rewizji rejestru należy scalić klasy będące synonimami taksonomicznymi (\textit{Macrolepiota procera}/\textit{Lepiota procera} oraz \textit{Collybia dryophila}/\textit{Gymnopus dryophilus}), odziedziczone po etykietach iNaturalist (zob.\ tab.~\ref{tab:zal2_gatunki}).

Drugim kierunkiem jest \textbf{poprawa jakości i różnorodności zbioru uczącego}. Analiza pominięć detekcji wykazała, że model radzi sobie z „katalogowymi” ujęciami z góry, lecz zawodzi na zdjęciach nietypowych: od spodu (widoczne pory lub blaszki), w dłoni, przy silnym przesłonięciu ściółką oraz na owocnikach starych i zniekształconych. Uzupełnienie korpusu uczącego o takie właśnie ujęcia -- pozyskane w terenie lub wygenerowane przez ukierunkowaną augmentację (przesłonięcia, obroty, symulacja zmiennego oświetlenia podszytu) -- powinno zredukować lukę dziedzinową między danymi iNaturalist a strumieniem kamery.

Trzecim kierunkiem jest \textbf{ponowne przemyślenie etapu lokalizacji owocnika}. Sam detektor YOLO11m działa dobrze w~warunkach zbliżonych do danych uczących (na izolowanym teście $mAP_{50}=82{,}75\%$), lecz w~terenie odsetek kadrów, na których zwraca co najmniej jedną ramkę, spada do ok.\ $82\%$ na całej próbie i~$87{,}7\%$ w~wariancie closed-set (miara na poziomie kadru, a~nie klasyczny recall obiektów). Jest to -- obok rozróżniania gatunków bliźniaczych -- jedno z~dwóch porównywalnych ograniczeń skuteczności terenowej, a~nie samodzielne wąskie gardło. Warto zbadać techniki podnoszące czułość lokalizacji niezależnie od klasyfikacji: augmentację w czasie testu (ang. \textit{test-time augmentation}), inferencję kafelkową i wieloskalową dla drobnych lub częściowo przesłoniętych owocników, obniżony próg pewności na samym etapie lokalizacji (z zachowaniem wyższego progu dla asercji statusu) oraz pętlę aktywnego uczenia (ang. \textit{active learning}), w której zebrane w terenie fałszywe negatywy trafiają z powrotem do zbioru uczącego. Mechanizm zapisu pełnego pakietu danych z Trybu Badawczego stanowi gotową infrastrukturę dla takiej pętli.

\section{Repozytorium projektu}
Kod źródłowy aplikacji mobilnej w języku Kotlin, skrypty potoku przetwarzania danych w języku Python oraz narzędzia eksportu modeli zostały udostępnione w publicznym repozytorium pod adresem: \url{https://github.com/horobetk/mushroom-detection}. Wersję aplikacji opisaną w pracy (plik APK) udostępniono w zakładce \textit{Releases} tego repozytorium; nazwy skryptów przywoływanych w tekście odpowiadają plikom w katalogu \texttt{scripts/}.

W przygotowaniu niniejszej pracy wykorzystano narzędzia oparte na metodach sztucznej inteligencji jako wsparcie redakcji oraz korekty językowej i stylistycznej tekstu; zaproponowane zmiany były każdorazowo weryfikowane i akceptowane przez autora.
Narzędzia oparte na metodach sztucznej inteligencji wykorzystano również pomocniczo przy opracowaniu i korekcie kodu źródłowego; wszystkie sugestie były każdorazowo weryfikowane i akceptowane przez autora.
